% concatenated from Critical-7-27-2026.tex
\documentclass[a4paper,11pt]{article}
\usepackage[utf8]{inputenc}
\usepackage{amsfonts, amssymb, amsmath, amsthm,tikz}
\usepackage{color}
%\usepackage{pgfplots}

%\tikzexternalize
%\usepackage[colorlinks=true,backref=page]{hyperref}fdsymbol

\usepackage{srcltx}

\usepackage{epsfig,graphicx}
\usepackage[margin=1.25in]{geometry}

\newtheorem{theorem}{Theorem}
\newtheorem*{theorem*}{Theorem}

\newtheorem{theoremName}{Theorem}

\renewcommand{\thetheorem}{\Alph{theorem}}

\newtheorem{lemma}{Lemma}
\newtheorem{claim}[lemma]{Claim}
\newtheorem{question}[lemma]{Question}
\newtheorem{thm}{Theorem}
\newtheorem*{thmB}{Theorem B'}

%\newtheorem*{namedthm}{Theorem}

\newtheorem{conjecture}[lemma]{Conjecture}
\newtheorem{corollary}[lemma]{Corollary}
\newtheorem{proposition}[lemma]{Proposition}
\newtheorem{remark}[lemma]{Remark}
\newtheorem{remarks}[lemma]{Remarks}
\newtheorem*{main-conjecture}{Main conjecture}
\newtheorem{addendum}[lemma]{Addendum}
\newtheorem{quest}{\bf Question}


\theoremstyle{definition}
\newtheorem{definition}[lemma]{Definition}

\def\DD{{\mathbb D}}  
\def\TT{{\mathbb T}}  

\def\ZZ{{\mathbb Z}}
\def\AA{{\mathbb A}}
\def\RR{{\mathbb R}}
\def\CC{{\mathbb C}}
\def\NN{{\mathbb N}}

\def\cC{\mathcal{C}}
\def\cO{\mathcal{O}}
\def\cI{\mathcal{I}}
\def\cJ{\mathcal{J}}
\def\cD{\mathcal{D}}
\def\La{\Lambda}
\def\la{\lambda}
\def\e{\epsilon}
\def\eps{\epsilon}

\def\ga{\gamma}
\def\Ga{\Gamma}
\def\de{\delta}
\def\De{\Delta}
\def\La{\Lambda}
\def\be{\beta}
\def\eps{\epsilon}
\def\al{\alpha}
\def\om{\omega}
\def\Ga{\Gamma}
\def\si{\sigma}

\def\supp{{\rm {supp}}}
\def\emp{\emptyset}
\def\diff{\operatorname{Diff}}
\def\diffdi{\diff_{\operatorname{Diss}}}
\def\diffdiss{\diff_{\operatorname{\chi-SDiss}}}
\def\diffdisschi{\diff_{\operatorname{\chi-SDiss}}}
\def\diffpdiss{\diff_{\operatorname{PerSDiss}}}
\def\interior{\operatorname{Interior}}

\def\noi{\noindent}


\date{\today}
  \title{Critical set for surface diffeomorphisms revisited}
\author{Sylvain Crovisier, Enrique Pujals}



\begin{document}
\maketitle
\begin{abstract}
We propose a notion of \emph{critical set} for surface diffeomorphisms, defined intrinsically in terms of the projective action of the derivative, that plays a role analogous to that of critical points in one-dimensional dynamics. Under a non-degeneracy assumption on the derivative cocycle, a compact invariant set admits a dominated splitting if and only if it contains no critical points. 
For dissipative diffeomorphisms of the disk we derive structural
consequences from the critical set and its recurrence: a dichotomy for
the closure of unstable manifolds of saddle periodic points, and, when the critical set is
non-recurrent, a finite decomposition of the dynamics analogous to the
one-dimensional Misiurewicz case.




\end{abstract}
\section{Introduction}


% The m
% In this paper, our objective is to identify and define a {\em critical set } for two-dimensional  surface diffeomorphisms as an intrinsically defined object designed to play a role analogous to that of {\em critical points} in one-dimensional dynamics. 


The motivation of the present paper stems from the classical understanding that, in the one-dimensional setting, the presence of critical points (where the derivative vanishes)  is a key obstruction to hyperbolicity. Conversely, smooth  one-dimensional endomorphisms without critical points are either hyperbolic, conjugate to hyperbolic or conjugate to an irrational rotation; see \cite{M}, where a real one dimensional version of Fatou's hyperbolicity criteria for holomorphic endomorphisms of the Riemann sphere was proved. In this sense, any compact invariant set is ``essentially hyperbolic" if and only if it does not contain critical points.
%within the limit set
Extending this perspective to two dimensions, we pursue a similar criterion: to isolate the phenomena that {\em obstruct hyperbolicity for surface diffeomorphisms} and, dually, {\em those whose absence guarantees it}. In this context, the critical set plays a role analogous to that of critical points in one-dimensional dynamics, and designed to capture precisely those dynamical features that impede hyperbolic behavior.


Since it was proved in \cite{PS1} that any compact invariant set exhibiting a {\it dominated splitting} for a generic $C^2$ surface diffeomorphism is a hyperbolic set, it follows that searching for dynamical obstructions to hyperbolicity is essentially equivalent to identifying obstructions to domination. The critical set we propose is therefore naturally aligned with this goal: {\it a well defined set whose presence obstructs domination, and whose absence corresponds to  the existence of a dominated splitting}.

A similar task was carried out in \cite{PRH}; we follow that perspective while simplifying the notion of the critical set introduced there, removing certain cumbersome features that limited its applicability and obtaining a more accessible and intrinsic description of criticality in surface dynamics. 
This concept was further extended in \cite{V} to the context of automorphisms of $\mathbb{C}^2$, underscoring its broader relevance in the study of dynamical systems beyond the real two-dimensional setting.



Informally, to exhibit a dominated splitting means that there is an invariant cone field that is forwardly contracted and backwardly expanded; in that sense, a critical point is a point with {\it ``a cone that is not contracted neither for the future nor for the past"}. More precisely, this notion is recast in terms of the projective dynamics induced by the derivative (see paragraph a- in the present section for a formal definition). Moreover, under the assumption of ``far from homotheties'' (which is implied by sufficient dissipation, see paragraph b), the absence of critical points is equivalent to the existence of a dominated splitting, and no full orbit consists entirely of critical points.

 

 
Tangencies associated with the stable and unstable manifold of periodic points (see paragraph i- in section \ref{examples} for the precise definition and discussion) are a natural obstruction to dominated splitting. On the one hand, however, tangencies are invariant under iteration and in that sense would not qualify as a good definition for a critical point. 


However, there is a robust reason why the critical set cannot be reduced to tangencies
of periodic points. Kupka--Smale diffeomorphisms (those for which all
periodic points are hyperbolic and all intersections between their stable and
unstable manifolds are transversal) form a residual subset of
$Diff^r(M)$; hence, generically, there are no tangencies associated with
periodic points at all. On the other hand, by the results of Newhouse
\cite{N} there exist \emph{open} sets of surface diffeomorphisms whose
non-wandering set is not hyperbolic and does not admit a dominated splitting.
By theorem \ref{p.existence}, every diffeomorphism in such an open set
carries critical points. Critical points are therefore present on open sets
of diffeomorphisms on which no periodic tangency exists: the notion is not a
reformulation of tangencies, but a strictly more general, intrinsically
defined obstruction (see paragraphs i-, iii- and iv- in section
\ref{examples} for the precise relations between the two notions).



% Moreover, generically there are no tangencies of periodic points, yet by the results of Newhouse (see \cite{N}) there exist open sets of surface diffeomorphisms whose non-wandering set is not hyperbolic (and does not admit a dominated splitting) and which therefore must contain {\it critical points}. Moreover, under strong conditions on the critical point, one can associate to its critical direction a stable and a (center-)unstable manifold that are tangent, so in that sense one recovers the notion of tangency but not necessarily for periodic points.
 
The relation between tangencies and critical points is made precise in section \ref{examples}: lemma \ref{lem:tang-crit} shows that associated with any tangency there is a critical point whose critical direction coincides with the tangent line to the stable and unstable manifolds at the tangency point. In the reverse direction, non-recurrent critical points (see remark \ref{rm:crit-dom} and section \ref{sec:misiu}) are related to tangencies between stable and (center-)unstable manifolds of points in sets with dominated splitting.


As in the one-dimensional setting, one can impose conditions on critical points to gain a more detailed understanding of the non-wandering set, for instance, obtaining a finite decomposition into transitive components or ensuring the existence of SRB (Sinai–Ruelle–Bowen) measures. A notable class of such systems is given by Misiurewicz maps, where the critical points are non-recurrent. These maps exhibit a finite number of transitive pieces and often display rich chaotic behavior, including the presence of  SRB measures. In the two-dimensional context, it is natural to expect that a suitably defined notion of criticality could play an analogous role. In this paper, we explore this possibility by introducing and analyzing a class of systems we refer to as two-dimensional Misiurewicz diffeomorphisms, where a version of non-recurrence of the critical set enables similar decomposition into a finite number of transitive pieces.

That approach opens the possibility to consider other settings where a more general condition on the critical set could have strong dynamical implications. This perspective is  carried out in \cite{CLPY} where a strong condition on the critical sets for infinite renormalizable dissipative diffeomorphisms  allows one to obtain a general description of the dynamics (see paragraph v- in section \ref{examples}).


% The goal is to define a critical set for surface diffeomorphisms, analogous to critical points in one-dimensional dynamics, as obstructions to hyperbolicity.

% In 1D, the presence of critical points inside the limit set prevents hyperbolicity; their absence implies the dynamics are essentially hyperbolic.

% In 2D, hyperbolicity is obstructed by the failure of dominated splitting. Thus, defining a critical set amounts to identifying the points where domination fails.

% The proposed definition is simpler and more intrinsic than previous ones, and it generalizes to contexts like complex dynamics.

% Critical points correspond to directions that are neither contracted forward nor backward. Their absence implies domination, while tangencies are a related but weaker obstruction.
% By imposing conditions on critical sets (e.g., non-recurrence), one can classify dynamics more precisely, paralleling the role of Misiurewicz maps in 1D.

% This framework extends to other settings, allowing structural descriptions of dynamics under strong assumptions on critical sets.











\bigskip

\paragraph{a -- The critical set.}


Consider a compact boundaryless riemannian surface $M$ and a $C^1$-diffeomorphism $f$.
The derivative of $f$  induces a dynamics  on the unit tangent bundle,  $G_x\colon T^1_xM \to T^1_{f(x)}M$,  usually called the {\it projective tangent bundle map}, 
 defined by $$G_x(v)=
\frac{D_xf(v)}{|D_xf(v)|},$$ and the family $\{G_x\}_{x\in M}$ induces a projective cocycle over $f$.  Let  $g$ be the derivative of $G$ along the fibers
$T_x^1M$ of the unit tangent bundle; more precisely, 
$g_x(v)= |D_wG_x|_{w=v}|.$
We denote the composition $G^n_x= G_{f^{n-1}(x)}\circ\dots \circ G_{x}$ and
to simplify the notation, we denote $D_wG^n_x$ by  $g^n_x$, abbreviated as $g^n$ when no confusion arises. Observe that $g^n_x(v)= \Pi g_{f^i(x)}(G^i_x(v)).$ It is easy to show that 

\begin{eqnarray}\label{eq_f-g}
    |g^n(v)|=  \frac{|\det(Df^n)|}{|Df^n(v)|^2}.
\end{eqnarray}
Let $\Lambda\subset M$ be an invariant compact set. That set  may contain attracting, repelling or elliptic  periodic points. Unlike  in \cite{PRH}, it is not assumed that $f_{|\Lambda}$ is  dissipative or that  $\La$ is contained in the limit set of $f$.


\begin{definition}\label{def:crit} Let $\Lambda\subset M$ be an invariant compact set of  a surface diffeomorphism $f.$ It is said that $x$ is a {\it critical point} of $f$ in $\La$ if there exists $v\in T_xM$ such that 
$|g^n(v)|\geq 1$ for all $n\in \ZZ$. The direction $v$ is called the critical direction and 
 $${\rm Crit}(f,\Lambda)=\{x\in\Lambda, \; \exists v\in T_x^1M,\;
\forall n\in \ZZ, \; |g^n(v)|\geq 1\},$$ is called the critical set of $f$ in $\Lambda.$ Throughout the paper, we refer to elements of ${\rm Crit}(f,\Lambda)$ as
critical points.
\end{definition}

% In the end of section \ref{proof-existence}  is discussed the relation between the definition of critical points as provided in \cite{PRH} and the one develpoed here.

In short,  a critical point is a point having a direction $v$ such that $g^n(v)$ is bounded from below for all integers.  By equation (\ref{eq_f-g}), the condition $|g^n(v)|\geq 1$ for all $n\in\ZZ$ is
equivalent to
$$\|Df^n(v)\|^2\;\leq\; |\det(D_xf^n)|, \qquad \forall\, n\in \ZZ.$$
In other words, a critical direction is a direction that is contracted at
least as strongly as the conformal rate $|\det(D_xf^n)|^{1/2}$ (the geometric
mean of the singular values of $D_xf^n$) in \emph{both} time directions: $v$
never carries the dominant expansion, neither in the future nor in the past.
In the dissipative case, where $|\det(D_xf)|\leq b<1$ on $\La$, the forward
inequality gives
$$\|Df^n(v)\|\leq b^{n/2},\qquad n\geq 0,$$
so the critical direction is exponentially contracted in the future, while
for $n\leq 0$ the inequality forces the complementary directions to carry an
expansion of order at least $|\det(D_xf^{n})|^{1/2}$. This is precisely the
behavior of the common tangent direction at a point of tangency between a
stable and an unstable manifold, contracted in the future, dominated by a
transverse expanded direction in the past, and it is the two-dimensional
counterpart of the one-dimensional picture, where the orbit of a critical
point is decoupled from any expansion of the derivative.
 
% \begin{definition} Let $f\in Diff^1(M)$ and $\La$ be a  compact invariant set of $f$. It is said that $f_{|\La}$ is {\it far from homotheties}, if there exists $\de>0$ such that  there is no point $x$ and vector $v\in T^1_xM$ satisfying 
% \begin{equation}\label{FFH}
%     (1-\de)^n< g^n(v)<(1+\de)^n, \,\,\,\forall n\geq 0.
% \end{equation}


% \end{definition}

\begin{definition}\label{def:FFH} Let $f\in Diff^1(M)$ and let $\La$ be a
compact invariant set of $f$. It is said that $f_{|\La}$ is {\it far from
homotheties} if there exists $\de>0$ such that there is no point $x$ and
vector $v\in T^1_xM$ satisfying
\begin{equation}\label{FFH}
(1-\de)^n< g^n(v)<(1+\de)^n, \qquad \forall\, n\geq 0.
\end{equation}
\end{definition}



Observe that if (\ref{FFH}) does not hold for any point $x$ and any vector $v$ then obviously also holds for any $\de'<\de$.
% previous definition is equivalently to the following: {\it every power $f^M_{|\La}$, $M\geq 1$, satisfies the condition 
% for $N=1$, whose projective derivative cocycle denoted as $g_M$ coincides with  $g^M$.}

By compactness, it is equivalent to say that  there is a positive integer $N_0$ and $\de$ such that for any $x\in \La$ and $v\in T^1_xM$ there is a positive integer $n\leq N_0$ such that either $g^{n}(v)\geq (1+\de)^n$ or $g^{n}(v)\leq (1-\de)^n.$

% \begin{remark}\label{rmk:FFH-powers}
% The property of definition \ref{def:FFH} is invariant under taking powers:
% $f_{|\La}$ is far from homotheties if and only if $f^{N}_{|\La}$ is, for
% any $N\geq 1$. Indeed, the powers of $f^N$ form a subset of the powers of
% $f$, which gives one implication; for the other, given $m\geq 1$, if some
% $(x,v)$ satisfies (\ref{FFH}) for $f^m$ with constant $\de$, then
% restricting to the times multiple of $N$ the same $(x,v)$ satisfies
% (\ref{FFH}) for $f^{mN}$ with rates $(1\pm\de)^{N}$, which are arbitrarily
% close to $1$ as $\de\to 0$; hence the condition for $f^{mN}$ (available
% since $mN$ is a power of $f^N$) yields the condition for $f^m$.
% \end{remark} consider the situation that the condition above rules out for the case $N=1.$:

 To motivate the notion of  ``far from homotheties''  suppose that for all $n \geq 0$ and all $v \in T^1_xM$ one has $(1-\delta)^n < g^n(v) < (1+\delta)^n$. Since $g^n(v) = |\det(Df^n)|/\|Df^n(v)\|^2$, dividing the window conditions for two vectors $v$ and $w$ gives
$$
\left(\frac{1-\delta}{1+\delta}\right)^{n/2} <
\frac{\|Df^n(v)\|}{\|Df^n(w)\|} <
\left(\frac{1+\delta}{1-\delta}\right)^{n/2},
$$
so that all directions are scaled by $Df^n$ comparably, with discrepancy at most an exponential factor in $\delta$. Taking $\delta\to 0$, the ratio $\|Df^n(v)\|/\|Df^n(w)\|$ is forced to $1$ for all pairs of directions and all iterates, an homothety for all iterates.  Being $\delta$-far from homotheties rules this out uniformly: for every $x \in \Lambda$ there must exist some direction $v \in T^1_xM$ and some $n \in \mathbb{Z}$ such that $g^n(v)$ exits the window $((1-\delta)^n,(1+\delta)^n)$, meaning the derivative cocycle must exhibit a definite discrepancy in the scaling of different directions at some iterate.



Recall that a compact invariant set $\La$ has a dominated splitting if there exists a subbundle decomposition $T_\Lambda M = E \oplus F$ with $\dim(E) = \dim(F) = 1$ and a positive constant $\la < 1$ such that $\frac{\|Df_{|E}\|}{\|Df_{|F}\|} < \la$; the subbundle $E$ is called the center-stable subbundle and $F$ the center-unstable subbundle.
 
In terms of the projective dynamics, this is equivalent to saying that $E$ is a repeller and $F$ is an attractor under $g$, i.e. $g(E) > 1 > g(F)$, or equivalently that for any vector $v$ either $g^n(v) \to 0$ or $g^{-n}(v) \to 0$ as $n \to +\infty$; equivalently again, $\La$ admits a contracting and a repelling cone field. The next result establishes a tight relation between the existence of critical points and dominated splitting.
 
\begin{theorem}\label{p.existence}
Let $f\in Diff^1(M)$  and $\La$ be a  compact invariant set of $f$. If  $f^N_{|\La}$ is far from homotheties for all $N\geq 1,$ it follows that $\La$ admits a dominated splitting $T_\Lambda=E\oplus F$ with $\dim(E)=\dim(F)=1$ if and only if  ${\rm Crit}(f,\Lambda)$ is empty.

\end{theorem}


Combined with the main theorem of \cite{PS1}, the previous result yields a
two-dimensional version of Ma\~n\'e's theorem \cite{M} quoted at the beginning
of the introduction: in the absence of critical points, the dynamics is
hyperbolic.

\begin{corollary}\label{cor:mane2d}
For a generic $f\in Diff^2(M)$ the following holds: any compact invariant set
$\La$ such that $f^N_{|\La}$ is far from homotheties for all $N\geq 1,$ is a hyperbolic set if and
only if ${\rm Crit}(f,\La)=\emptyset$.
\end{corollary}

This dichotomy orients the whole paper. When ${\rm Crit}(f,\La)=\emptyset$
the set $\La$ admits a dominated splitting and is, in the generic $C^2$
setting, hyperbolic (corollary \ref{cor:mane2d}); this is the
``essentially hyperbolic'' regime, the two-dimensional counterpart of a
one-dimensional map without critical points. The substance of the theory
lies in the opposite regime, ${\rm Crit}(f,\La)\neq\emptyset$, where
domination fails and the critical set records precisely how: the remainder
of the paper studies what its presence, and conditions on its dynamics,
impose on the global behavior of $f$.


If one considers $C^2-$diffeomorphisms, not necessarily generic, it follows from \cite{PS2} that for any compact invariant set with a dominated splitting,  the limit set  can be decomposed into a finite number of normally hyperbolic closed curves with dynamics conjugate to irrational rotation, finite number of normally hyperbolic periodic intervals and compact sets conjugate to a hyperbolic ones such that the period of non-hyperbolic periodic points is bounded. For a detailed discussion of results about dominated splitting needed in this paper, see section \ref{sec:preliminaries}.  

\begin{remark}\label{rmk:why-powers}
Observe that the hypothesis of theorem \ref{p.existence} requires the far from
homotheties property not only for $f$ but for all its powers. The main reason is that  condition (\ref{FFH}) for $f$ alone does not pass
to the iterates. Consider an attracting periodic point $q$ of period $n$
whose derivative at the period $D_qf^{n}$ has its two eigenvalues of equal
modulus (a homothety up to a rotation; sufficient dissipation only forces
$|\det D_qf^{n}|\leq b^{n}$ and does not exclude this). Along the period the
individual derivatives $D_{f^j(q)}f$ need not be conformal, so $q$ may
satisfy (\ref{FFH}) for $f$ with some $\de$, while at scale $n$ one has
$g^{nk}(v)=1$ for every $v$ and every $k$, so the analogue of (\ref{FFH})
for $f^{n}$ fails at $q$ for every $\de$. A periodic domain, in
particular one of large period, can thus be far from homotheties for $f$
yet not for an iterate, and the obstruction to domination it carries is
visible only at the coarse time scale: this is exactly what the all-powers
hypothesis captures.
\end{remark}


One may ask whether the existence of a critical point in an invariant set is a property preserved under conjugacy. This is not the case for $C^0-$conjugacies: one can have two $C^0-$conjugate surface diffeomorphisms such that the invariant set of one of them is a hyperbolic set (and therefore does not contain critical points) and the other exhibiting a (cubic) tangency (see paragraph ii- in section \ref{examples} for more details). However, having critical points is a property invariant by $C^1-$conjugacies, since  dominated splittings are preserved by diffeomorphisms. 

One may further ask whether the notion of critical point is independent of the choice of metric; to address that issue, one could define  $x$ to be  a \emph{critical point}
if there exists $v\in T^1_xM$ and $K>0$ such that
$\forall n\in \ZZ, \; |g^n(v)|\geq K.$ The threshold $|g^n(v)| \geq 1$ depends on the choice of Riemannian metric, but the critical set $\mathrm{Crit}(f,\Lambda)$ does not. Since $M$ is compact, any two smooth Riemannian metrics are Lipschitz equivalent: there exists $C > 1$ such that $C^{-1}\|w\| \leq \|w\|' \leq C\|w\|$ for all tangent vectors $w$. As $g^n(v) = |\det(Df^n)|/|Df^n(v)|^2$, passing from one metric to the other multiplies $g^n(v)$ by a factor uniformly bounded between $C^{-2}$ and $C^2$, independently of $n$. Hence $|g^n(v)| \geq 1$ for all $n \in \mathbb{Z}$ in one metric implies $|g^n(v)| \geq C^{-2} > 0$ for all $n \in \mathbb{Z}$ in the other, and $x \in \mathrm{Crit}(f,\Lambda)$ with respect to one metric if and only if it belongs to $\mathrm{Crit}(f,\Lambda)$ with respect to the other. Throughout the paper we fix an arbitrary smooth metric and adopt the threshold $K=1$ for notational convenience.



% It follows immediately now that the notion of critical point is independent of the metric and, 
% similarly, one can adapt the proof of theorem \ref{p.existence} with the above definition of critical point.



%OLD paragraph b

% \paragraph{b-- Hypothesis that imply that  $f_{|\La}$ is far from homotheties.}



% The main result of this section asserts that, under a quantitative form of dissipation that we will make precise below, $f_{|\La}$ is far from homotheties off the basins of finitely many attracting periodic points. The underlying mechanism is the following. If a direction $v$ at a point $x$ satisfies inequalities (\ref{FFH}),
% then by (\ref{eq_f-g}), $|Df^n(v)|^2= |\det(Df^n)|/g^n(v)$ is squeezed between $|\det(Df^n)|/(1+\de)^n$ and $|\det(Df^n)|/(1-\de)^n$. When $f$ is sufficiently dissipative, that asymptotic close to homotheties, forces the orbit into the basin of an attracting periodic point whose local basin contains a ball of definite size. Compactness of $M$ then bounds the number of such attractors, and the theorem follows. The condition we now introduce is the explicit numerical form of``sufficiently dissipation" that makes this argument go through.




% \begin{definition} Let $f\in \diff^1(M^2)$ and let $\La$ be a compact invariant set such that $f_{\La}$ is dissipative. If $f$ satisfies  
% \begin{equation}\label{cond-diss-FFH}
%    2\sqrt{b}\,M_0 + 2b + \frac{2b^{5/2}}{M_0^3} <1,
% \end{equation}where $b=\max_{x\in \La}|det(D_xf)|$, $M_0=\max_{x\in \La}||D_xf||$, we  say that {\it $f$ is sufficiently dissipative}.
    
% \end{definition}

% It was proved in \cite{PRH} that if  $f_{|\La}$ is dissipative and  $\La$ does not contain periodic orbits, then $f_{|\La}$ is far from homotheties (see main lemma in \cite{PRH}). A similar statement can be obtained assuming that the dynamic is area expanding and there are no repellers in $\La$.
%A more general result is the following


% \begin{theorem}\label{measure-FFH}
    
%  Let $f\in Diff^1(M)$ and let $\La$ be a compact invariant set. If there exists $\ga>0$ such that for any invariant measure, the difference of the Lyapunov exponents at any regular points is larger than $\gamma$, then $f_{|\La}$ is far from homotheties. 

% \end{theorem} 
%   The previous result can be reformulated for  the case of dissipative dynamics, meaning that  there is  $0<b<1$ such that for any $x\in \La$ it holds  $|det(D_xf)|<b$. 


  
% \begin{theorem}\label{dissipative-FFH} 
% Let $f\in \diff^1(M^2)$ and let $\La$ be a compact invariant set such that $f_{|\La}$ is sufficiently dissipative, then  there exists a finite number of attracting periodic points such that $f$ restricted to the intersection of $\La$ with the complement of the basins of attraction of those attracting periodic points is  far from homotheties.  

% \end{theorem} 





% The idea of the proof is that for $\de$ sufficiently small  and $N$ sufficiently large there is $x$ and a vector $v$  such that inequalities (\ref{FFH}) holds for any positive integer up to $N$ then $x$ is in the basin of attraction of an attracting  periodic point that its local basin of attraction is uniformly large; since in a compact manifold there are at most finite number of attracting periodic points satisfying that their  local basin of attraction is uniformly large, the theorem follows. 
% In short, the  suficiently dissipativity condition 
% forces $g^n(v) = |\det(Df^n)|/\|Df^n(v)\|^2$ to exit the homothety window unless
% the orbit is being exponentially contracted, which forces it into the basin of
% attraction of an attracting periodic point.


% We believe that the sufficiently  dissipative condition needed for theorem \ref{dissipative-FFH} can be relaxed to just dissipation; this depends on a more general version of  proposition \ref{prop:pinning}. This is discussed in subsection \ref{subsec:relaxing}. 

% In what follows, we discuss the sufficiently condition for the case of H\'enon maps of the plane.

% \begin{remark}\label{Henon-diss}{\bf H\'enon maps and attracting domains. } Let $f_{\mu,b}$ be the classical H\'enon family of diffeomorphisms of the   real plane $$f_{\mu, b}(x,y)=(x^2+\mu+b.y,x).$$ For dissipative H\'enon maps (i.e., $|b|<1$), there exists a connected open disk  $P$ of the parameter space $\{(\mu,b): |b|<1,\, \mu\in \RR\}$ such that for each $(b, \mu)\in P$ it follows that there is a disk $\DD_{\mu,b}$ such that:

% \begin{itemize}
%     \item[--] $f_{\mu,b}(\DD_{\mu, b})\subset \DD_{\mu,b},$
%     \item for any $(x,y)\in \RR^2$ either there is a forward iterate that is contained in $\DD_{\mu,b}$ or the forward iterates converge to $-\infty.$ 
% \end{itemize}
    
% \end{remark}


% \begin{remark}{\bf H\'enon maps and far from homotheties. } If $\sqrt{|b|}< 0.2247,$ it follows that ${f_{\mu,b}}_{|D_{\mu,b}}$ satisfies that is sufficiently dissipative.
    
% \end{remark}

% Theorems \ref{measure-FFH} and \ref{dissipative-FFH} provide two natural sufficient conditions for being far from
% homotheties, covering complementary situations. Theorem \ref{measure-FFH} addresses the general
% case: a uniform Lyapunov exponent gap prevents any direction from having its
% projective derivative $g^n(v)$ trapped in the homothety window
% $((1-\delta)^n,(1+\delta)^n)$ for all $n \in \mathbb{Z}$, since by Oseledets theorem
% the asymptotic exponential growth rate of $g^n(v)$ is controlled by the difference
% of Lyapunov exponents.





% One can apply previous lemma to the critical point:

% \begin{corollary}Let $f\in Diff^1(M)$ and let $\La$ be a compact invariant and $\ga>0$ such that for any invariant measure, the difference of the lyapunov exponents is larger than $\ga$. Then, for any $x$, $v$, $\eps>0$ and  $N$ there exists $n>N$ such that either $g^n(v)< e^{n(-\ga+\eps)}$ or $g^n(v)> e^{n(\ga-\eps)}.$

% In partiular, if  $(x, v)$ a critical point,  then there exist two infinite sequences of integers $(n_i^+), (n_i^-)$ such that 
% \begin{enumerate}

% \item[--]$g^{n_i^+}(v)> e^{\ga.n_i^+}$,

% \item[--]  $g^{-n_i^+}(v)> e^{\ga.n_i^+}$.


% \end{enumerate}

 

% \end{corollary}


%It is natural to wonder if the condition on critical point can be improved, meaning that the growth $(g^n(v))_{n\in \ZZ}$ is exponential. Due to the next proposition, this is not true in general.



%they are related to what holds if a critical point is periodic. This can be done even in the case that the dynamics is not dissipative.

% Revised paragraph b- 
% Changes:
%  (i) grammar fixes in the new opener ("that asymptotic close to homotheties", "sufficiently dissipation")
%  (ii) trimmed the post-theorem redundancy: the "idea of proof" + "in short" pair is replaced by one short paragraph
%  (iii) reworded the "relaxing" sentence to be more active
%  (iv) added a label to the second Hénon remark
%  (v) minor grammar fixes ("sufficiently condition" -> "sufficient dissipation condition", etc.)


% Restructured paragraph a (Option 3).
% Changes:
%  (i) The geometric-content passage is set off as a numbered Remark
%      rmk:homothety-window, subdivided internally with italic lead-in phrases.
%  (ii) The metric-independence passage is demoted from a numbered Remark
%       to a free-standing \paragraph block (no label, since it is not
%       intended to be cited from elsewhere), also subdivided internally
%       with italic lead-in phrases.
%  (iii) The conjugacy paragraph remains as standalone running prose.
%  (iv) No mathematical content changed; no content deleted.


\paragraph{b-- Hypothesis that imply that  $f_{|\La}$ is far from homotheties.}

The main result of this section asserts that, under a quantitative form of dissipation that we will make precise below, any iterate  $f^N_{|\La}$ is far from homotheties off the basins of any attracting periodic points.  The mechanism (for the particular case of $N=1$) is the following: if a direction $v$ at a point $x$ satisfies inequalities (\ref{FFH}), then by (\ref{eq_f-g}), $|Df^n(v)|^2= |\det(Df^n)|/g^n(v)$ is squeezed between $|\det(Df^n)|/(1+\de)^n$ and $|\det(Df^n)|/(1-\de)^n$. When $f$ is sufficiently dissipative, this asymptotic proximity to homotheties forces the orbit into the basin of an attracting periodic point whose local basin contains a ball of definite size. Compactness of $M$ then bounds the number of such attractors, and it  follows that $f$ is far from homotheties. The condition we now introduce is the explicit numerical form of sufficient dissipation that makes this argument go through. A similar argument holds for every iterate of $f$, and therefore removing all the basins of attraction of all sinks, the result follows.


\begin{definition} Let $f\in \diff^1(M^2)$ and let $\La$ be a compact invariant set such that $f_{\La}$ is dissipative. If $f$ satisfies  
\begin{equation}\label{cond-diss-FFH}
   \sqrt{b}\,M_0 + 2b  <1/2,
\end{equation}where $b=\max_{x\in \La}|det(D_xf)|$, $M_0=\max_{x\in \La}||D_xf||$, we  say that {\it $f$ is sufficiently dissipative}.
    
\end{definition}


\begin{theorem}\label{dissipative-FFH} 
Let $f\in \diff^1(M^2)$ and let $\La$ be a compact invariant set such that
$f_{|\La}$ is sufficiently dissipative. Then any power of  $f$ restricted to the set of
points of $\La$ that do not belong to the basin of attraction of any
attracting periodic orbits is  far from homotheties. 
\end{theorem} 

% %Moreover, for each
% $N\geq 1$ there is a finite number of attracting periodic orbits such that
% the condition (\ref{FFH}) at scale $N$ already holds off the basins of
% those finitely many orbits. 

The proof proceeds by a finite-time version of the argument sketched above, whose main ingredient is proposition \ref{prop:pinning}. The sufficient dissipation hypothesis can plausibly be weakened to mere dissipation; this would require a stronger version of proposition \ref{prop:pinning} and is discussed in subsection \ref{subsec:relaxing}. 

We now illustrate the sufficient dissipation condition for the  H\'enon maps of the plane.

\begin{remark}\label{Henon-diss}{\bf H\'enon maps and attracting domains. } Let $f_{\mu,b}$ be the classical H\'enon family of diffeomorphisms of the   real plane $$f_{\mu, b}(x,y)=(x^2+\mu+b.y,x).$$ For dissipative H\'enon maps (i.e., $|b|<1$), there exists a connected open disk  $P$ of the parameter space $\{(\mu,b): |b|<1,\, \mu\in \RR\}$ such that for each $(b, \mu)\in P$ it follows that there is a disk $\DD_{\mu,b}$ such that:

\begin{itemize}
    \item[--] $f_{\mu,b}(\DD_{\mu, b})\subset \DD_{\mu,b},$
    \item for any $(x,y)\in \RR^2$ either there is a forward iterate that is contained in $\DD_{\mu,b}$ or the forward iterates converge to $-\infty.$ 
\end{itemize}
    
\end{remark}


\begin{remark}\label{Henon-FFH}{\bf H\'enon maps and far from homotheties. } If $\sqrt{|b|}< \frac{\sqrt2-1}{2}\approx 0.2071,$ it follows that ${f_{\mu,b}}_{|D_{\mu,b}}$ is sufficiently dissipative.
    
\end{remark}

 
\paragraph{c-- Properties of the critical set.} In what follows we collect several properties of the critical set. The first group concerns recurrence behavior, paralleling classical results on critical points in one-dimensional dynamics. The second concerns the dependence of the critical set on the diffeomorphism under small perturbations.



 

\begin{proposition}\label{properties}
Let $f\in Diff^1(M^2)$ and let $\La$ be a compact invariant set, the following properties hold
 \begin{itemize}



    % \item[--]  If $f_{|\La}$  is far from homotheties, then the critical direction associated with a critical point is unique; moreover, the critical directions of critical points vary continuously. 


  
  \item[--] If a critical point is periodic, then the derivative at the period has only one eigenvalue: the derivative at the period is either a homothety, elliptic, or parabolic (up to scalar multiplication). In particular, in the dissipative context, the periodic point is an attracting periodic point.

\item[--]  If $f_{/\La}$  is far from homotheties, then it cannot hold that all iterates of the critical point and the critical direction  are  critical points and critical directions.


  

  
    \item[--] The critical set is closed and depends semi-continuously on $f$: for any  neighborhood $U$ of ${\rm Crit}(f,\Lambda)$ there exists a neighborhood $V$ of $\La$ such that if $g$ is sufficiently close to $f$ and $\La'$ is a compact invariant set of $g$ contained in $V$ then ${\rm Crit}(g,\Lambda') \subset U.$ 
  
   

   
  

    
 \end{itemize}





\end{proposition}


The second item of the proposition for dissipative dynamics is similar to  an obvious result for one-dimensional endomorphisms: if a critical point is periodic, then it is an attracting periodic point. In the dissipative context, the proof proceeds by considering the return of a critical direction under the iterate that makes the critical point periodic: if the direction is invariant, then the periodic point is an attracting periodic point such that the derivative has   either one or two invariant directions (with the same rate of contraction); if the direction is not invariant, then the derivative has a complex eigenvalue with modulus smaller than one.  

The following remark and proposition  are about relating the critical direction of a critical point with the subbundles of a set with dominated splitting that exists in the complement of the critical set. The next remark, describes this relation under the assumption that the critical point is not forward (backward) recurrent.
To be precise,  let us recall first that a dominated splitting $E\oplus F$ of a compact invariant set $\La$ can be extended to an invariant dominated splitting $\tilde E\oplus \tilde F$ on a small neighborhood $U$ of $\La$ in the sense that if $x\in U$ and $f(x)\in U$ then $D_f(\tilde E_x)= \tilde E_{f(x)}$ and $D_f(\tilde F_x)= \tilde F_{f(x)}$. The extension $\tilde F$ is unique for the points that all its backward iterates remain in $U$ and similarly the extension $\tilde E$ is unique for the points that all its forward iterates remain in $U$. Moreover, if $U$ is a sufficiently small neighborhood it holds that $\tilde F (\tilde E)$ is close to $F$ in the sense that given $x\in U$ there is $x'\in \La$ such that the slope between $\tilde F_x$ and $F_{x'}$ is small (respectively with $\tilde E_x$ and $\tilde E_{x'}$). 





\begin{remark}\label{rm:crit-dom} Given a critical point $c$, if the  $\omega-$limit of $c$ (denoted as $\omega(c)$) is disjoint from the critical set, then $c$ belongs to the center-stable manifold of $\omega(c)$ and the forward iterates of the critical direction coincide with $\tilde E$, which is uniquely defined on the forward orbit of $c$. If  the $\alpha-$limit of $c$  (denoted as $\alpha(c)$) is disjoint from the critical set, then $c$ belongs to the center-unstable manifold of $\alpha(c)$ and the backward iterates of the critical direction coincide with $\tilde F$ (uniquely defined on the backward orbit of $c$). 
    
\end{remark}
This follows immediately from the fact that if the critical direction does not coincide with the center-unstable subbundle $\tilde F$ defined in a neighborhood of $\La$,  then the backward iterates converge to $\tilde E$ and therefore $g^{-n}(v)\to 0,$ contradicting the definition of critical direction. Under the non-recurrence assumption above, a critical point behaves as the point in the orbit of tangency where the unstable and stable directions are interchanged (see paragraph i- of \ref{examples}).

A more general statement, which allows for recurrence of critical points to the critical set, is the following compatibility result between the subbundles of a dominated splitting and the critical directions of a critical point.


\begin{proposition}\label{pr:compatibility} Let $f\in \mathrm{Diff}^1(M^2)$ and let $\Lambda$ be a compact invariant set. Suppose there exists $\delta_0>0$ such that $\Lambda_{\delta_0}$ is non-empty, where $\Lambda_{\delta_0}=\bigcap_{n\in \mathbb{Z}} f^n\!\left(\Lambda\setminus B_{\delta_0}(\mathrm{Crit}(f, \Lambda))\right)$ .

Then, for any $\varepsilon>0$ and $\delta<\delta_0$, there exists a positive integer $N_0$ such that the following holds. If $c$ is a critical point satisfying $f^{-j}(c)\notin B_{\delta}(\mathrm{Crit}(f,\Lambda))$ for all $0<j\leq n$ with $n>N_0$, then there exists a positive integer $n^*\leq n$ such that:
\begin{itemize}
    \item[--] $f^{-n^*}(c)\in B_\varepsilon(\Lambda_{\delta})$, and
    \item[--]  $\mathrm{slope}(G^{-n^*}(v), F)< \varepsilon$, where $F$ is the center-unstable subbundle of $\Lambda_{\delta}$.
\end{itemize}
Similarly, if $c$ is a critical point satisfying $f^{j}(c)\notin B_\delta(\mathrm{Crit}(f,\Lambda))$ for all $0< j\leq n$ with $n>N_0$, then there exists a positive integer $n^*\leq n$ such that:
\begin{itemize}
    \item[--] $f^{n^*}(c)\in B_\varepsilon(\Lambda_\delta)$, and
    \item[--]   $\mathrm{slope}(G^{n^*}(v), E)< \varepsilon$, where $E$ is the center-stable subbundle of $\Lambda_\delta$.
\end{itemize}




\end{proposition}




In short, the previous proposition states that the critical direction behaves like a center-unstable direction when iterated backward (provided the critical point takes a long time to return backward to a neighborhood of the critical set), and like a center-stable direction when iterated forward (provided it takes a long time to return forward).
% Abusing the language and interpretation that can be understood as the critical point being ``either cubic or quadratic" (see section \ref{examples} paragraph v- for some related interpretation).

A simple consequence of that proposition is that for points close to a set with dominated splitting whose forward (backward) iterate approaches a critical point, the forward (backward) iterate of the center-unstable (stable) direction approaches the critical direction of that critical point. 


\begin{corollary}\label{cor:compatibility} 
Let $f\in \mathrm{Diff}^1(M^2)$ and let $\Lambda$ be a compact invariant set. Suppose there exists $\delta_0>0$ such that the set $\Lambda_{\delta_0}=\bigcap_{n\in \mathbb{Z}} f^n\!\left(\Lambda\setminus B_{\delta_0}(\mathrm{Crit}(f, \Lambda))\right)$ is non-empty.

For any $\varepsilon>0$ and $\de_1<\de_0$, there exist a positive integer $N_0$, and constant $\delta_2>0$ such that the following holds. If  $f^{j}(x)\notin B_{\delta_1}(\mathrm{Crit}(f,\Lambda))$ for all $0\leq j\leq n-1$ with $n\geq N_0$, and $f^n(x)\in B_{\delta_2}(\mathrm{Crit}(f,\La))$, then 

\begin{itemize}
    \item[--] there is a positive integer $n^*<n$ such that $f^{n^*}(x)\in B_\e(\La_{\de_1}),$
    \item[--]  
there exists $c\in \mathrm{Crit}(f, \Lambda)$ such that $\mathrm{slope}(G^{n-n^*}(\widetilde{F}_{f^{n^*}(x)}), v_c)<\varepsilon.$
\end{itemize}

Similarly, if  $f^{-j}(x)\notin B_{\delta_1}(\mathrm{Crit}(f,\Lambda))$ for all $0\leq j\leq n-1$ with $n\geq N_0$, and $f^{-n}(x)\in B_{\delta_2}(\mathrm{Crit}(f, \Lambda)$ then 

\begin{itemize}
    \item[--] there is a positive integer $n^*<n$ such that $f^{-n^*}(x)\in B_\e(\La_{\de_1}),$
    \item[--] there exists $c\in \mathrm{Crit}(f, \Lambda)$ such that $\mathrm{slope}(G^{-(n-n^*)}(\widetilde{F}_{f^{-n^*}(x)}), v_c)<\varepsilon.$
\end{itemize}


\end{corollary}





In one dimension, exponential self-recurrence of a critical point forces it into the basin of an attracting periodic point; Proposition \ref{exponential recurrence} below is the two-dimensional analog, valid under no exponential growth of the norm.

\begin{proposition}\label{exponential recurrence}
Let $f$ be a $C^2$ surface diffeomorphism and $\La$ be a dissipative invariant compact set,  satisfying $\underset{k\to +\infty}\limsup \tfrac 1 k \log \|Df^k|_\La\|\leq 0$ which contains a critical points $c$
with exponential recurrence, i.e. $\underset{k\to +\infty}\liminf \tfrac 1 k \log d(f^k(c),c)< 0.$
Then $\La$ contains a sink.
\end{proposition}



\paragraph{d-- Critical points for dissipative dynamics of the disk.}





In what follows, we assume that $\DD$ is an open disk that is properly invariant by $f$, that is $\overline{f(\DD)}\subset \DD$ and $f: \DD\to f(\DD)$ is a diffeomorphism.To simplify notation, we denote this class of maps by $ Diff^r(\DD)$ and observe that any map in that class can be extended to a diffeomorphism of the sphere.

In dimension one, every unstable branch of a repelling periodic point either contains a critical point or lies in a periodic interval on which the dynamics is a diffeomorphism. Theorem \ref{disk2} is the two-dimensional analog, for mildly and sufficiently dissipative diffeomorphisms of the disk; the second option of the 1D statement is replaced by the existence of a normally hyperbolic invariant arc, since the map is already a diffeomorphism.


\begin{definition}
      Given a boundaryless surface $S$ and a dissipative  smooth diffeomorphisms $f:S\to f(S)\subset S$,  it is said that $f$ is {\it mildly dissipative} if for any ergodic measure $\mu$ which is not supported on an attracting periodic point,and   for $\mu$-almost every point, each of its  two stable branches   meets the boundary of $S$.
 \end{definition}

In \cite{CP} it was shown
that mild dissipation is satisfied for large classes of systems. For instance, it holds for some $C^2$
open sets of diffeomorphisms of the disc including diffeomorphisms close to one-dimensional endomorphisms, and also for polynomial automorphisms of $\RR^2$ whose
Jacobian is  close to zero, including the diffeomorphisms from the H\'enon family
with Jacobian of modulus less than $1/4$ (up to restricting to an appropriate trapped disc).  This class captures certain properties of one-dimensional maps, but keeps two-dimensional
features showing all the well-known complexity of dissipative surface diffeomorphisms.
The dynamics of the new class, in some sense, is intermediate between one-dimensional
dynamics and general surface diffeomorphisms.


The two-dimensional result that resembles the one-dimensional stated before  is proved for the closure of an unstable branch; moreover, the second option that holds in the one-dimensional situation does not make much sense, given that the map is already a diffeomorphism. So, that option is replaced by saying that the accumulation set  of the unstable branch is in the basin of attraction of a periodic point or it is an invariant closed arc $\ga$ that is {\it normally hyperbolic}; i.e. $\ga$ has a dominated splitting $E^s\oplus F$ such that $F$ is tangent to the curve $\ga$ and $E^s$ is a contracting direction. To be precise, given a saddle periodic point $p$, which we may assume without loss of generality to be fixed,  the accumulation set of one of its unstable branches (which can also be assumed fixed) is the accumulation set of the union of the iterates of a fundamental domain of that unstable branch. This set is an invariant connected set that contains $p$.





 \begin{theorem}\label{disk2} Let $f\in Diff^2(\DD)$ such that $f_{|\DD}$ is mildly and sufficiently dissipative. Given an unstable branch  $\Ga^u(p)$ of a saddle fixed point $p$ either:

\begin{enumerate}
   

% \item[--] it is a saddle periodic point and the branch coincides with one of the stable branches of the saddle periodic point,

%\item[--] either it is a (semi-)attracting periodic point and the branch is contained in the basin of that periodic point, 
\item[--]  $\Ga^u(p)$ is a normally hyperbolic attracting  invariant arc with one endpoint $p$ and the other is a (semi-)attracting fixed point, 

 \item[--] or there is a critical point either in  $\Ga^u(p)$ or in the accumulation set of $\Ga^u(p)$.
    

    \end{enumerate}

    
    
    
    % the saddle periodic point is contained in a non-trivial homoclinic class, then it follows that the unstable branch that intersects the homoclinic class accumulates on a critical point.


 %or it has dominated splitting such that at least one of the transitive components is a (semi)attracting periodic point.
 
\end{theorem}

Observe that in the previous theorem it is not claimed that the critical point belongs to the limit set. Examples of that situation are considered in  section \ref{examples}, paragraph {\em vi}. %Also, if the unstable branch  intersects the stable manifold, then the unstable branch accumulates at a critical point.

% Other example can be found in \cite{E} where it is shown a surfaces diffeomorphisms in the boundary of Anosov,  exhibiting a cubic tangency associated with a heteroclinic point. 
% \begin{theorem}\label{disk} Let $f\in Diff(\DD)$ such that $f_{|\DD}$ is dissipative. It holds  that  the accumulation set of any unstable branch of a periodic point satisfies that  

% \begin{enumerate}
%     \item[--] either it contains a critical points that is also in the non-wandering (limit) set,

%     \item[--] or it is an attracting set such that the non-wandering set has dominated splitting; moreover it holds that 

%     \begin{enumerate}
%     \item[--] either it contains a  critical point in the stable manifold of set with dominated splitting

% \item[--] or the attracting set is a normally hyperbolic periodic arc. 
    

%     \end{enumerate}

% \end{enumerate}

%  %or it has dominated splitting such that at least one of the transitive components is a (semi)attracting periodic point.
 
% \end{theorem}


%Similar result holds for the stable manifold



%\begin{lemma} If $(x,v)$ is a critical point and $g^n(v)>(1+\delta)^n$ for any $n>0$, it holds that $(x,v)$ is a forward $(\de, \frac{D}{1+\de})-$regular critical point.
 
 
 
%\end{lemma}

%\begin{corollary} If $(\de, \la, m, D)$ verifies, then the set of $(\de, la)-$regular points has a uniform stable manifold.
 
%\end{corollary}








%\paragraph{e-- Mild dissipative diffeomorphisms.}





%\paragraph{f--  Mild dissipative diffeomorphisms of the disk.}






%\paragraph{g-- Regular critical points and the  $\chi-$dissipative case.}


%\begin{definition} Given a smooth dissipative diffeomorphism $f$ and $\chi\in (0,1)$ it is said that it is $\chi-$dissipative if for  $x,y\in S$ and any unit vector $u\in T_yS$,
%\begin{equation}\label{e.dissipation2}
%|\det Df(x)|\leq  \|Df(y).u\|^{\chi}.
%\end{equation}
%\end{definition}


%\begin{theorem}\label{chi critical} Let $f$ be a smooth $\chi-$dissipative diffeomorphism. associated with any regular $(\de,\la)-$critical point, there exists a local stable manifold.
%\end{theorem}


% 
% \begin{theorem} Let $f$ be $\frac 9{10}-$strong dissipative. Them, there is $\e>0$ and a positive integer $n_0$  such that  of any  critical point $x$, the critical direction $v$  satisfies $$g^n(v)>(1+\e)^n,\,\,\ n>n_0.$$ 
%  
% \end{theorem}
% 
% \begin{proof} If not, there is a sequences $(n_i)$ such that $(1-\e)^{n_i}< g^{n_i}(v)< (1+\e)^{n_i}$.  Using the $\frac 9{10}-$condition and Pliss lemma (as stated in \cite{CP}) simultanously for both inequalities; it follows that there is a point $x'$ and a direction $v'$ such $(1-\e)^{n}< g^{n_i}(v)< (1+\e)^{n},$ 
%  for all positive $n$. By inequality (\ref{min contraction}) it follows that for any vector $w$ it holds that $(1-\e)^{n}< g^{n_i}(v)< (1+\e)^{n},$ and since $||Df^n(w)||=\frac{ det(Df^n)}{g^n(w)}$  follows  that $||Df^n||< \lambda^n$ for any positive integer $n$ and some $\lambda<1.$ 
%  
%  
%  
% \end{proof}


\paragraph{f-- Misiurewicz diffeomorphisms.}

A smooth one-dimensional endomorphism is Misiurewicz if all its critical points are non-recurrent; in that case the dynamics is uniformly expanding away from the critical set. We introduce the two-dimensional analog (Definition below) and prove, in Theorem \ref{thm:finite classes}, that for mildly and sufficiently dissipative Misiurewicz diffeomorphisms of the disk the number of generalized homoclinic classes is finite.

% Revisit \cite{PS2}

% From now on given a compact invariant set $\La$ of a surface diffeomorphisms $f$, we assume that the limit set of $f$ restricted to $\La$ is isolated in the limit set of $f.$

\begin{definition}
 Let $\La$ be a compact invariant set of a $C^{2}$ surface diffeomorphism $f$. We say that $f_{|\La}$ is a  {\it Misiurewicz diffeomorphism} if there exists a neighborhood $V$ of the critical set ${\rm Crit}(f, \La)$ such that  any $c\in {\rm Crit}(f, \La)$ satisfies $\al(c)\cap  V=\emp$ and $\omega(c)\cap V =\emp$.
\end{definition}


% Theorems about dominated splitting

% \begin{theorem}\label{misiu1} Let $\La$ be a compact invariant set for a $C^{2}$ surface diffeomorphisms $f$. If $f_{|\La}$ is Misiurewicz then it follows that any critical point belong to the stable  manifold  and to the unstable manifold  of a set exhibiting a dominated splitting. Moreover,  the critical point 
% direction coincides with the tangent at the critical point of the stable and unstable manifold of the dominated set.   
% \end{theorem}



%In what follows, we consider planar manifolds, i.e. that any component of $M$  can be embedded in the two-sphere.

%We  use the following result by Pixton (\cite{Pi}):

\vskip 3pt

% \begin{theoremName}
%  
% \end{theoremName}


% \noi{\bf Pixton's theorem}
% {\it Suppose $M$ is planar, $1 \leq r \leq \infty$. Then there is a residual
% set ${\cal R}\subset  Diff^r(M)$ such that, for $f\in {\cal R}$ follows 
% \begin{enumerate}
%  \item a saddle point is an accumulation point of $\Omega(f)$ if and only if it has a
% transversal homoclinic point;
% \item if $p, q$ are saddle points then  one branch of the unstable manifold of $q$ accumulates on the stable of $p$  if and
% only if the unstable and stable manifold intersects transversely.
% \end{enumerate}}


Differently that in the one dimensional case where the Misiurewicz condition only involves the forward orbit
of the critical points, in the two-dimensional case one has also to consider the backward iterates.   This is not
a technical precaution but reflects the structure of a critical point
described in remark \ref{rm:crit-dom} and proposition \ref{pr:compatibility}:
the critical direction behaves as a center-stable direction in the future and
as a center-unstable direction in the past, so the critical orbit interacts
with the rest of the dynamics through both of its $\alpha-$ and $\omega-$ limits, exactly as a point
of tangency, whose forward orbit follows a stable manifold and whose backward
orbit follows an unstable one. Accordingly, the proofs of section
\ref{sec:misiu} use the condition $\omega(c)\cap V=\emptyset$ and
$\alpha(c)\cap V=\emptyset$ in an entirely symmetric way.
 
The known examples of Misiurewicz maps are described in section \ref{examples}, in paragraph ii. Naturally, one can wonder if these are the only natural classes of examples. More broadly, one can ask how the dynamics on the critical set constrains the global dynamics on the non-wandering set; the natural question is the following: 







\begin{question} \label{misiu3}Let $\La$ be a compact invariant set of a $C^{2}$  surface diffeomorphisms $f$. If  $f_{|\La}$ is   Misiurewicz does it follow that $\Lambda$ can be decomposed into a finite number of transitive sets and a finite number of normally hyperbolic arcs?
 
\end{question}

If one  assumes that all the periodic points are hyperbolic, and  if the previous question is answered positively, then the dynamics  decomposes into finitely many transitive pieces. If one allows non-hyperbolicity,  one can have infinitely many  transitive isolated sets; for instance,  a normally hyperbolic arc such that the restriction to that arc is a one-dimensional diffeomorphism exhibiting infinitely many attracting fixed points (as in the case of the map $x\to x^r\sin(\frac 1{x})$). 

In the following theorem \ref{thm:finite classes}, we provide  a partial answer  to question \ref{misiu3} for mildly dissipative diffeomorphisms of the disk  focusing on non-trivial {\it homoclinic classes} instead of general transitive sets. In this paper, a homoclinic class is the closure of the set of  {\it crossing} points of the stable and unstable manifolds of a saddle periodic point (see definition \ref{defi:crossing}).    In the classical definition, only transversal intersections are considered: an intersection is transversal when the tangent spaces of the stable and unstable manifolds at the intersection point are not collinear. We instead allow crossing intersections, meaning that the unstable manifold meets both connected components of any local-stable ``band'' around the intersection point and for brevity we keep calling  these classes as homoclinic classes (see definition \ref{defi:generalizedHC} in section \ref{sec:misiu} for details) .


\begin{theorem}\label{thm:finite classes} Let $f$ be a mildly and sufficiently dissipative     $C^2-$diffeomor\-phism of the disk, and let
$\La=\bigcap_{n\geq 0}f^n(\DD)$ be its maximal invariant set, which is
compact and invariant. Suppose that  $f_{|\La}$ is Misiurewicz. Then, the number of non-trivial  homoclinic classes is finite.

Moreover, there exists a positive integer $n_0$ such that any saddle periodic point  with period larger than $n_0$ not contained in any of the finitely many non-trivial homoclinic classes,  belongs to  a non-trivial  normally hyperbolic periodic arc.
    
\end{theorem}

 Observe that there could exist infinitely many saddle periodic orbits that do not belong to a non-trivial homoclinic class and there  could also exist infinitely many normally hyperbolic arcs. 

We note that in the one-dimensional case, Misiurewicz maps are not $C^2$-open:
they are typically accumulated by maps for which the critical point is recurrent
but not periodic. The analogous question for surface diffeomorphisms is beyond
the scope of the present paper. Here our goal is more modest: to show that the
non-recurrence condition on the critical set has strong dynamical consequences,
namely the finiteness of the number of homoclinic classes, in the spirit of the
one-dimensional case where non-recurrence of critical points implies a finite
decomposition into transitive pieces.



\paragraph{A dictionary.} The following table summarizes how the results of
this paper transpose the classical role of critical points in one-dimensional
dynamics to surface diffeomorphisms.

\begin{center}
\begin{tabular}{p{0.42\textwidth}|p{0.42\textwidth}}
{\bf One-dimensional endos} & {\bf Surface diffeomorphisms}\\
\hline\hline
critical point: $f'(c)=0$ &
critical point: definition \ref{def:crit}, via the projective cocycle
\\
\hline
no critical points $\Rightarrow$ essentially hyperbolic (Ma\~n\'e \cite{M}) &
theorem \ref{p.existence}; corollary \ref{cor:mane2d} (generic $C^2$,
far from homotheties)\\
\hline
a periodic critical point is attracting &
proposition \ref{properties}, first item\\
\hline
exponentially recurrent critical point lies in the basin of an attracting
periodic orbit &
proposition \ref{exponential recurrence}\\
\hline
an unstable branch of a repelling periodic point contains a critical point,
or lies in a periodic interval where the map is a diffeomorphism &
theorem \ref{disk2}: critical point in the closure of the branch, or
normally hyperbolic invariant arc\\
\hline
Misiurewicz: non-recurrent critical orbits $\Rightarrow$ finitely many
transitive pieces &
Misiurewicz diffeomorphisms (paragraph f-); theorem
\ref{thm:finite classes}: finitely many generalized homoclinic classes\\
\end{tabular}
\end{center}


% \paragraph{A dictionary.} The results of this paper transpose to surface
% diffeomorphisms, statement by statement, the classical role of critical
% points in one-dimensional dynamics. The critical point of an endomorphism
% (where the derivative vanishes) corresponds to definition \ref{def:crit},
% formulated through the projective cocycle. Ma\~n\'e's theorem \cite{M},
% absence of critical points implies essential hyperbolicity,  corresponds to
% theorem \ref{p.existence} together with corollary \ref{cor:mane2d}. The
% elementary fact that a periodic critical point must be attracting corresponds
% to the first item of proposition \ref{properties}; the fact that an
% exponentially recurrent critical orbit lies in the basin of an attracting
% periodic point corresponds to proposition \ref{exponential recurrence}. The
% one-dimensional dichotomy for unstable branches of repelling periodic points, either the branch meets a critical point or it lies in a periodic
% interval on which the map is a diffeomorphism, corresponds to theorem
% \ref{disk2}. Finally, the Misiurewicz condition, non-recurrence of the
% critical orbits, which in dimension one yields a decomposition of the
% dynamics into finitely many transitive pieces, corresponds to the
% definition of paragraph f- and theorem \ref{thm:finite classes}.

\paragraph{Organization of the paper.}

The paper is organized as follows. Section \ref{examples} collects examples that calibrate the definition of the critical set against known constructions: tangencies between stable and unstable manifolds, Misiurewicz maps in the H\'enon family, recurrent critical points without exponential growth of the projective cocycle, the constructions of Benedicks--Carleson and Wang--Young, odometers and unicritical sets, and critical points outside the limit set. Section \ref{sec:preliminaries} gathers, for reference, the results on dominated splitting that are used in the sequel; this section is mostly expository and follows \cite{PS1,PS2}. Section \ref{proof-existence} proves Theorem \ref{p.existence}, the characterization of dominated splitting by the absence of critical points under the far-from-homotheties hypothesis. Section \ref{sec:dissipative} proves Theorem \ref{dissipative-FFH}, that under sufficient dissipation every power of $f$ is far from homotheties off the basins of the attracting periodic orbits; this is the longest and most technical part of the paper, organized around a singular-frame factorization of the derivative cocycle, and concludes in subsection \ref{subsec:relaxing} with a discussion of how the sufficient dissipation hypothesis might be weakened. Section \ref{sec:properties} establishes the properties of the critical set listed in Proposition \ref{properties}, including the compatibility between critical directions and the subbundles of a nearby dominated splitting (Proposition \ref{pr:compatibility}). Section \ref{sec:disk} proves Theorem \ref{disk2} on the dichotomy for unstable branches of saddle points for mildly and sufficiently dissipative diffeomorphisms of the disk. Section \ref{sec:misiu} introduces a generalized notion of homoclinic class and proves Theorem \ref{thm:finite classes} on the finiteness of such classes for mildly and sufficiently dissipative Misiurewicz diffeomorphisms. Section \ref{sec:questions} collects some questions left open by this work, concerning the structure of the critical set, the strengthening of the critical-point condition, and the extraction of SRB measures from conditions on the critical orbit.


% Why next theorem is significant:  
% \begin{theorem}\label{misiu2} Let  $\La$ be a compact invariant set such that $f_{|\La}$ is a Misiurewicz map. Then, there exists $g$ $C^r-$arbitrary close to $f$ exhibiting a homoclinic tangency with orbit close to $\La$. In particular, there exists a Newhouse domain ${\cal U}_r$ such that $f\in \bar{\cal U}_r$ and $\bar U\cap \La\neq\emp.$
 
% \end{theorem}


% \paragraph{i-- Collet-Eckman condition.}



% \begin{definition}
%  Let $\La$ be a compact invariant set for a $C^{2}$  surface diffeomorphisms $f$ such that $f_{|\La}$ is dissipative. It is said that $f_{|\La}$ satisfies a  Collet-Eckman condition if there exist positive constants $\delta, \la, \si$ with $\la<1, \si>1$ such that any $c\in Crit(f, \La)$ is a forward $(\delta, \la)-$regular critical point and there is  $w\in T_cM$ such that  $|Df^n(w)|>\si^n$.
% \end{definition}

% \begin{theorem}\label{CE} Let $f\in Diff^2(\DD)$ be a mild dissipative diffeomorphis without semi-attracting periodic points and satisfying the Collet-Eckman condition with constants $(\de, \la, \si)$. If (relation between the constants)  Then $f$ satisfies the almost Misiurewicz condition and therefore the nonwandering is decomposed in a finite number of transitive pieces.
 
% \end{theorem}









% \paragraph{j-- Twist maps,  Green bundles, dominated splitting and critical points.}
% \label{twist-sec}



% We want to relate the notion of critical points  with the tools developed in \cite{Ar},  properly adapted  for dissipative diffeomorphisms. 

% \begin{definition}
    
% Let $f\in Diff^{1}(M)$ and a compact invariant set $\La$ it is said that $f_{|\La}$ is {\it twist} if  there exists a continuous oriented subbundle $x\to v_x$ defined  in a neighborhood of $\La$   and  $\al_0>0$ such that   $slope(G(v_x), v_{f(x)}> \al_0$ and the component of $G(v_x)$ over $v_{f(x)}$ is non-negative. Usually, the subbundle $v$ is called the vertical direction.
% \end{definition}

% % given that we consider line bundles in the tangent, “projective space” of ${\mathbb P}(T_xM)$ is the same thing as the unit tangent circle at $x$ modulo antipodal identification. In that case,

% Recall that the unitarian tangent bundle $(T_xM)_1$ of a surface is a circle;  given  a vector $v$ in $(T_xM)_1$ ,  $(T_xM)_1\setminus\{v, -v\}$ has two connected components that are intervals  and each of them has a natural order compatible with the positive (anti-clockwise) orientation of the unitarian tangent bundle.

% Now, given  $f\in Diff^{1}(M)$ and a compact invariant set $\La$ such that $f_{|\La}$ is twist, one can define the following sequences  $\{G^n(v_{f^{-n}(x)})\}_{n>0}$ and $\{G^{-n}(v_{f^{n}x})\}_{n>0}$. If both sequences of vectors are in the same component of $(T_xM)_1\setminus\{v_x, -v_x\}$ it is said that $f_{|\Lambda}$ {\it do not have conjugate points}. In that case, it follows that $\{G^n(v_{f^{-n}(x)})\}_{n>0}$ is an decreasing sequence, $\{G^{-n}(v_{f^{n}(x)})\}_{n>0}$ is increasing, and $G^n(v_{f^{-n}(x)})> G^{-m}(v_{f^{m}(x)}) $ for any $n, m$. Therefore, in this situation, the following limits exist
% $$G^+_x= \lim_{n\to +\infty} G^n(v_{f^{-n}(x)}),\,\,\,\,\, G^-_x= \lim_{n\to +\infty} G^{-n}(v_{f^{n}(x)}),$$
% and  those functions are called {\it Green's bundles}; moreover,  $$G^-_x\leq G^+_x, \,\, \mbox{for any}\,x\in \Lambda.$$







% The next theorem recasts the definition of dominated splitting using Green's bundles.







% \begin{theorem}\label{green domination}{\bf Green bundles and dominated splitting}. Let $\La$ be a compact
% invariant set  such that $f_{|\La}$ is twist.  Then it holds that $\La$ has a dominated splitting $E\oplus F$ if and only if  $G^-_x$ and $G^+_x$ are transversal directions; moreover, it follows that $E_x=G^-_x $ and $F=G^+_x.$
% \end{theorem}


% Now we relate the Green bundles to the critical direction at the critical point. 



% % \begin{definition} Let  $f\in \operatorname{Diff}^{1}(M)$ and let $\La$ be a  compact invariant set such that $f_{|\La}$ is twist; 

% % \end{definition}

% \begin{theorem}\label{green critical} {\bf Green bundles and critical points}. Let  $f\in \operatorname{Diff}^{1}(M)$ and let $\La$ be a  compact invariant set such that $f_{|\La}$ is twist, dissipative and there are not conjugate points. Then if $\La$ has not a dominated splitting  then the critical points are contained in the sets of points  where $G^+=G^-$; moreover, both directions coincide with the critical direction $v$. 
 
% \end{theorem}
 




% \paragraph{k-- Twist minimal dissipative dynamics in the annulus.}
% In this section, we revisit the critical point using the techniques developed in previous paragraph for the specific case of dissipative annulus dynamics (denoted as $\operatorname{Diff}^{1+\alpha}_{diss}(\AA)$)  with zero entropy. More precisely, we consider diffeomorphisms from $f:\AA\to f(\AA)$ where $\AA$ is an annulus such that $\overline{f(\AA)}\subset {\AA}$ and $|det(D_xf)|<\lambda<1$ for all $x\in \AA;$ let us denote with $\La$ the attractor in the annulus, i.e.,  $\La:=\cap_{n>0}f^n(\AA)$.
% % we also assume that $f_{|\La}$ is twist. 


% % \begin{remark}\label{graph bound} Under the hypothesis of theorem \ref{twist}, it follows that if $\gamma$ is one of the boundaries of $\AA$ then $f^m(\gamma)$ is smooth graph respect to the vertical provided for any positive   $m$.
 
% % \end{remark}




% \begin{theorem}\label{thm:annulus} Let  $f\in \operatorname{Diff}_{diss}^{1+\al}(\AA)$, let us assume that  $f_{|\AA}$ is twist and  $\La$ is  a  minimal  aperiodic  attractor. Then $\La$ is a closed invariant curve that is a Lipschitz graph and  it follows that either

% \begin{itemize}
%     \item[--] $\La$ is  a normally hyperbolic curve with $f_{|\La}$ is semiconjugate to an iirational rotation or 

%     \item[--] there is a tangency; i.e. there is $x\in \La$ such that the graph is differentiable at $x$ and its tangent coincide with a contractive direction. Moreover, one of the iterates of that orbit is a critical point.
 
% \end{itemize} 
% \end{theorem}


% In previous theorem, it follows that if $f$ is $C^2$ and $\La$ is normally hyperbolic, then the curve is also  $C^2$ and $f_{|\La}$ is conjugate to an irrational rotation. 











% % \begin{definition} It is say that $\La$ is a normally hyperbolic curve if there is an invariant splitting $E^s\oplus F$ such that $E^s$ is uniformly contractive and $F=T_x\La.$ Observe that there is $\ga_0$ such that $\angle(E^s, F)\geq \ga_0$.
 
% % \end{definition}

% % \begin{remark} If $\La$ is a normally hyperbolic curve then $G^+=F$ and $G^-=E^s.$
 
% % \end{remark}




% % \begin{theorem} \label{goes to tangency} Let $f_n\to f$ be a sequences of dissipative twist annulus diffeomorphisms  such that $\La_n$ is a minimal  normally hyperbolic curve and $\La$ is a not a normally hyperbolic Lipschitz graph. Then $sl(E^s_n, F)\to 0.$
 
% % \end{theorem}



% %\paragraph{m-- Arnold's family.}

%   A special example of dissipative twist maps in the annulus is the one induced by  the Arnold's family:
% $$\hat f_{a,\omega}(x)= x+a\sin(x)+\omega$$ where $\omega$ is a real number.
% The two-dimensional version is given by  $$f_{a,\omega, b}(x,y)= (\hat f_{a, \omega}(x)+y, b(\hat f_{a, \omega}(x)-x+y))$$
% It has been proved  in \cite{CKP} the following







% \begin{theorem}\label{arnoldthm2} For any $0<b<1$ and $\omega_0$, there exists curve $a\to \omega_{\omega_0}(a)$ such the attractor $\La_{a, \omega_{\omega_0}(a), b}$ is $Lipschitz-$graph such that the rotation number (restricted to the attractor) is $\omega_0.$
 
% \end{theorem}

% Combining previous theorem with theorem \ref{thm:annulus} it follows that  either $\La_{a, \omega_{\omega_0}(a), b}$ is a normally hyperbolic curve or it contains critical point where the graph is differentiable and the contractive direction is tangent to the graph.




\section{Examples}
\label{examples}

In this section we present examples of non-hyperbolic diffeomorphisms for which the critical set can be identified.


\paragraph{{\em i-} Tangencies and critical points.}

Homoclinic and heteroclinic tangencies have played a fundamental role establishing the existence of open sets of non-hyperbolic surface diffeomorphisms. In our context, these systems have critical points.  So, tangencies may provide a natural candidate for critical points.

\begin{definition} We say that $f$
has a {\it heteroclinic tangency} if there is a pair of periodic points $p, q$ such that
there is a point $x\in W^s(p)\cap W^u(q)$ with $T_xW^s(p)=
T_xW^u(q).$ In this case, we say that $x$ is a point of tangency.
\end{definition}


Lemma 1.3.2 in \cite{PRH}  relates the tangencies with the critical points.
Observe that if $x$ is a point of tangency, then any iterate of $x$ is also a point of
tangency. However, the lemma says that in the orbit of tangency there is exactly one   point $x$ and an integer $k$  such that $|g^{-n}_x(v)|\geq 1$ for any $n\geq 0$  and $|g^n_{f^k(x)}(G^k(v))|\geq 1$ also for any $n\geq 0,$ where $v$ spans  $T_xW^s(p)$ with $T_xW^u(q)$. In that sense, a tangency seems to be  a critical point up to an iterate. However, in the following lemma it is shown that given a tangency, there is an iterate that is a critical point in the sense of definition \ref{def:crit}.

\begin{lemma}\label{lem:tang-crit} Let $x$ be a point of  heteroclinic or homoclinic tangency then, there is an iterate of $x$ that it is a critical point.
    
\end{lemma}

\begin{proof} We provide two proofs, one that relies on theorem \ref{p.existence} and another one that follows from an easy claim about finite products of positive numbers (which also highlights the idea beyond the proof of theorem \ref{p.existence}). For simplicity, let us assume that the tangency is associated with a pair of fixed points.

For the first proof, consider the closure of the orbit of the tangency; this set consists of the orbit of the tangency together with the two fixed points where it accumulates; that set does not have dominated splitting and so by theorem \ref{p.existence} it contains a critical point that cannot be one of the fixed points.

The second proof follows from the next claim:

\begin{claim}\label{cl:product} Let $a_1,\dots, a_N$ be a sequence of positive numbers such that $\prod_{i=1}^N a_i\geq 1;$ then there exists $K\leq N $ such that $\prod_{i=0}^{j} a_{K-i}^{-1}\geq 1$ for $j=1,\dots K-1$ and $\prod_{i=1}^{j} a_{K+i}\geq 1 $ for $j=1,\dots, N-K.$
    
\end{claim}

The proof of this claim is simple; consider all cumulative products up to the index $j,$ i.e. $P_j=\prod_{i=1}^j a_i$; take  $K$ to be an index at which the cumulative product is minimal..

To apply the claim, observe that one can take a point $x'$ in the orbit of the tangency, a vector $v'$ (the common tangent direction of the stable and unstable manifolds at point $x'$) and a positive integer $N$ such that $|g^{n}(v')|\geq 1$ for any $n<0$ or $n>N$ and $g^n(G^k(v'))\geq 1$ for any $n>0$. Observing that $|g^N(v)|=\prod_{i=0}^{N-1}|g(G^i(v))|\geq 1$, by the claim one gets a positive integer $K$ such that $ g^{-j}(G^K(v'))=\prod_{i=0}^{j} g^{-1}(G^{K-i}(v))\geq 1$ for $j=1,\dots K-1$ and $g^j(G^K(v'))=\prod_{i=1}^{j} g(G^{K+i}(v'))\geq 1 $ for $j=1,\dots N-K;$ therefore, $g^{|n|}(G^k(v'))\geq 1$ for any $n.$ 
    
\end{proof}

Observe that in case of tangencies, the notion of critical point is stronger; in fact., if $x$ is a critical point associated with an orbit of tangencies, then for its critical direction $v$ it follows that there exist $\la>1$ and a positive integer $N_0$ such that for any  $n>N_0$ and $k>0$ holds $|g^k(G^n(v))|>\la^k$ and $|g^{-k}(G^{-n}(v))|>\la^k.$ In another interpretation, one can consider the derivative along the orbit of the tangency; this yields a sequence of matrices (indexed by $\ZZ$) $(H_n)$ such that: for any $n\neq 0$ the matrix $H_n$ is diagonal with the horizontal being contracted and the vertical being expanded for any $n<0$, and conversely,  with the horizontal being expanded  and the vertical being contracted for any $n>0$, and $H_0$ being a matrix that sends the vertical into the horizontal. 

% More precisely,  it was proved  that given  a point of tangency, then there is a  unique iterate of that point such that it  is a critical value.

The notion of tangencies can be extended beyond periodic points; in fact, one can consider tangencies between the stable and unstable manifold of points in a hyperbolic set (as was considered in the paper by Newhouse; see \cite{N}). As noted in the introduction, Kupka–Smale diffeomorphisms, those for which all periodic points are hyperbolic and all intersections between stable and unstable manifolds of periodic points are transversal, form a residual set and therefore exhibit no heteroclinic tangencies. Nevertheless, by Newhouse's results there is an open set of surface diffeomorphisms that are not hyperbolic and do not admit a dominated splitting; the diffeomorphisms in this open set therefore contain critical points. One can ask, though, whether if they could be related to heteroclinic tangencies associated with non-periodic points in a hyperbolic set (see paragraphs iii- and  iv- in the present section).

% therawe will see in other examples, the existence of critical points does not imply that there is a tangency associated with periodic points.




\paragraph{{\em ii-}  Misiurewicz diffeomorphisms.}  In \cite{R} the author introduces an example of a homoclinic class that is semiconjugate to a hyperbolic set exhibiting a quadratic homoclinic tangency taking place in the limit set.  In \cite{E}, a surface diffeomorphism on the boundary of the Anosov class (and topologically conjugate to an Anosov diffeomorphism) is constructed, exhibiting a cubic tangency associated with a heteroclinic point.
  Both are  examples  of Misiurewicz diffeomorphisms.


Another example of a Misiurewicz diffeomorphism can be found in the H\'enon family $ f_{\mu,b}(x,y)=(x^2+\mu+by,x)$.  For sufficiently large negative $\mu$, this map has a hyperbolic Cantor set. Increasing $\mu$, one can find a first bifurcation parameter $\mu_0$ at which the hyperbolicity of the Cantor set is destroyed due to an internal tangency of the stable and unstable manifolds of a saddle point (see \cite{T} and \cite{BS2}).
Moreover, for the H\'enon family with $|b|<1$, the parameter set on which the topological entropy attains its maximal value $\log 2$ has the following structure (see Theorem 2 of \cite{BS}): in its interior, the set of non-escaping points to infinity forms a hyperbolic horseshoe; on its boundary, the diffeomorphism exhibits an internal tangency.

% Another example is provided by a transitive diffeomorphism that is conjugate to a linear Anosov and exhibiting a homoclinic tangency of a fix point; more precisely, this example  is  obtained by isotopy from a linear Anosov and they are conjugate to a linear Anosov.  
% This last example can
% be tweaked to another example exhibiting infinitely many topological transversal tangencies associated with periodic points. In the similar way, the example in \cite{E} can be also adapted to show a cantor set of (cubic) heteroclinics tangencies between compact pieces of  the stable and unstable  manifolds of two disjoint hyperbolic sets.

All the examples above lie on the boundary of the Axiom A class. To show that there is exactly one critical point in each example, the one associated with the tangency, it suffices to show that at every point outside the orbit of the critical point there is a well-defined cone field contracted by some iterate of $f$.





\paragraph{{\em iii-} Recurrent critical points and no exponential growth of the projective cocycle.} In \cite{HMS} it is shown that there exist diffeomorphisms in the boundary of a Axiom A surface diffeomorphisms such that the lack of hyperbolicity arises from the presence of a cubic tangency between a stable manifold and an unstable manifold, one of which is not associated with a periodic point; moreover, that tangency provides a critical point which is recurrent and so this critical point is not a heteroclinic tangency associated with points in a hyperbolic set.  


The following example shows that given a horseshoe and a transverse homoclinic point (not necessarily associated with a periodic point), the diffeomorphism can be deformed so that the transverse intersection becomes a tangency, thereby producing tangencies, and hence critical points, at recurrent points. Moreover, the construction can be carried out so that the critical point does not exhibit exponential growth of the projective cocycle.

\begin{proposition} \label{recurrent critical}

There exists a $C^1$-diffeomorphism  with an invariant compact set
conjugate to the shift on $\{0,1\}^\ZZ$ and with a single critical orbit.
The critical orbit can be taken  at any non-periodic point of $\{0,1\}^\ZZ$.

Moreover, the construction can be performed in such a way that the derivative of the projective cocyle does not grow exponentially. 
\end{proposition}

\begin{proof}
    Let us consider a dissipative diffeomorphism $f_0$ exhibiting a  classical horseshoe $K_0$ 
which is the maximal invariant set in an open domain $U$
and let us fix any non-periodic point $x\in K_0$.
We compose $f_0$ with a fibered rotation centered at $x$ with arbitrarily small
support and which almost sends $E^u_x$ to $E^s_x$.
We obtain a diffeomorphism $f_1$.
If the support $B_0$ of the rotation is small enough, a (narrow) cone field is still preserved,
the dynamics is still hyperbolic and conjugated to the shift,
the conjugacy is $C^0$-close to the identity,
the bundles $E^u,E^s$ are almost unchanged outside
the union of the $N_0$-iterates of the support $B_0$ of the rotation (where $N_0$ is large,
independent of the size of the support).
The new horseshoe $K_1$ is still the maximal invariant set in $U$.

We repeat the construction and get a sequence of diffeomorphisms $f_n$,
a sequence of horseshoes $K_n$, and a sequence of perturbation balls $B_n$.
%towers $T_nf_n^{-N_n}(B_n)\dots, B_n,\dots, f_n^{N_n}(B_n)$.
One can suppose that the diameter of $B_n$ and the angle of the rotation that we perform at step $n$
are less than $2^{-n}$ (after step $n$ the angle between $E^s$ and $E^u$
at $x$ is less than $2^{-n-1}$, so the rotation at step $n+1$ by an angle
smaller than $2^{-n-1}$ can still close the cone).
The convergence is thus strong enough, $f_n$ converges $C^1$ to
a diffeomorphism $f$ and $K_n$ to an invariant compact set $K$,
still the maximal invariant set in $U$,
and topologically conjugate to the shift.
Of course we still have dissipation.

The diffeomorphism $f_0$ preserved two transverse cone fields.
Outside these cones $g$ contracts (in the future or in the past depending on the cone).
Outside $B_0$ the cones are disjoint, so any point outside $B_0\cup f_0(B_0)$
does not satisfy the definition of critical set, for any diffeomorphisms $f_n$
(they all coincide there).
We want to repeat this argument so that after step $n$,
the critical set is contained in $B_n$. At the limit the critical set is reduced
to $x$ (which can be checked to be critical).
After step $n$, we consider two transverse invariant cone fields,  $C^s$ and $C^u$,
that are very thin, so that for $m_n$ large enough,
$g^{m_n}$ is expanding only inside $C^s$ and $g^{-m_n}$ is expanding only inside $C^u$.
Moreover after perturbation, we suppose that $f_{n+1}$ still sends $C^u$
disjoint from $C^s$, so that no point outside $B_{n+1}$ can be critical after step $n+1$.
\end{proof}



\paragraph{{\em iv-} Critical points for non-uniform hyperbolic diffeomorphisms  in the  works of Benedicks-Carlesson and Wang-Young.}

The idea of identifying a critical set for two-dimensional maps, a set designed to play the role analogous to that of critical points in one dimension, goes back to \cite{BC}. In fact, this approach is implemented in the landmark paper of Benedicks-Carleson \cite{BC} to prove the existence of strange attractors for a positive measure set of parameters for the H\'enon family. In their setting the critical set is defined as a certain type of tangencies between the unstable manifold of a fix point and stable ones that are obtained after an inductive process and for certain parameters of the  H\'enon family (in particular, that tangency is  not a heteroclinic tangency associated with points in a hyperbolic set).
In \cite{WY}, the inductive construction of \cite{BC} is reworked, with built-in geometric properties for the critical set incorporated as part of the induction. 
The critical set introduced in \cite{WY} is intrinsically defined, characterized as the set of points on the attractor at which stable and unstable directions are interchanged. It is shown to have a special geometric structure: it is the intersection of a nested sequence of rectangles whose geometric and dynamical properties are explicit. For a precise description of the geometric properties of the critical set in 
this setting and its relation to the interchange of stable and unstable 
directions, we refer the reader to section~2 of \cite{WY}. 
Recently, the critical set has been used in \cite{BoS} to define a kneading theory for some H\'enon maps.

% In a more general context, one can consider non-uniform hyperbolic measures; for instance, one may consider an homoclinic class that remains non-hyperbolic by small perturbation (again, that exists  due to \cite{N}) and by the result in \cite{BCS} the exhibits measure of maximal entropy that cannot be hyperbolic. The natural question that follows is to try to identify the critical points for these non-hyperbolic homoclinic classes.


\paragraph{{\em v-} Critical points in Odometers and unicritical sets.}

In \cite{CLPY}, the dynamics of odometers and their renormalization is studied, introducing a class of infinitely renormalizable unicritical diffeomorphisms of the disk (with a non-degenerate ``critical point”). Odometers appear naturally for mildly dissipative diffeomorphisms of the disk. Since they do not admit a dominated splitting (see \cite{CPT}), they contain critical points. In \cite{CLPY} a it is assumed a that there is a {\it unique critical point} in the odometer and  the definition is strengthen by requiring the projective cocycle to grow exponentially along the critical direction in both time directions, i.e. $|g^n(v)| > (1+\de)^{|n|}$ for all $n$ and some $\de>0$. Under these conditions, together with unique ergodicity, one can show the existence of local stable and center-stable manifolds that are tangent at the critical point; this is used to prove that under renormalization, the dynamics converges to a unimodal one-dimensional endomorphism.


\paragraph{{\em vi-} Critical points outside the Limit set.} The classical Smale's horseshoes dynamic in the sphere provides a simple example of critical points that  are not in the limit set: the dynamics decomposes into a repelling fixed point and a trapping disk; inside the disk, the limit set is the union of a hyperbolic Cantor set and a fixed attracting point. The critical point is the ``fold'' of the unstable manifold of the hyperbolic set, lying in the basin of attraction of the attracting fixed point. This is in fact the typical situation for non-trivial dissipative Axiom A diffeomorphisms on a trapping disk. 

Another easy example arises by considering the unstable branch of a fixed hyperbolic saddle $p$ that is attracted by an attracting fixed point $q$: if the attracting fixed point has two different eigenvalues, it holds that associated with the smallest one there is a unique invariant strong-stable foliation in the local basin of attraction $q$ and observe that $|g(v)|>1$ for any vector tangent to that foliation. If the unstable manifold of $p$ is at some point tangent to this strong-stable foliation, then along the orbit of that point there is a critical point (for details, see paragraph i- in the present section and claim \ref{cl:product}). Similarly, if both eigenvalues of $q$ are equal, then for any vector $v$ in the tangent bundle of a point in the local basin of $q$ satisfies $|g^n(v)|\leq 1.$




% \paragraph{{\em vii-} Higher dimension.}
% A proposition: a point $x$ is critical if there exists two non-transverse subspaces
% $E,F$ such that $\dim(E)+\dim(F)=\dim(M)$ and
% in the space of grassmannian of dimension $\dim(E)$,
% $g^n(E)\geq 1$ and in the space of grassmannian of dimension $\dim(F)$,
% $g^{-n}(F)\geq 1$, for all $n\geq 1$.\\


% \begin{quest}
% Do we have a similar result in higher dimension?
% \end{quest}



\section{Preliminaries about dominated splitting}
\label{sec:preliminaries}
This notion was  introduced independently
by Ma\~n\'e, Liao, and Pliss,as a first step toward proving the stability conjecture: that structurally stable systems are hyperbolic. Simple examples of invariant sets of surface diffeomorphisms that exhibit a dominated splitting but are not hyperbolic include: normally hyperbolic periodic arcs; normally hyperbolic closed invariant curves on which the dynamics is conjugate to an irrational rotation; and homoclinic classes associated with non-hyperbolic periodic points. In \cite{PS2} a general description of the limit set was provided under the assumption of dominated splitting.


   
If $f_{|\Lambda}$ is dissipative and admits a dominated splitting $E \oplus F$ then $E$ is uniformly contracted by $Df$. Since the present paper deals with dissipative diffeomorphisms, we will write dominated splittings in the form $E^s \oplus F$, with $E^s$ uniformly contracted.

For a set $\La$ with a dominated splitting $E^s\oplus F$, there are two complementary $C^1$ local laminations: one is  the  local stable  lamination $\{W^{s}_\eps(x)\}_{x\in \La}$ satisfying $f(W^s_\e(x))\subset W^s_{\la.\e}(f(x))$ for some positive $\la$ smaller than one, and the other is the local center-unstable $\{W^{cu}_\eps(x)\}_{x\in \La}$ tangent to $F$. On the other hand, for any point $x$ and $\eps>0$ one can define   the local unstable set of size $\eps$, denoted as   $W^{u}_\eps(x)$, as the set of points $y$ such that $dist(f^{-n}(x), f^{-n}(y))<\eps$ for any positive integer $n$ and  $dist(f^{-n}(x), f^{-n}(y))\to 0$ when  $ n\to +\infty$.  It holds that in a set with dominated splitting, the local unstable is contained in the  center-unstable. 

The following lemma has been proved in \cite{PS1}.

\begin{lemma}\label{bounded-arcs} Let $\La$ be a set with dominated splitting $E\oplus F$ of a $C^2$ surface diffeomorphisms. There exist $\e>0$  and $\de>0$ such that for any  arc $J$ transversal to $E$ such that the forward iterates of $J$ remains in an $\e-$neighborhood of $\La$ and $length(f^n(J))<\de$ for all $n>0$ it holds that there is a normally hyperbolic periodic arc $I$ such that $J\subset W^s(I).$
    
\end{lemma}




% For a  dissipative invariant  set $\La$ with dominated splitting,  it follows that one of the subbundles of the dominated decomposition is a uniformly contracted one; therefore, there exists $\eps>0$ such that  any local stable set of size $\eps$ is a submanifold that all together form a $C^1-$lamination. 

The previous lemma has a stronger version when the set $\La$ is a continuous set. 
% Now we provide a version of lemma \ref{centerunstable} in a particular context. 

 
 \begin{lemma}\label{centerunstable2} Let $f\in Diff^2(M)$ and let $\La$ be a connected   compact invariant  with a dominated splitting $T_\La M=E^s\oplus F$. Let $U$ be a neighborhood of $\La$ where the dominated splitting can be extended. Let $J$ be a $C^1-$compact arc transversal to $E^s$ such that the forward iterates of $J$ are contained in $U.$ Then

 \begin{itemize}
     \item[--] either $length(f^n(J))\to \infty$ as $n\to +\infty,$

     \item[--] or the   length of a subsequence of  forward iterates of $J$ is uniformly bounded  above and   

     \begin{itemize}

       \item[--] either the forward iterates of $J$ converges to a normally hyperbolic periodic arc or 
         \item[--] $length(f^n(J))\to 0$ and  $J$ lies in the basin of attraction of a (semi)attracting periodic point. 
     \end{itemize}

     
 \end{itemize}


 
     
 \end{lemma}

% A similar statement can be formulated for arcs transversal to $F$ direction that its backward iterates remains in a neighborhood of $\La$ with bounded length concluding in that case that its backward iterates converges also to a normally hyperbolic arc.

The local center-unstable manifold $W^{cu}_\e(x)$, minus the point $x$, has two components that are called local center-unstable branches and denoted as $\Ga^{cu}_\e(x)$. As a consequence of lemma \ref{bounded-arcs} one can obtain the following dynamical characterization of the center  unstable  manifolds.

\begin{lemma}\label{centerunstable}
For any center-unstable branch $\Ga^{cu}_\e(x)$ of a point $x$ one has the following characterization:






\begin{itemize}

    \item[--] either $x$ is a (semi)attracting periodic point and there is $\e(x)$ such that $\Ga^{cu}_{\e(x)}$ is contained in the basin of attraction of that point, 


    \item[--] or there exists $\e(x)$ and $N$ such that for any $n> N$ it holds that 
 
    
    $f^{-n}(\Ga^{cu}_{\e(x)}(x))\subset \Ga^{cu}_\e(f^{-n}(x))$ and either
    \begin{itemize}
        \item $\Ga^{cu}_{\e(x)}(x)\subset W^u_\e(x)$ or

        \item $x$ is a periodic point, $\Ga^{cu}_{\e(x)}(x)$ is normally hyperbolic  and there is a sequence of periodic points of the same period of $x$ contained in $\Ga^{cu}_{\e(x)}(x)$ and accumulating on $x.$
        
    \end{itemize}
    

    
\end{itemize}

\end{lemma}

Previous lemma has a stronger version summarized in following remark.

\begin{remark} For recurrent non-periodic points, a stronger statement holds: both local unstable branches are ``large", or at least one is large and the other extends to an intersection with the stable manifold of a semi-attracting periodic point.

 More precisely, there is $\e>0$ such that for any recurrent non-periodic point one of the following holds
 
 \begin{itemize}
     \item[--] $W^{cu}_{\e}(x)=W^{u}_{\e}(x)$

     \item[--] one center-unstable branch, called $\Ga^{cu,+}_{\e}(x)$, coincides with $W^{u,+}_{\e}(x)$ and in the other branch $\Ga^{cu,-}_{\e}(x)$ contains a point $x'$ in $ W^s(q)$ for some   semi-attracting periodic point $q$, with $\Ga^{cu,-}_{[x, x']}(x)= W^{u,-}_{\e}(x)$, where $\Ga^{cu,-}_{[x, x']}(x)$ denotes the arc of center-unstable branch with extreme points $x$ and $x'$.   
 \end{itemize} 
    
\end{remark}




% Similar statement can be proved for center stable branches, replacing (semi)attracting by (semi)repelling.

As a consequence of previous statement, one can get the following proposition that resembles the local product structure for hyperbolic sets but in the dominated splitting  case, for recurrent non-periodic points.






\begin{proposition}\label{prop:LPE} Let $\La$ be a compact invariant  set with dominated splitting  of a $C^2$ surface diffeomorphism. There exists a positive constant $\de$ such that for  any recurrent points   $x$ and $y$ at distance smaller than  $\de$, it holds that $W^{\si}_\eps(x)\cap W^{\tau}_\eps(y)\neq \emptyset$ for both $(\si,\tau)=(u,s)$ and $(\si, \tau)=(s,u)$.

\end{proposition}






The following proposition follows from theorem A and corollary 4.1.1 in \cite{PS2}.

\begin{proposition}\label{PS2-1} Let $\La$ be a compact invariant set with dominated splitting  of a $C^2$ surface diffeomorphism. There exists a positive integer $N_0$ and a positive constant $\de$ such that any pair of periodic points $p, q\in \La$ is homoclinically related provided that their periods are larger than $N_0$ and their distance is smaller than  $\de$.
    
\end{proposition}





Now, we state a few simple results (used later in section \ref{sec:misiu}) about sets with dominated splitting contained in the limit set without assuming that the whole limit set has dominated splitting.

Let $\La$ be a set with dominated splitting; it is said that $z\in \La$ is an {\it interior point} if there exist recurrent points $z', z'' \in \La$ arbitrarily close to each other such that 

\begin{itemize}
    \item[--] $W^{s}_\eps(z')\cap W^u_\eps(z'')\neq \emptyset $ and $W^{u}_\eps(z')\cap W^s_\eps(z'')\neq \emptyset;$ 
    %belong to different connected components of $W^{s}_\eps(z)\setminus \{z\};$

    \item[--] $z$ is in the interior of the rectangle bounded by the local stable and unstable arcs of $z'$ and $z''.$ 
    %$W^{u}_\eps(z)\cap W^s_\eps(z') $ and $W^{u}_\eps(z)\cap W^s_\eps(z'') $ belong to different connected components of $W^{u}_\eps(z)\setminus \{z\}.$
\end{itemize}

 
    
It is said that a set $\La$ with dominated splitting is a {\it heterocycle} if it contains a finite number of  saddle periodic orbits and any other point in $\La$ is a transversal intersection point between the stable and unstable manifolds of some  periodic orbits in $\La.$  

The following remark  is about characterizing heterocycles and is used in section \ref{sec:misiu}.

\begin{remark}\label{rmk:heterocycle}
Observe that a heterocycle provides a multi-rectangle $S$ bounded by  the stable and unstable arcs of finite number of periodic points and the points in the heterocycle are either those periodic points or are contained in   those arcs. By the classical $\la-$lemma it follows that one can get an arbitrarily small rectangle $R$ containing a periodic point of the heterocycle in one corner and bounded by two parallel stable arcs and two parallel unstable arcs of that point. Moreover, if $\omega(x)$ is a heterocycle, any sufficiently large forward iterates of $x$ are in the interior of the multi-rectangle $S$ and there are forward iterates of $x$ in the interior  of $R.$
\end{remark}


% \begin{proposition} Let $f\in Diff^2(M)$ let $x\in M$ such that $\omega(x)$ has a dominated splitting. Then either

% \begin{itemize}
%     \item[--] $\omega(x)$ is a periodic orbit and if $x\in L(f)$, the periodic orbit a saddle and one of its unstable branch intersects $L(f)$ (in particular, the periodic orbit cannot be contained in the interior of a normally hyperbolic arc);

%     \item[--] $\omega(x)$ is contained in a homoclinic class and there exists $\La$ a locally maximal invariant set with dominated splitting arbitrary close to $\omega(x).$
    
% \end{itemize}
    
% \end{proposition}

The next proposition establishes a dichotomy about the possible limit sets with dominated splitting.


\begin{proposition}\label{dicho-omega} Let $f\in Diff^2(M)$ and let $x\in M$ such that $\omega(x)$ ($\alpha(x)$) has a dominated splitting. Then either

\begin{itemize}
  

    
    \item[--] $\omega(x)$ ($\alpha(x)$)  is a periodic orbit,  moreover, either  $x$ coincides with   the periodic orbit or it belongs to the stable set of that periodic orbit;
    % moreover, if $x\in L(f)$, and it is not the periodic orbit, then  the periodic orbit a saddle periodic orbit and one of its unstable (stable) branch intersects $L(f)$ (in particular, the periodic orbit cannot be contained in the interior of a normally hyperbolic arc);

    \item[--] $\omega(x)$ ($\alpha(x)$)  contains interior points;
    
    \item[--] $\omega(x)$ ($\alpha(x)$) is a heterocycle.
    
\end{itemize}
    
\end{proposition}


% \begin{proof} It follows that either $\omega(x)$ is finite or not. If it is finite, then $\omega(x)$ is a periodic orbit  and $x$ either is the periodic orbit or it belongs to the stable set of that periodic orbit. 
% % If $x$ is not the periodic orbit and belongs to the limit set, if follows that the periodic orbits cannot be an attracting periodic point, so it is a saddle and if both unstable branches are not in $L(f)$ then any point in a small  neighborhood of the periodic orbit is wandering and so $x,$ a contradiction.

% If  $\omega(x)$ is infinite,  let us consider the forward orbit of any point $z\in \omega(x)$ that is recurrent. If the forward orbit is finite, then $z$ is a periodic point. If it is infinite, one can  consider three  large forward iterates of $z$ that are  arbitrarily  close to each other. Using  proposition \ref{prop:LPE}  either one can find one iterate {\it that lies between the other two} as it is described in the definition of being an interior point, or  two of them share the same (local) stable manifold, or two of them share the same (local) unstable manifold. 
% In the first case, one gets an interior point. In the second case, one concludes that $z$ belongs to the stable manifold of a periodic point. In the last case, since they belong to the local unstable manifold of a periodic point, but they are themselves arbitrarily large forward iterates,  $z$ must be that periodic point.
% Therefore, one concludes that any $z\in \omega(x)$ is either a periodic point or belongs to the stable manifold of a periodic orbit; similarly, analyzing the backward orbits of $z$ one concludes that any $z\in \omega(x)$ is either a periodic point or belongs to the unstable manifold of a periodic orbit. If there are infinitely many periodic orbits, then at least one is an interior point.

% In the remaining case, it follows that there are finitely many periodic orbits in $\omega(x)$ (that are saddles) and any other point belongs to the stable and unstable of those periodic points. Since $\omega(x)$ has dominated splitting, the intersection points are transversal and so one gets a heterocycle.



    
% \end{proof}


\begin{proof}
Either $\omega(x)$ is finite or infinite. If it is finite, then $\omega(x)$
is a periodic orbit and $x$ either coincides with that periodic orbit or
belongs to its stable set.

Assume now that $\omega(x)$ is infinite, and let $z\in\omega(x)$ be a
\emph{recurrent} point; such a point always exists, since $\omega(x)$ is
compact and invariant and therefore contains a minimal set, all of whose
points are recurrent. Observe that every forward iterate of $z$ is again a
recurrent point, so that proposition \ref{prop:LPE} applies to any pair of
forward iterates of $z$.

If the forward orbit of $z$ is finite, then $z$ is a periodic point. If it
is infinite, one can consider three large forward iterates of $z$ that are
arbitrarily close to each other. Using proposition \ref{prop:LPE}, either
one of the three iterates \emph{lies between the other two} as described in
the definition of interior point, or two of them share the same local stable
manifold, or two of them share the same local unstable manifold. In the
first case one gets an interior point. In the second case, $z$ belongs to
the stable manifold of a periodic orbit; since $z$ is recurrent, $z$ belongs
to $\omega(z)$, which is that periodic orbit, and so $z$ is periodic. In the
last case, since the two iterates belong to the local unstable manifold of a
periodic point but are themselves arbitrarily large forward iterates of $z$,
again $z$ must be that periodic point.

Therefore, either $\omega(x)$ contains an interior point, or every recurrent
point of $\omega(x)$ is a periodic point. Let us assume from now on that the
second option holds. In that case there are only finitely many periodic
orbits in $\omega(x)$: otherwise, by proposition \ref{PS2-1}, two of them
with sufficiently large periods and sufficiently close to each other would
be homoclinically related, producing transversal homoclinic intersections in
$\omega(x)$ and hence interior points.

Given now an arbitrary point $y\in\omega(x)$, the recurrent points of
$\omega(y)\subset \omega(x)$ are periodic and belong to the finite family of
periodic orbits above; hence $\omega(y)$, having finitely many periodic
orbits as its only recurrent points, forces $y$ to be either one of those
periodic orbits or a point of the stable manifold of one of them. Applying
the symmetric argument to the backward orbit of $y$ (the points of
$\alpha(y)$ and proposition \ref{prop:LPE} for $f^{-1}$), one concludes that
$y$ is either periodic or belongs to the unstable manifold of one of the
periodic orbits of the family.

In conclusion, $\omega(x)$ consists of finitely many periodic orbits (which
are saddles) together with points belonging simultaneously to the stable
manifold of one of them and to the unstable manifold of another (possibly
the same). Since $\omega(x)$ has a dominated splitting, these intersections
are transversal, and one gets a heterocycle.
\end{proof}

To finish, we discuss the dynamical decomposition of the limit set provided by Theorem A in \cite{PS2}. It states  that if the limit set of a $C^2$ diffeomorphisms of a compact surfaces has dominated splitting, then it can be decomposed into a finite number of isolated periodic points,  finite number of normally hyperbolic arcs, a finite number of normally hyperbolic invariant closed curves with dynamics conjugate to irrational rotation and a finite number of homoclinic classes. 
 
\begin{remark}\label{decomposition}
 
 The same result stated in Theorem A in \cite{PS2} can be obtained if one restricts the limit set to an attracting open domain. More precisely, if $f$ is a $C^2$ surface diffeomorphism and $B$ is an open set such that $\overline{f(B)}\subset B$ and the limit set of $f$ restricted to $B$ has dominated splitting, then the same thesis holds. Moreover, if $B$ is a disk and $f_{|B}$ is dissipative, then there are no normally hyperbolic closed periodic arcs.
    
\end{remark}




% \begin{theorem} Any  compact invariant set with dominated splitting of smooth surface diffeomorphisms either 
% \begin{itemize}
%     \item[--]  is a normally hyperbolic closed curve with dynamic conjugate to an irrational rotation;
%     \item[--] it is conjugate to a hyperbolic set  contained in an homoclinic class and such that any periodic point of a large enough period is a hyperbolic periodic point.
% \end{itemize}
    
% \end{theorem}



% \begin{theorem} Any locally maximal compact invariant set with dominated splitting of a smooth surface diffeomorphisms can be decomposed in a finite number of 

% \begin{itemize}
%     \item[--]   normally hyperbolic closed curve with dynamic conjugate to an irrational rotation;
%     \item[--] normally hyperbolic closed arc;
%     \item[--] sets  conjugate to a hyperbolic set, it is contained in an homoclinic class and any periodic point of a large enough period is a hyperbolic periodic point.
% \end{itemize}
    
% \end{theorem}



% \begin{theorem} Any compact invariant set of type 3 is contained in a locally maximal invariant set.

% \end{theorem}

% \begin{claim} Let $x\in L(f)$, if $\omega(x)$ has dominated splitting then it follows that $\omega(x)$ is contained in a homoclinic class.


% \end{claim}


% \begin{claim} Given a compact invariant set with dominated splitting of a smooth surface diffeomorphisms, there is a locally maximal invariant  set containing it and arbitrary close to it.
    
% \end{claim}
% Properties of locally maximal invariant   sets with dominated splitting:

% \begin{itemize}
%     \item[--] Boundary points; boundary points are in stable/unstable manifolds of periodic points (corner points?)
%     \item[--] Markov partition given by rectangles.
%     \item[--] Interior point.
%     \item[--] Quadrants.
    
% \end{itemize}



% \noi{Proof of theorem \ref{misiu1}:} From \cite{PS1} follows that $\omega(c)$ is a dominated set: either contained in a homoclinic class, or  a normally hyperbolic set; so $c$ belong   to the stable manifold of a point in $\La_\omega$. Arguing in the same way, follows that $c$ belong to the unstable manifold of some point in a hyperbolic set $\La_\alpha.$ Observe that if $E=T_xW^s(\La_\omega)$ and $F=T_cW^u(\La_\al)$ then $g^n(E)>1$ and $g^{-n}(F)>1$ for any $n$ sufficiently large so $E=F=v$, where $v$ is the critical direction.
 

% \qed



% \noi{Proof of theorem \ref{misiu2}:}
%   For each critical point $c$ let $\La_{\om,c}$ be a locally maximal invariant hyperbolic set containing $\omega(c)$ and $\La_{\al,c}$ be a locally maximal invariant hyperbolic set containing $\al(c)$. Since, there exists a neighborhood $V$ of $Crit(f,\La)$ such that $\La_h=\cap_{n\in \ZZ}f^n(V^c)$ is a locally maximal hyperbolic set it follows that  its admits a spectral decomposition in a finite number of localy maximal invariant transitive pieces $\La_1\dots \La_n$. 

% \begin{claim}
% For any $\La_i$ it follows that $W^u(\La_i)\cap Crit(f, \La)\neq \emp.$
% \end{claim}
% Observe   that  $[W^u(\La_i)\setminus \La_i]\cap\La\neq \emp$ otherwise $\La$ is not transitive. Let us consider the sets $\La_i$ such that its unstable manifold do not accumulate on the critical points and let us denote them as $\La_{i'}$.  Reducing the neighborhood $V$ if needed, follows that the parts of the unstable manifold of $\La_{i'}$ inside $\La$  is also inside $\La_h$ (i.e. $[W^u(\La_{i'})\setminus \La_{i'}]\cap\La\subset \La_h$), therefore  $\cup_{i'}W^u(\La_{i'})\cap \La$ is hyperbolic and contained in $\La_h$ so, there is a finite number of hyperbolic set $\La_{i'_1}\dots \La_{i'_k}$ such that $\cup_{i'}W^u(\La_{i'})\cap \La\subset\cup\La_{i'_j}$ but in particular, this implies that the unstable of $\La_{i'_j}$ is contained in the set itself; a contradiction. This finish the proof of the claim.

% From the fact that $\La_h$  is decomposed in a finite number of hyperbolic sets, there exist a finite number of critical points $c_1\dots c_n$ such that for any $c$ follows that $\La_{\al,c}\subset \La_{\al,c_i}\subset \La_{i'}$  and $\La_{\om,c}\subset \La_{\om,c_j}\subset \La_{j'}$ for some $i,i', j,j'$; let us denote $\Ga_{i}=\La_{\al,c_i}\cup \La_{\om,c_i}.$ From the claim, it follows that the unstable of any  $\La_{\om,c_j}$ accumulate on the critical set, in particular, its accumulate on the stable manifold of some $\La_{\al,c_i}$; therefore, there exist a sequence $i_1,i_2\dots i_k$ forming a cycle in the sense that  $W^u(\La_{\om, c_{i_j}})$ accumulates on the stable manifold of $\La_{\al, c_{i_{j+1}}}$ and the unstable of $\La_{\om, c_{i_k}}$ on the stable of $\La_{\al, c_{i_1}}$. From the fact that each of that pieces are localy maximal, it follows that their stable and unstable boundaries are given by the stable and unstable 
% manifold of a finite number of periodic points; therefore, if the unstable manifold of one piece accumulate on the stable one of another piece, it follows that the unstable manifold of a periodic point accumulate on the stable of another periodic point.

% By Pixton  the unstable of $\La_{\om, c_{i_k}}$ intersects transversely the stable of $\La_{\al, c_{i_1}}$ without unfolding the connection between stable and unstable associated with the critical points.

% Now, we can $C^r-$unfold the tangency associated with the hyperbolic set and we get the homoclinic tangency.


% \qed



% \noi{\em Proof of theorem \ref{misiu3}:} The goal is to prove that there is a finite number of basic sets such that any critical point belongs to the stable and unstable manifold of a basic set. For that, it is built boxes around critical point using stable and unstables manifold of the hyperbolic set; it is show that these classes are finite and the return to the critical region is hyperbolic compose with parabolic.

% \qed


\section{Proof of theorem \ref{p.existence}}
\label{proof-existence}


The proof of theorem \ref{p.existence}  uses the  following two lemmas. The first (Lemma \ref{l.H}) provides, under the far-from-homotheties hypothesis, points with bounded backward (resp. forward) projective growth in some direction. The second (Lemma \ref{DScond}) characterizes the existence of a dominated splitting in terms of a uniform decay condition on the projective cocycle. We then deduce Theorem \ref{p.existence} by showing that failure of the dominated splitting condition produces a critical point.

\begin{lemma}\label{l.H}  Let $f\in Diff^1(M)$ and let $\Lambda$ be a compact invariant set such that $f_{|\La}$ is far from homotheties, then the following sets are non-empty:
$$H^+(f,\Lambda)=\{x\in\Lambda, \; \exists \,\,v\in T_x^1M,\;
\forall n\in \NN, \; |g^{n}(v)|\geq 1\},$$
$$H^-(f,\Lambda)=\{x\in\Lambda, \; \exists\,\, v\in T_x^1M,\;
\forall n\in \NN, \; |g^{-n}(v)|\geq 1\}.$$
\end{lemma}
% \begin{proof} It is a direct application of Oseledets's theorem: by the far from homotheties hypothesis any ergodic measure has two distinct Lyapunov exponents.  
%  $H^+(f,\Lambda)$ is empty.
%Any unitary tangent vector $v$ of $\Lambda$ has a positive iterate $n$ such that
%$|g^n(v)|<1$. This implies by compactness that for some $N\geq 1$ one has
%$$\forall v\in T^1M,\; |g^N(v)|<1.$$
%This contradicts that the fibers $T^1M$ have constant length.
% \end{proof}
\begin{proof}
Let us assume by contradiction that $H^+(f,\Lambda)$ is empty. Then every
unit vector $v\in T^1_\Lambda M$ has some positive iterate $n$ with
$|g^{n}(v)|<1$. By compactness of the unit tangent bundle over $\Lambda$
and continuity, there exist $N_0\geq 1$ and $\theta<1$ such that every
$v\in T^1_\Lambda M$ has some $n\leq N_0$ with $|g^{n}(v)|\leq\theta$.

We iterate this. Given $v$, choose $n_1\leq N_0$ with
$|g^{n_1}(v)|\leq\theta$; then choose $n_2\leq N_0$ with
$|g^{n_2}(G^{n_1}(v))|\leq\theta$; and so on. After $k$ blocks the total
time $T_k=n_1+\dots+n_k$ satisfies $k\leq T_k\leq kN_0$ and, by the cocycle
relation, $|g^{T_k}(v)|\leq\theta^{k}\leq\theta^{T_k/N_0}=(\theta^{1/N_0})^{T_k}$.
Writing $a=\sup_{T^1_\Lambda M}\max(|g|,|g|^{-1})$ to control the iterates
between consecutive return times $T_k$, one gets, for every $N\geq 0$,
$$|g^{N}(v)|\leq a^{N_0}\,(\theta^{1/N_0})^{N},$$
so for $N$ large enough (depending only on $N_0$ and $\theta$),
$$\forall v\in T^1_\Lambda M,\qquad |g^{N}(v)|<1.$$
This is impossible: $G^{N}_x$ is a diffeomorphism of the fiber circle
$T^1_xM$ onto $T^1_{f^{N}(x)}M$ with derivative modulus $g^{N}$, so
$\int_{T^1_xM} g^{N}\,d\theta$ equals the length of $T^1_{f^{N}(x)}M$, which
equals that of $T^1_xM$; if $g^{N}<1$ everywhere this integral would be
strictly smaller, a contradiction. Hence $H^+(f,\Lambda)\neq\emptyset$. The
non-emptiness of $H^-(f,\Lambda)$ follows by applying the same argument to
$f^{-1}$.
\end{proof}
\bigskip
The second lemma is a slight variation of a proposition  proved in \cite{PRH} (see the Main Proposition in that paper); for completeness, we provide the proof.

\begin{lemma}\label{DScond} Let $f\in Diff^1(M)$ and let $\Lambda$ be a compact invariant set such that $f_{|\La}$ is far from homotheties. It follows that
$\Lambda$ admits a dominated splitting $T_\Lambda=E\oplus F$, with $\dim(E)=\dim(F)=1$
if and only if the following condition holds.
\begin{description}
\item[(*)] For any  $\delta>0$ arbitrary small, and any  $N\geq 1$ it holds that for any $x\in \Lambda$, there exists $v\in T^1_xM$, such that
for any $m\geq 0$ one has $|g^N(G^m(v))|<(1+\delta)^N$.
\end{description}
\end{lemma}

We prove the `if' direction, the `only if' being straightforward. In terms of the `if' direction, the following remark is essential:

\begin{remark}\label{rmk:single-pair}
The proof of the `if' part of lemma \ref{DScond} only uses condition (*): if there exist $\de$ with $2\de<\de_0$
(where $\de_0$ is the constant in the far from homotheties hypothesis) and
an integer $N\geq 1$ such that for every $x\in\La$ some $v\in T^1_xM$
satisfies $|g^{N}(G^m(v))|<(1+\de)^{N}$ for all $m\geq 0$, then $\La$ admits
a dominated splitting. Consequently, if $\La$ does not admit a dominated
splitting, then for \emph{every} $\de$ with $2\de<\de_0$ and \emph{every}
$N_0\geq 1$ there exists $x\in\La$ such that for any $v\in T^1_xM$ there is
$m\geq 0$ with $|g^{N_0}(G^{m}(v))|\geq (1+2\de)^{N_0}$.
\end{remark}



Before proving the previous lemma, we show how it is used to conclude theorem \ref{p.existence}. The proof needs a classical lemma due to Pliss, of which we state a short version  (for a more general statement see lemma 3.1 in \cite{CP}):
    \vskip 3pt

\paragraph{\bf Pliss's lemma.}{\it    
Given $0 < \gamma_0 < \gamma_1$ and $a > 0$, there exist $n_0 = n_0(\gamma_0,\gamma_1,a)$ and
$l = l(\gamma_0,\gamma_1,a) > 0$ such that for $n\geq n_0$ and  any sequence of numbers $\{a_i\}, 0\leq i\leq n$  satisfying 
$a^{-1} < a_i < a$ and $\prod ^n_{i=0} a_i > \gamma_1^{n}$,  there exist $1 \leq n_1 < n_2 < ....< n_r$ with $r > ln$ such that
$$\prod ^k_{i=n_j+1}  a_i > \gamma_0^{k-n_j},\,\,\,\, n_j < k < n.$$}
    \vskip 3pt



%\begin{lemma}
%Let $\Lambda$ be an invariant set without sinks  admiting a dominated splitting $T_\Lambda=E\oplus F$, with $\dim(E)=\dim(F)=1$
%if (and only if) the following condition holds.
%\begin{description}
%\item[(*)] There are $\delta>0$ and smaller that $\frac{4}{5}b_0$, $N\geq 1$ and, for any $x\in \Lambda$, there is $v\in T^1_xM$, such that
%for any $m\geq 0$ one has $|g^N(G^m.v)|<(1+\delta)^N$.
%\end{description}
%\end{lemma}

\begin{proof}[Proof of theorem ~\ref{p.existence}]




The ``only if" part of the theorem is clear, we will prove the ``if" one.

\medskip
 Let us fix $\delta>0$ small and let us set $N_0$  large. Since
$\La$ is assumed to have no dominated splitting, lemma \ref{DScond}, in
the single-pair form established at the end of its proof,
provides a point $x\in\Lambda$ satisfying the following:

\begin{description}
\item[(**)]
for  any $v\in T^1_xM$
there is $m^+_0\geq 0$ such that $|g^{N_0}(G^{m^+_0}(v))|\geq (1+2\delta)^{N_0}$.
\end{description}
That is possible by remark \ref{rmk:single-pair}: failure of a dominated splitting yields the failure of (*) at every pair $\delta,N_0$ with $2\delta<\delta_0$.

One can also apply lemma~\ref{l.H} to the $\alpha$-limit set of $x$. So, there is a point $y\in \alpha(x)$ and $u\in T_y^1M$ such that for any positive integer $k$ it holds that $|g^{-k}(u)|\geq 1$; so, it also holds that   $|g^{-k}(u)|\geq (1-\de)^k$. Choosing $m^-$ sufficiently large in such a way that $f^{-m^-}(x)$ is sufficiently close to $y$,  and by continuity  one can get a vector $w\in T^1_{f^{-m^-}(x)}M$ close to $u\in T^1_yM$  and $N$ large such that $\forall\, 0\leq k\leq N,\; |g^{-k}(w)|\geq (1-2\delta)^k.$ Moreover, $N$ is arbitrary large provided that $f^{-m-}(x)$ is sufficiently close to $y.$

Now, we take the vector $v:=G^{m^-}(w)\in T^1_xM$. On one hand, it holds that

\begin{equation}\label{e.negative}
\forall\,\, 0\leq k\leq N,\; |g^{-k}(G^{-m^-}(v))|\geq (1-2\delta)^k.
\end{equation}

On the other hand, one can  apply property (**) to $v\in T^1_xM$, and so there is $m^+_0\geq 0$ such that $|g^{N_0}(G^{m^+_0}(v))|\geq (1+2\delta)^{N_0}$. From the fact that $N_0$ was choosen large, one can apply  Pliss lemma and  one deduces that there is $m^+\geq 0$ such that
\begin{equation}\label{e.positive}
\forall\,\, 0\leq k\leq N,\;|g^k(G^{m^+}(v))|\geq (1+\delta)^k.
\end{equation}


Consider the largest integer $m$ satisfying that  $-m^-\leq m\leq m^+$ and \eqref{e.negative} still holds, that means, $m$ satisfies
that 

\begin{equation}\label{e.negative-bis}
 \forall  0\leq k\leq N+ m-m^- , \; |g^{-k}(G^{m}(v))|\geq (1-2\delta)^k. 
\end{equation}



We show by induction on $k=0,\dots, m^+-m$ that $|g(G^{m+k}(v))|> (1+\delta):$ for $k=0$  if $|g(G^{m}(v))|<(1+\delta)$  then $|g^{-1}(G^{m+1}(v))|>(1+\delta)^{-1}$  and since $(1+\de)^{-1}> 1-2\de$, combining with  (\ref{e.negative-bis}) this gives (\ref{e.negative-bis}) but starting at  $m+1$ contradicting the maximal property of $m.$ The inductive step for $k > 0$ is identical, applied at  $m+k$ instead of $m$. Combining with   \eqref{e.positive} this gives
$$\forall \,\, 0\leq k\leq m^+-m+N,\; |g^{k}(G^{m}(v))|\geq (1+\delta)^k.$$
Using the last inequality and (\ref{e.negative-bis}) we have thus found a point $x_{\delta,N}=f^m(x)$ and a direction $v_{\delta,N}$ which satisfy
$$\forall \,\, 0\leq k\leq N,\; |g^{k}(G^{m}(v_{\delta,N}))|\geq (1+\delta)^k,$$
$$\forall\,\, 0\leq k\leq N,\; |g^{-k}(G^{m}(v_{\delta,N}))|\geq (1-2\delta)^k,$$
Considering a limit point $z_N$ and a limit direction $v_N$ (obtained by compactness of $\La$ and the unit tangent bundle) when $\delta\to 0$ one deduces that 
$$\forall \,\, k\in \{-N,\dots, N\},\; |g^{k}(v_N)|\geq 1.$$
By letting $N\to +\infty$ one gets a critical point, more precisely, a diagonal extraction as $N \to \infty$ yields $(z, v_\infty) \in T^1 \Lambda$ with $|g^k(v_\infty)| \geq 1$  for all $k \in \mathbb{Z}$ so $z\in \mathrm{Crit}(f, \Lambda)$.
\end{proof}


\begin{proof}[Proof of lemma ~\ref{DScond}] Recalling how the notion of dominated splitting can be recast in terms of the projective cocycle, to prove that a set $\La$ has a dominated splitting, it is enough to prove that there exist positive integers  $n_0$ and $m_0$ and $\de>0$ such that for any $x \in \La$ it holds that there is $v\in T^1_xM$ such that \begin{equation}\label{dom}
|g^{n_0}(G^m(v))| < (1-\de)^{n_0},\,\,\forall m>m_0.
\end{equation}
Now, by  (*),  taking $\de$  small and $N=1$ (below we consider the general case for any $N\geq 1$),   for any $x$ there is  $v$ such that  $|g(G^m(v))| < (1+\de)$ for any $ m>0.$ Observe that in particular, it holds that  \begin{equation}\label{dom1}|g^n(G^m(v))| < (1+\de)^n,\,\,\,\forall n>0, \,m>0.
\end{equation}
The goal is to show that  this vector actually satisfies equation (\ref{dom}) for some $n_0$. If it is not the case, one gets sequences $(n_i), (m_i)$ such that $|g^{n_i}(G^{m_i}(v))| \geq  (1-\de)^{n_i}$; by Pliss's lemma, for some $\de'>\de$ but close to it, there is a sequence $(m_i')$ such  that $m_i-m_i'\to \infty$ and 
\begin{equation}\label{dom2}|g^n(G^{m_i'}(v))| > (1-\de')^n,\,\,\,\, 0\leq n\leq m_i-m_i'.
\end{equation}
Taking an accumulation point  $z, w$ of the sequence $(f^{m_i'}(x)), (G^{m_i'}(v))$, one get from equations (\ref{dom1}) and (\ref{dom2}) that $(1-\de')^n\leq g^n(w)\leq (1+\de)^n$ and that contradicts the fact that $f_{|\La}$ is far from homotheties.  Choosing $\delta$
(and hence $\delta'$) less than $\delta_0$, where $\delta_0$ is the constant in the far-from-homotheties hypothesis, one gets a contradiction.


Observe that the argument above only used condition (*) for 
$\de,N=1$. For $\de,N$ with $N>1$, the same
proof applies in the same way  to $F=f^{N}$: the projective derivative cocycle of
$F$ is $g_F=g^{N}$, so the hypothesis reads as condition (*) with
 $\de_F$ and $N=1$ for $F$, where $1+\de_F=(1+\de)^{N}$, and
$F_{|\La}$ is far from homotheties by the hypothesis of theorem \ref{p.existence}. A dominated
splitting for $F$ is a dominated splitting for $f$. In particular, if
$\La$ admits no dominated splitting, then condition (*) fails for every
 $\de$ and $N>0$ with $\de$ small.

\end{proof}



\section{Far from homotheties. Proof of theorem \ref{dissipative-FFH}}
\label{sec:dissipative}

Sufficient dissipation is inherited by every power of $f$. Write
$b=\max_{x\in\La}|\det(D_xf)|$ and $M_0=\max_{x\in\La}\|D_xf\|$; so by
submultiplicativity the corresponding quantities for $f^{N}$ satisfy $b_N\leq  b^{N}$ and    $M_{0,N}\leq M_0^{N}$. Therefore, 
$$
\sqrt{b_N}\,M_{0,N}+2b_N
\leq \sqrt b\,M_0+2b<\;\frac{1}{2},
$$
so $f^{N}_{|\La}$ is sufficiently dissipative for every $N\geq 1$.


After this observation, to prove theorem \ref{dissipative-FFH}, it is shown first that for each $N$, $f^N_{|\La}$ is far from homotheties off the basin of finite number of attracting periodic points.   In other words, it is proved  the following weak  version of 
theorem \ref{dissipative-FFH} for any power  $f^N$ of $f$:

\begin{thmB}\label{dissipative-FFH2}

Let $f\in \diff^1(M^2)$ and let $\La$ be a compact invariant set such that
$f_{|\La}$ is sufficiently dissipative. Then  there is a finite number of attracting periodic points of  $f$ such that $f$ restricted to the set of
points of $\La$ that do not belong to the basin of attraction of these finite
attracting periodic orbit is far from homotheties. 
\end{thmB}


Theorem \ref{dissipative-FFH} now follows immediately: {\it since  for each positive integer $N$ it is true that $f^N_{|\La}$ is sufficiently dissipative, one can apply theorem B' and after removing a finite number of attracting periodic points, it is true that $f^N_{|\La}$ is far from homotheties; so, after removing all the attracting periodic points of $f_{|\La}$ one finds that it satisfies that all its powers are far from homotheties.}





One can restate Theorem B' in the following equivalent form, which we then prove: there exist $\de$ and $N_0$ such that for any $x\in \La$ 
\begin{itemize}
    \item[--] if there is $v\in T^1_M$ satisfying   $g^n(v)>(1-\de)^n$ for $n=1\dots N_0$ then there is $0\leq k\leq N_0$ such that $g^k(v)>(1+\de)^k,$ 

    \item[--] if there is $v\in T_M$ satisfying   $g^n(v)<(1+\de)^n$ for $n=1\dots N_0$ then there is $0\leq k\leq N_0$ such that $g^k(v)<(1-\de)^k.$ 
\end{itemize}



The structure of the section is as follows. We first show (lemma \ref{finite-basin}) that a finite orbit with frequent contraction lands in the basin of an attracting periodic point with definite local basin; by compactness, there are only finitely many such attractors (remark \ref{finite-attracting}). The technical heart of the section is proposition \ref{prop:pinning}, which gives a bound on $\|A_2 A_1\|$ for two-step compositions under a pinning condition; the bound is small precisely when the sufficient dissipation condition holds. We deduce theorem \ref{dissipative-FFH} by applying the proposition to consecutive pairs of derivatives along an orbit satisfying \eqref{FFH}, extracting frequent contraction, and invoking lemma \ref{finite-basin}. Subsection \ref{subsec:relaxing} discusses how the threshold $\sqrt{b}M_0 < 1$ might be weakened by considering longer blocks of matrices.

\begin{lemma}\label{finite-basin}
    
Let $f\in Diff(M)$ and let $\la<1$, 
there exists $r>0$ and a positive integer $m$ such that given a finite orbit $\{f^j(x)\}_{j=0}^n$ with $n$ larger than $m$ and a subsequence of iterates $x_{j_i}$ such that   either 
\begin{itemize}
    \item[--] $f(x_{j_i})=x_{j_{i+1}}$ and $||D_{x_{j_i}}f||<\la$ or

    \item[--] $f^2(x_{j_i})=x_{j_{i+1}}$ and $||D_{x_{j_i}}f^2||<\la.$
\end{itemize} 
Then $x$ is in the basin of attraction of a periodic point of period smaller and equal than $m;$ moreover, the basin of attraction of that point contains a ball of radius $r.$
\end{lemma}

\begin{proof} By the continuity of $Df$  follows that there exist $r>0$ and $\la<\la'<1$ such that if $f(x_{j_i})=x_{j_{i+1}}$ then $f(B_r(x_{j_i}))\subset B_{\la'.r}(x_{j_{i+1}})$ and if $f^2(x_{j_i})=x_{j_{i+1}}$ then $f^2(B_r(x_{j_i}))\subset B_{\la'.r}(x_{j_{i+1}})$; since one can cover the set $\La$ with finite attracting  balls of size $r$, the lemma follows.
    
\end{proof}


\begin{remark}\label{finite-attracting} For any $r$, there exists a finite number of attracting periodic points that their basin contains a ball of radius at least $r.$ 
    
\end{remark}

% \noi{\em Proof of theorem \ref{measure-FFH}:}
% We first claim that for any $x\in \Lambda$ and $N\geq 1$,
% there exists $n\geq N$ and $u\in T^1_xM$
% such that $g^n(u)> b^{-n}$.
% If this is not the case, there exits $x\in \Lambda$
% such that for any invariant measure $\mu$ that is limit
% of $\frac 1 n \sum_{k=0}^{n-1} \delta_{f^k(x)}$,
% and any invariant measure $\nu$ which lifts $\mu$ on $T^1_{\lambda} M$,
% we have $\int \log (g) d\nu\leq -\log(b)$. But if $\nu^+,\nu^-$
% are the two ergodic measures which lift $\mu$,
% we get a contradiction:
% $$\int \log (g) d\nu^+=-\int \log (g) d\nu^-=\lambda^+(\mu)-\lambda^-(\mu)>-\log(b).$$

% Let us assume now that the conclusion of the theorem fails.
% For any $n\geq 1$ we have
% $1\leq g^n(v)\leq \lambda^n$ but for infinitely many $n$ we have
% $g^n(u)\geq b^n$. This shows that the angle between $u,v$ is small
% whereas the angle between $G^n(u)$ and $G^n(v)$ is bounded away from zero.

% Let us consider $0<m<n$ such that
% $g^m(u)$ is close to $b^{n/2}$.
% Since $\lambda<b^{1/2}$,
% one deduces that the angle between $G^m(u)$
% and $G^m(v)$ is bounded away from zero.
% As a consequence $g^{n-m}(G^mv)$ is close to
% $1/g^{n-m}(G^mu)$, that is to $b^{-1/2n}$. But
% $$g^{n-m}(G^mv)=g^n(v)/g^m(v)\geq \lambda^{-m}> b^{-1/2},$$
% which is a contradiction.


% \qed


% \noi{\em Proof of theorem \ref{dissipative-FFH}:}  To prove the theorem is enough to show that if there is $x$  such that there exists $v$ satisfying 
% $(1-\de)^n< g^n(v)< (1+\de)^n,\;(**)$ for all $n\leq 0$, then $x$ is in the basin of attraction of a (semi)attracting periodic point.   Let us fix  $\la_1$ such that $\la_1^2< b$ and let $\de=\frac{b}{\la_1^2}-1$. Let us consider the set of points. 


%  Let us take ${(TM)}_1=\{(x,v)\in TM:
% v\in T_xM, |v|=1\}$ and consider the map
% $F\colon {(TM)}_1\to {(TM)}_1$ defined as $F(x,v)=(f(x), G_x(v)).$ Observe
% that this map is a skew product usually called
% the ``dynamical cocycle''. Consider a point
% $(x,v)$ satisfying $(**)$. Let us define the
% measures $\mu_N= (\mu_N^1, \mu^2_N)$ where
% $$\mu_N^1=\frac{1}{N} \sum_{i=0}^N \de_{f^i(x)},$$
% $$\mu_N^2=\frac{1}{N} \sum_{i=0}^N \de_{G^i_{f^i(x)}(v)}.$$
% These measures have a convergent subsequence to an
% invariant measure $\mu$  of $F$ (observe that this measure is a
% measure supported byTM)_1$).   If the projection on $M$ of the
% support of the  measure $\mu$  is supported on the orbit of
% a sink, then the lemma follows. So, let us assume that it
% is not the case. Let us take $g\colon {(TM)}_1\to \RR$ defined by
% $g(x,v)=g_x(v),$ the derivative of $G_x.$ It follows that
% \begin{eqnarray}
% \label{lim1}
% \int \log(g_z(w)) d\mu(z,w)= \lim_{N\to\infty} \int
% \log(g_z(w)) d\mu_N(z,w),
% \end{eqnarray}
% and from the definition of the measure $\mu_N=(\mu^1_N, \mu^2_N),$ it
% follows that
% \begin{eqnarray}
% \label{lim2}
% \int \log(g_z(w)) d\mu_N(z,w)=
% \frac{1}{N} \sum_{i=0}^N \log(g_{f^i(x)}(G^i(v)))=\frac{1}{N}
% \log(g_x^N(v)).
% \end{eqnarray}
% From the fact that
% $$(1-\de_2)^n\leq g^n(v) \leq (1+\de_2)^n$$ for any $n$ large enough,
% it follows from (\ref{lim1}) and (\ref{lim2}) that
% \begin{eqnarray}
% \label{lim3}
% \log(1-\de_2) \leq\int \log(g_z(w)) d\mu(z,w)\leq \log(1+\de_2).
% \end{eqnarray}
% From  the Oseledets Theorem and the equation \ref{eq_f-g}  follows that either
% \begin{enumerate}
% \item $\int \log(g_z(w)) d\mu(z,w)= \la_F-\la_E$ or
% \item $\int \log(g_z(w)) d\mu(z,w)= \la_E-\la_F,$
% \end{enumerate}
% where $\la_E$ is the smallest Lyapunov exponent and $\la_F$ is the largest Lyapunov exponent. Observe that the measure can not be supported in the orbit on a (semi)attracting periodic point and so $\la_F\geq 0,$ then
% it follows that either 
% \begin{enumerate}
% \item $\int \log(g_z(w)) d\mu(z,w)\geq-\la_E\geq  \log(\frac{1}{b})$ or
% \item $\int \log(g_z(w)) d\mu(z,w)\leq \la_E\leq  \log(b).$
% \end{enumerate}
% From the election of $b_0$, $b$ and $\de_2,$  the last two
% inequalities contradict inequality (\ref{lim3}). 





% \qed


\begin{definition}\label{defi:pinning} Let $A_1, A_2$ be a pair of matrices in $Gl(2,\RR), $ we say that they satisfies the \emph{pinning condition} if there exists a unit vector $v \in S^1$ such that
\[
g_{A_1}(v) = 1 \qquad \text{and} \qquad g_{A_2}\bigl(G_{A_1}(v)\bigr) = 1,
\] where $G_{A_i}$ and $g_{A_i}$ are the projective map and projective derivative on the unit circle for $A_1, A_2.$   
    
\end{definition}



The next proposition establish that given a pair of matrices $A_1, A_2$ satisfying the pinning condition and a relation between their norms, then $A_2.A_1$ is a contraction.  The hypotheses on their norms, $M_1=||A_1||, M_2=||A_2||$, below ($M_1 \geq 1$, 
$M_1 M_2 \geq 1$) are dictated by the application in the proof of Theorem \ref{dissipative-FFH}: the proposition will be invoked precisely in the case where neither $A_1$
 alone nor the pair $A_2A_1$ is already contracting.

\begin{proposition}\label{prop:pinning}
Let $b$ and $M_0$ be positive real numbers with $b < \frac{1}{M_0^2}< 1 < M_0.$
Let $A_1, A_2 \in \mathrm{GL}_2(\mathbb{R})$ be invertible matrices satisfying the pinning condition with $|\det A_i| = b$, and singular values $M_i$ (maximal) and $b/M_i$ (minimal) also verifying
\[
M_i \leq M_0 \quad (i = 1, 2), \qquad M_1 \geq 1, \qquad \frac 1{M_1} \leq  M_2.
\]
% For an invertible $2 \times 2$ matrix $A$, define the projective map and projective derivative on the unit circle by
% \[
% G_A(v) = \frac{Av}{\|Av\|}, \qquad g_A(v) = \frac{|\det A|}{\|Av\|^2}.
% \]
 Then it follows that 
\begin{equation}\label{boundA1A2}
\|A_2 A_1\| \leq \sqrt{b}\,(M_1 + M_2) + b\left(\frac{M_1}{M_2} + \frac{M_2}{M_1}\right) +  \frac{b^{5/2}(M_1+M_2)}{(M_1 M_2)^2}.
\end{equation}

In particular,
\begin{eqnarray}\label{boundbM0}
    \|A_2 A_1\| \leq 2\sqrt{b}\,M_0 + 4b.
\end{eqnarray}
\end{proposition}


The statement of the previous proposition may  appear surprising. Each individual matrix $A_i$ has top singular value $M_i$, which under our hypotheses may be as large as $M_0$, so the trivial submultiplicative bound gives only $\|A_2 A_1\| \leq M_1 M_2 \leq M_0^2$, a quantity that is in general much larger than $1$. The theorem replaces this naive bound by an estimate of order $\sqrt b\,(M_1 + M_2)$, which is small under the assumption $2\sqrt b\,M_0 < 1$. The mechanism behind this gain is geometric, and we describe it informally before turning to the proof strategy.


\noindent\emph{The geometric mechanism.} The pinning condition $g_{A_1}(v_1) = 1$ means that $\|A_1 v_1\| = \sqrt b$, a value strictly between the minimal contraction $b/M_1$ and the maximal expansion $M_1$. In the singular frame of $A_1$ (that is described in detail later), such a vector $v_1$ must lie at a very specific angle from the contracting direction $e^1_1$, namely $\arctan(\sqrt b/M_1)$, which is small: $v_1$ is close to the direction of \emph{maximum contraction}. Correspondingly, its image $v_2 = G_{A_1}(v_1)$ in the output frame of $A_1$ is close to the direction of \emph{maximum expansion} $f^2_1$, since the small component of $v_1$ along $e^2_1$ gets amplified by $M_1$ while the large component along $e^1_1$ is crushed by $b/M_1$.

The second part of the pinning, $g_{A_2}(v_2) = 1$, forces $v_2$ to also be close to the minimal contraction direction $e^1_2$ of $A_2$ (by the analogous argument applied to $A_2$). Combining the two observations:

\begin{quote}
The direction of \emph{maximal expansion output} of $A_1$ is approximately aligned with the direction of \emph{minimal contraction input} of $A_2$   (see figure \ref{pinning}).
\end{quote}

This alignment provides the conclusion: The most expanded image of $A_1$ is close to the most contracting direction of $A_2$ (which crushes it by a factor $b/M_2$). Their product is $M_1 \cdot b/M_2 = bM_1/M_2$, a quantity comparable to $b$ rather than $M_1 M_2$. The opposite pairing, bottom output of $A_1$ feeding into top input of $A_2$ — gives the symmetric small quantity $bM_2/M_1$. Both singular values of $A_2 A_1$ end up of order $b$, not $M_i^2$.


\begin{figure}[h]
\centering
\includegraphics[scale=0.7]{fig_pinning.pdf}
\caption{Pinning argument}
\label{pinning}
\end{figure}

\noindent The proof formalizes this alignment by isolating the ``perfect alignment'' configuration and treating the actual configuration as a small perturbation of it.




\vskip 5pt



Let's move now to the  preparative material for the proof of proposition \ref{prop:pinning}.

\paragraph{Singular frames and projective dynamics}



Let $A$ be an invertible $2\times 2$ real matrix with $|\det A| = b > 0$ and singular values $b/M$ and $M$ for some $M > \sqrt{b}$. The singular value decomposition gives orthonormal frames defined below.

\begin{definition}[Singular frames]
The \emph{right singular frame} of $A$ is the orthonormal pair $(e^1, e^2)$ of unit vectors such that:
\begin{itemize}
\item $e^1$ is the direction of \emph{minimal contraction}: $\|A e^1\| = b/M$,
\item $e^2$ is the direction of \emph{maximal expansion}: $\|A e^2\| = M$.
\end{itemize}
The \emph{left singular frame} is the orthonormal pair $(f^1, f^2)$ defined by
\[
f^1 = \frac{A e^1}{\|A e^1\|} = \frac{M}{b}\, A e^1, \qquad f^2 = \frac{A e^2}{\|A e^2\|} = \frac{1}{M}\, A e^2.
\]
\end{definition}

% In these frames, $A$ admits the diagonal representation
% \[
% A = M\, f^2 (e^2)^T + \frac{b}{M}\, f^1 (e^1)^T,
% \]
% equivalently, expressing vectors in the right singular basis and outputs in the left singular basis,
% \[
% A\bigl(\alpha e^1 + \beta e^2\bigr) = \frac{b\alpha}{M}\, f^1 + M\beta\, f^2.
% \]

\paragraph{Projective derivative and projective map in singular frame}

For $v \in S^1$ a unit vector, recall the definition of the projective map and its projective derivative associated to a linear transformation $A$:
$G_A(v) = \frac{Av}{\|Av\|}, \qquad g_A(v) = \frac{|\det A|}{\|Av\|^2} = \frac{b}{\|Av\|^2}.$
Writing $v = \cos\theta\, e^1 + \sin\theta\, e^2$ with $\theta \in [0, 2\pi)$, then
% $\[
% \|Av\|^2 = \frac{b^2}{M^2}\cos^2\theta + M^2 \sin^2\theta,
% \]
% and consequently
% \[
$g_A(v) = \frac{M^2}{b} \frac{1+\tan^2\theta}{1 + \frac{M^4}{b^2}\tan^2\theta}.$

% \begin{proof}
% Direct computation: $Av = (b/M)\cos\theta \, f^1 + M\sin\theta\, f^2$, and the formula follows since $(f^1, f^2)$ is orthonormal.
% \end{proof}

In particular, let $A$ be as above with $b/M < \sqrt{b} < M$. The equation $g_A(v) = 1$, equivalently $\|Av\|^2 = b$, has exactly four solutions on $S^1$, given by $v = \pm v_+$ and $v = \pm v_-$, where
\[
v_\pm = \cos\theta\, e^1 \pm \sin\theta\, e^2 \quad \text{with} \quad \tan\theta = \frac{\sqrt{b}}{M}, \quad \theta \in (0, \pi/2).
\]
% Equivalently,
% $\cos^2\theta = \frac{M^2}{M^2 + b}, \qquad \sin^2\theta = \frac{b}{M^2 + b}.$

% \begin{proof}
% By Lemma~\ref{lem:gA-formula}, $g_A(v) = 1$ amounts to
% \[
% \frac{b^2}{M^2}\cos^2\theta + M^2\sin^2\theta = b.
% \]
% Writing $c = \cos^2\theta$ and $s = 1 - c$:
% \[
% \frac{b^2}{M^2} c + M^2 (1 - c) = b
% \quad \Longrightarrow \quad
% c\left(\frac{b^2}{M^2} - M^2\right) = b - M^2
% \quad \Longrightarrow \quad
% c = \frac{M^2 - b}{M^2 - b^2/M^2} \cdot \frac{M^2}{M^2}.
% \]
% Simplifying the right-hand side:
% \[
% c = \frac{M^2(M^2 - b)}{M^4 - b^2} = \frac{M^2(M^2 - b)}{(M^2 - b)(M^2 + b)} = \frac{M^2}{M^2 + b}.
% \]
% The hypothesis $\sqrt{b} < M$ ensures $M^2 - b > 0$, so the cancellation is valid. The four solutions $\pm v_\pm$ correspond to the four sign choices of $\cos\theta, \sin\theta$.
% \end{proof}

The image $w := G_A(v)$ in terms of the image singular frame is given by
\[
w = \cos\phi\, f^1 + \sin\phi\, f^2 \quad \text{with} \quad \tan\phi = \frac{M}{\sqrt{b}}, \quad \phi \in (0, \pi/2).
\]
% Equivalently, if $\phi = \pi/2 - \theta$ then 
% $\cos^2\phi = \frac{b}{M^2 + b}, \qquad \sin^2\phi = \frac{M^2}{M^2 + b}.$

% \begin{proof}
% From $Av = (b/M)\cos\theta\, f^1 + M\sin\theta\, f^2$ and $\|Av\| = \sqrt{b}$,
% \[
% w = \frac{Av}{\sqrt{b}} = \frac{b\cos\theta}{M\sqrt{b}}\, f^1 + \frac{M\sin\theta}{\sqrt{b}}\, f^2.
% \]
% Writing $w = \cos\phi\, f^1 + \sin\phi\, f^2$ in the left singular frame:
% \[
% \tan\phi = \frac{M\sin\theta/\sqrt{b}}{b\cos\theta/(M\sqrt{b})} = \frac{M^2 \sin\theta}{b\cos\theta} = \frac{M^2}{b}\cdot \frac{\sqrt{b}}{M} = \frac{M}{\sqrt{b}}.
% \]
% Since $\tan\theta\cdot \tan\phi = (\sqrt{b}/M)(M/\sqrt{b}) = 1$, we have $\phi = \pi/2 - \theta$, and the squared cosines/sines follow.
% \end{proof}

\begin{remark}
The angle $\theta$ that $v$ makes with the contracting direction $e^1$ is small (of order $\sqrt{b}/M$), while the angle $\phi$ that $w = G_A(v)$ makes with the bottom output direction $f^1$ is large (close to $\pi/2$). Equivalently, $v$ is close to the contracting direction $e^1$ and $w$ is close to the expanding output direction $f^2$.
\end{remark}

\paragraph{Reduction to canonical coordinates (revised)}


For any orthogonal $Q$ the projective derivative satisfies $g_Q \equiv 1$, so the chain rule
\[
g_{AB}(v) = g_A\bigl(G_B(v)\bigr)\, g_B(v)
\]
implies that the pinning condition is preserved under such modifications.

We exploit this freedom as follows. Choose orthogonal $Q_1$ on the input side of $A_1$ so that the right singular frame of $A_1$ is mapped to the standard basis with the convention $
e^1_1 = (1,0) \quad (\text{minimal contraction}), \, e^2_1 = (0,1) \quad (\text{maximal expansion}).$

Choose orthogonal $Q$ between $A_1$ and $A_2$ so that the left singular frame of $A_1$ is also mapped to the standard basis: $
f^1_1 = (1,0), \qquad f^2_1 = (0,1).$ 
With these choices, $A_1$ takes the diagonal form
\[
A_1 = \begin{pmatrix} b/M_1 & 0 \\ 0 & M_1 \end{pmatrix}.
\]

For $A_2$, choose orthogonal $Q_2$ on the output side so that the left singular frame is mapped to the standard basis with the \emph{swapped} convention
$f^1_2 = (0,1)$ (image of minimal contraction direction),$f^2_2 = (1,0)$ (image of maximal expansion direction).

The right singular frame $(e^1_2, e^2_2)$ of $A_2$ remains a free orthonormal pair, to be determined by the pinning condition.


\paragraph{The matching rotation $O$:} Now we introduce an  orthogonal $O$ is defined to implement the {\it alignment} between maximal expansion of $A_1$ and minimal contraction of $A_2$.  More precisely, let $O \in SO(2)$ be the rotation that maps the vertical direction $(0,1) = f^2_1$ (the image under $A_1$ of its maximal expansion direction) to the minimal contraction direction $e^1_2$ of $A_2$: $ O(0,1)^T = e^1_2.$ By orthogonality of $O$, this also determines $O\,(1,0)^T = e^2_2$ (with sign fixed by the rotation orientation).

Under the pinning condition $g_{A_1}(v_1) = 1$ and $g_{A_2}(G_{A_1}(v_1)) = 1$, with the canonical-coordinate chosen, the rotation $O$ is given by
\[
O = \begin{pmatrix} \cos\theta & \sin\theta \\ -\sin\theta & \cos\theta \end{pmatrix},
\qquad
\theta = \arctan\!\frac{\sqrt{b}}{M_1} + \arctan\!\frac{\sqrt{b}}{M_2}.
\]
Equivalently,using that $\sin\arctan(x) = \frac{x}{\sqrt{1+x^2}},$ and $ \cos\arctan(x) = \frac{1}{\sqrt{1+x^2}},$ applied to  $x = \sqrt b/M_1$ and $x = \sqrt b/M_2$ it follows that 
\[
\sin\theta = \frac{\sqrt{b}\,(M_1 + M_2)}{\sqrt{(M_1^2 + b)(M_2^2 + b)}},
\qquad
\cos\theta = \frac{M_1 M_2 - b}{\sqrt{(M_1^2 + b)(M_2^2 + b)}}.
\]

In fact, let $v_2 := G_{A_1}(v_1)$. We locate $v_2$ from two viewpoints: The angle that $v_2$ makes with the bottom output direction $f^1_1 = (1,0)$ is $\phi_1 = \arctan(M_1/\sqrt b)$. Equivalently, the angle that $v_2$ makes with the vertical direction $f^2_1 = (0,1)$ is
$\frac{\pi}{2} - \phi_1 = \arctan\!\frac{\sqrt b}{M_1}.$ The angle that $v_2$ makes with the minimal contraction direction $e^1_2$ is $\theta_2 = \arctan\!\frac{\sqrt b}{M_2}.$

Choosing the orientation so that $v_2$ lies between the vertical $(0,1)$ and the direction $e^1_2$ on the unit circle, the total angle from $(0,1)$ to $e^1_2$ is the sum: $\theta = \arctan\!\frac{\sqrt b}{M_1} + \arctan\!\frac{\sqrt b}{M_2}.$ 
The matrix $O$ rotating $(0,1)$ to $e^1_2$ by this angle takes the stated form. 

\paragraph{The diagonal matrix $A_2 O A_1$}


Under the canonical coordinates described and with $O$ as defined before,
\begin{equation}\label{diagonal}
A_2 O A_1 = \begin{pmatrix} bM_2/M_1 & 0 \\ 0 & bM_1/M_2 \end{pmatrix}.
\end{equation}
In particular,
\[
\|A_2 O A_1\| = b\,\max\!\left(\frac{M_1}{M_2},\,\frac{M_2}{M_1}\right), \qquad |\det(A_2 O A_1)| = b^2.
\]

% \begin{proof}
% We compute the action of $A_2 O A_1$ on the standard basis.

% The diagonal form of $A_1$ gives
% \[
% A_1 (1,0)^T = (b/M_1)\,(1,0)^T, \qquad A_1 (0,1)^T = M_1\,(0,1)^T.
% \]
% By Definition~\ref{def:O}, $O\,(0,1)^T = e^1_2$ and $O\,(1,0)^T = e^2_2$. Applying $A_2$, and using the swapped convention for $A_2$'s left singular frame in Remark~\ref{rem:canonical}:
% \[
% A_2 e^1_2 = \frac{b}{M_2}\, f^1_2 = \frac{b}{M_2}\,(0,1)^T, \qquad A_2 e^2_2 = M_2\, f^2_2 = M_2\,(1,0)^T.
% \]

% Therefore:
% \[
% (A_2 O A_1)(1,0)^T = A_2\, O\,\bigl((b/M_1)(1,0)^T\bigr) = \frac{b}{M_1}\, A_2 e^2_2 = \frac{b M_2}{M_1}\,(1,0)^T,
% \]
% \[
% (A_2 O A_1)(0,1)^T = A_2\, O\,\bigl(M_1\,(0,1)^T\bigr) = M_1\, A_2 e^1_2 = \frac{b M_1}{M_2}\,(0,1)^T.
% \]
% This establishes the diagonal form.

% The operator norm is the larger diagonal entry, $b\,\max(M_1/M_2, M_2/M_1)$, and the determinant is the product of diagonal entries, $(bM_2/M_1)(bM_1/M_2) = b^2$.
% \end{proof}

\paragraph{Factorization of $A_2 A_1$}


With $D := A_2 O A_1$ then,
\[
A_2 A_1 = D \cdot \bigl(A_1^{-1} O^{-1} A_1\bigr).
\]

% \begin{proof}
% Direct computation:
% \[
% D \cdot \bigl(A_1^{-1} O^{-1} A_1\bigr) = A_2 O A_1 \cdot A_1^{-1} O^{-1} A_1 = A_2 \cdot O O^{-1} \cdot A_1 = A_2 A_1. \qedhere
% \]
% \end{proof}


With $A_1 = \mathrm{diag}(b/M_1, M_1)$ and $O$ rotation by angle $\theta$,
\[
A_1^{-1} O^{-1} A_1 = \begin{pmatrix} \cos\theta & -(M_1^2/b)\sin\theta \\[4pt] (b/M_1^2)\sin\theta & \cos\theta \end{pmatrix}.
\]

% \begin{proof}
% We compute step by step. First,
% \[
% O^{-1} = \begin{pmatrix} \cos\theta & -\sin\theta \\ \sin\theta & \cos\theta \end{pmatrix},
% \qquad
% A_1 = \begin{pmatrix} b/M_1 & 0 \\ 0 & M_1 \end{pmatrix},
% \qquad
% A_1^{-1} = \begin{pmatrix} M_1/b & 0 \\ 0 & 1/M_1 \end{pmatrix}.
% \]
% Then
% \[
% O^{-1} A_1 = \begin{pmatrix} \cos\theta & -\sin\theta \\ \sin\theta & \cos\theta \end{pmatrix}\begin{pmatrix} b/M_1 & 0 \\ 0 & M_1 \end{pmatrix}
% = \begin{pmatrix} (b/M_1)\cos\theta & -M_1 \sin\theta \\ (b/M_1)\sin\theta & M_1\cos\theta \end{pmatrix},
% \]
% and finally
% \[
% A_1^{-1} O^{-1} A_1 = \begin{pmatrix} M_1/b & 0 \\ 0 & 1/M_1 \end{pmatrix}\begin{pmatrix} (b/M_1)\cos\theta & -M_1 \sin\theta \\ (b/M_1)\sin\theta & M_1\cos\theta \end{pmatrix}
% = \begin{pmatrix} \cos\theta & -(M_1^2/b)\sin\theta \\[4pt] (b/M_1^2)\sin\theta & \cos\theta \end{pmatrix}.
% \]
% A consistency check: $\det(A_1^{-1} O^{-1} A_1) = \cos^2\theta + \sin^2\theta = 1$, as expected for the similarity transformation of a rotation.
% \end{proof}

In the canonical coordinates,
\begin{equation}\label{A2A1}
    A_2 A_1 = \begin{pmatrix} (bM_2/M_1)\cos\theta & -M_1 M_2\,\sin\theta \\[4pt] (b^2/(M_1 M_2))\,\sin\theta & (bM_1/M_2)\cos\theta \end{pmatrix},
    \end{equation}

with $\theta= \arctan\!\frac{\sqrt b}{M_1} + \arctan\!\frac{\sqrt b}{M_2}.$
% \begin{proof}
% By Lemmas~\ref{lem:factorization} and~\ref{lem:conjugate},
% \[
% A_2 A_1 = \begin{pmatrix} bM_2/M_1 & 0 \\ 0 & bM_1/M_2 \end{pmatrix}\begin{pmatrix} \cos\theta & -(M_1^2/b)\sin\theta \\[4pt] (b/M_1^2)\sin\theta & \cos\theta \end{pmatrix}.
% \]
% Computing entry by entry:
% \begin{align*}
% (A_2 A_1)_{11} &= \frac{bM_2}{M_1}\cdot \cos\theta + 0 = \frac{b M_2}{M_1}\cos\theta, \\[4pt]
% (A_2 A_1)_{12} &= \frac{bM_2}{M_1}\cdot\left(-\frac{M_1^2}{b}\sin\theta\right) + 0 = -M_1 M_2\,\sin\theta, \\[4pt]
% (A_2 A_1)_{21} &= 0 + \frac{bM_1}{M_2}\cdot\frac{b}{M_1^2}\sin\theta = \frac{b^2}{M_1 M_2}\sin\theta, \\[4pt]
% (A_2 A_1)_{22} &= 0 + \frac{bM_1}{M_2}\cdot\cos\theta = \frac{bM_1}{M_2}\cos\theta.
% \end{align*}
% Assembling these gives the stated matrix.

% A consistency check on the determinant:
% \[
% \det(A_2 A_1) = \frac{b M_2}{M_1}\cdot\frac{b M_1}{M_2}\cos^2\theta - \left(-M_1 M_2 \sin\theta\right)\cdot\frac{b^2}{M_1 M_2}\sin\theta = b^2\cos^2\theta + b^2\sin^2\theta = b^2,
% \]
% matching $|\det(A_2 A_1)| = |\det A_2||\det A_1| = b\cdot b = b^2$.
% \end{proof}
Observe that the largest entry, the off-diagonal $M_1 M_2 \sin\theta$, is bounded by $\sqrt b\,(M_1 + M_2)$ because the explicit value of $\sin\theta$ contains a compensating factor $1/(M_1 M_2)$: the smallness of the rotation angle $\theta$ exactly cancels the potentially large product $M_1 M_2$. The explicit estimates are done in what follows, using the Frobenius norm.

\begin{proof}[Proof of proposition \ref{prop:pinning}] First, we estimate the entries of $A_2.A_1$ given by equation \ref{A2A1}, and we bound each term.

\noindent\emph{Step 1: bound on each entry.}
Using $|\cos\theta| \leq 1$ and the explicit value of $\sin\theta$ ,
\[
|\sin\theta| = \frac{\sqrt b\,(M_1+M_2)}{\sqrt{(M_1^2+b)(M_2^2+b)}} \leq \frac{\sqrt b\,(M_1+M_2)}{M_1 M_2},
\]
since $(M_i^2 + b) \geq M_i^2$. Therefore:
\begin{align*}
\bigl|(A_2 A_1)_{11}\bigr| &\leq \frac{b M_2}{M_1}, \\[4pt]
\bigl|(A_2 A_1)_{12}\bigr| = M_1 M_2 |\sin\theta| &\leq \sqrt b\,(M_1+M_2), \\[4pt]
\bigl|(A_2 A_1)_{21}\bigr| = \frac{b^2}{M_1 M_2}|\sin\theta| &\leq \frac{b^{5/2}(M_1+M_2)}{(M_1 M_2)^2}, \\[4pt]
\bigl|(A_2 A_1)_{22}\bigr| &\leq \frac{b M_1}{M_2}.
\end{align*}

\medskip

\noindent\emph{Step 2: bound the operator norm by the Frobenius norm.} Using $\|A_2 A_1\| \leq \|A_2 A_1\|_F = \sqrt{\sum_{i,j} |(A_2 A_1)_{ij}|^2}$:
\[
\|A_2 A_1\|_F^2 \;\leq\; b(M_1+M_2)^2 \;+\; b^2\!\left(\frac{M_1^2}{M_2^2} + \frac{M_2^2}{M_1^2}\right) \;+\; \frac{b^5(M_1+M_2)^2}{(M_1 M_2)^4}.
\]

\medskip

\noindent\emph{Step 3: take the square root.} Using $\sqrt{x_1 + x_2 + x_3} \leq \sqrt{x_1} + \sqrt{x_2} + \sqrt{x_3}$ for $x_i \geq 0$,
\[
\|A_2 A_1\|_F \;\leq\; \sqrt b\,(M_1+M_2) \;+\; b\sqrt{\frac{M_1^2}{M_2^2} + \frac{M_2^2}{M_1^2}} \;+\; \frac{b^{5/2}(M_1+M_2)}{(M_1 M_2)^2}.
\]

\medskip

\noindent\emph{Step 4: simplify the middle term.} Using $\sqrt{a^2 + c^2} \leq a + c$ for $a, c \geq 0$,
\[
\sqrt{\frac{M_1^2}{M_2^2} + \frac{M_2^2}{M_1^2}} \;\leq\; \frac{M_1}{M_2} + \frac{M_2}{M_1}.
\]


Combining,
\[
\|A_2 A_1\| \;\leq\; \sqrt b\,(M_1+M_2) \;+\; b\!\left(\frac{M_1}{M_2} + \frac{M_2}{M_1}\right) \;+\; \frac{b^{5/2}(M_1+M_2)}{(M_1 M_2)^2}. 
\]

\noindent\emph{Step 5: bounding in terms of $b$ and $M_0$.}  A simple calculation shows that the right hand of the equation (\ref{boundA1A2})  on the rectangle $[1, M_0] \times [1/M_0, M_0]$ is realized at the corner $(M_0, M_0)$, so the maximum value is 
\[
f(M_0, M_0) = 2\sqrt b\, M_0 + 2b + \frac{2 b^{5/2}}{M_0^3}, 
\]
and from the fact that $M_0\geq \sqrt b,$ it holds that $\frac{2 b^{5/2}}{M_0^3}< 2. b$, and so 
\[
f(M_0, M_0) \leq 2\sqrt b\, M_0 + 4b. 
\]



\end{proof}

\vskip 5pt


\medskip

To understand the bound (\ref{boundA1A2}) observe that the term $\sqrt b\,(M_1+M_2)$  comes from the off-diagonal entry $M_1 M_2 \sin\theta$. The product $M_1 M_2$ is large, but $\sin\theta \approx \sqrt b\,(1/M_1 + 1/M_2)$ is small, and the cancellation gives $\sqrt b\,(M_1+M_2)$. This is the dominant term and reflects the imperfection of the alignment between the direction of maximal expansion of $A_1$ and minimal contraction of $A_2.$
 The  term $b(M_1/M_2 + M_2/M_1)$  comes from the diagonal entries, which are essentially the singular values of the perfect-alignment matrix $D$. This is the genuine size of the composition when the alignment is exact, and it dominates only when the matrices are very ``unbalanced'' relative to each other ($M_1 \gg M_2$ or vice versa).

 The condition $\sqrt b\,M_0 < 1$ ensures both terms remain manageable: it guarantees that the rotation angle is genuinely small (controlling the first term) and that the unbalance ratio $M_1/M_2$ is not too extreme (controlling the second) and this is done in the step 5 above.

 Now, before completing the proof of theorem \ref{dissipative-FFH}, we provide two remarks, the first one is about relaxing the pinning condition and the second one is about considering matrices that their determinant are not necessarily the same.


\begin{remark}\label{rmk:pinning} Theorem \ref{dissipative-FFH} assumes that the right-hand side of (\ref{boundbM0}) is smaller than one. Choose $\la_0<1$ and $\la_1$ also smaller than one but close to one and $\de>0$ small,  such that the right hand  (\ref{boundbM0}) is smaller than $\la_0$. In  inequality \ref{boundA1A2}, if one relaxes  the pinning condition to assume  that $1-\de < g_{A_1}(v)< 1+\de$ and $1-\de < g_{A_2}(G_{A_1}(v))< 1+\de$, and   if one  takes $M_1>\la_1$ and $M_2.M_1>\la_1$ it follows that the right-hand side of (\ref{boundA1A2}) is smaller than $\la_0.$ 
    
\end{remark}


\begin{remark}\label{rmk:det-normalization}
The proposition is stated for matrices with equal determinants $|\det
A_i|=b$; in the proof of theorem \ref{dissipative-FFH} it is applied to
derivatives along an orbit, whose determinants $b_i$ satisfy only $b_i\leq
b$. All the estimates in the proof are monotone non-decreasing in each
$b_i$ (in particular $\sin\theta=\frac{\sqrt{b_1}M_2+\sqrt{b_2}M_1}
{\sqrt{(M_1^2+b_1)(M_2^2+b_2)}}\leq\frac{\sqrt
b\,(M_1+M_2)}{M_1M_2}$), so the bounds (\ref{boundA1A2}) and
(\ref{boundbM0}) hold with $b=\max_{x\in\La}|\det(D_xf)|$.
\end{remark}


\begin{proof}[Proof of theorem B'] Let us consider $\la_0, \la_1$ as in the previous remark,  $\la=\max\{\la_0, \la_1\}$, and a finite orbit of a point $x$ such that there exists a vector $v$ satisfying  
\begin{equation}\label{flex-pinning}
    1-\de< g_{f^k(x)}(G^k(v))<1+\de,\,\,\mbox{for}\, k=0\dots m.
\end{equation}
    
 By induction, one can chose a subsequences of iterates  $\{x_{i_j}\}$ such that   either 
\begin{itemize}
    \item[--] $f(x_{i_j})=x_{i_{j+1}}$ and $||D_{x_{i_j}}f||<\la$ or

    \item[--] $f^2(x_{i_j})=x_{i_{j+1}}$ and $||D_{x_{i_j}}f^2||<\la,$
\end{itemize}  

Once $x_{i_j}$ has been chosen, we distinguish three cases according the norm of $D_{x_{i_j}}f$. Case (i), if $||D_{x_{i_j}}f||<\la_1$ then $x_{i_{j+1}}=f(x_{i_j})$. Case (ii),  if $||D_{x_{i_j}}f|| \geq \la_1$ and  $||D_{x_{i_j}}f||||D_{f(x_{i_j})}f||<\la_1<1$ then one takes $x_{i_{j+1}}=f^2(x_{i_j})$ and $||D_{x_{i_j}}f^2||<\la<1$. Case (iii)  if $||D_{x_{i_j}}f||||D_{f(x_{i_j})}f||\geq \la_1$ one can apply remark  \ref{rmk:pinning} and again it follows that $||D_{x_{i_j}}f^2||=||D_{x_{i_j}}f.D_{f(x_{i_j})}f||<\la_0<1.$

We now  apply lemma \ref{finite-basin}, it follows that if $n\geq m$ then  $x$ is in the basin of attraction of an attracting  periodic point that has a local basin containing a ball of radius $r$. By remark \ref{finite-attracting}, there are finitely many  attracting periodic points. So, if we considered $\La'$, the complement in $\La$ of the basins of attraction of these finite attracting periodic points, it holds that that in $\La'$ there are not finite trajectories of size larger than $m$ satisfying (\ref{flex-pinning}); therefore, taking $0<\de'$ smaller than $\de$, there is not trajectory of size larger than $m_0$ satisfying $(1-\de')^k< g^k(v)< (1+\de)^k$ for $k=0\dots n.$ In particular, no infinite forward orbit in $\Lambda'$ lies in the window, so $f|_{\Lambda'}$ is far from homotheties.


\end{proof}
\medskip


\begin{remark} Observe that the condition $\sqrt b\, M_0 < 1$ serves two purposes. In the context of proposition \ref{prop:pinning} it ensures that the rotation angle $\theta$ needed for alignment  is small enough for the optimization of $b$ and $M_0$. In the degenerate case that $\frac{b}{M_2}=M_2$, it implies that $A_2$ is $\sqrt b. O$ where $O$ and so $||A_2||=\sqrt b$  and so $||A_2.A_1||<1.$
    
\end{remark}

\subsection{Relaxing the  sufficiently dissipative condition}\label{subsec:relaxing}


The sufficient dissipation condition (\ref{cond-diss-FFH}) is used in this
section only through proposition \ref{prop:pinning}, whose threshold
$\sqrt b\,M_0<1$ is what forces the quantitative form of (\ref{cond-diss-FFH}).
We conjecture that this can be improved to only require dissipation.

\begin{conjecture}\label{conj:relaxing}
Theorem \ref{dissipative-FFH} holds assuming only that $f_{|\La}$ is
dissipative, that is, $b=\max_{x\in\La}|\det(D_xf)|<1$, without the
quantitative condition (\ref{cond-diss-FFH}).
\end{conjecture}



The mechanism behind proposition \ref{prop:pinning} is that a pinned pair of
matrices, after the optimal alignment, contracts at a rate controlled by $b$
rather than by the individual norms $M_i$. The threshold $\sqrt b\,M_0<1$
measures how much the imperfection of the alignment is allowed to relax over a
\emph{single} pair. The following non-rigorous computation suggests that, by
grouping the cocycle into longer blocks, the relaxation is getting better: the
 threshold should relax from $\sqrt b\,M_0<1$ towards to  $b<1$ as the block length grows.


\emph{The heuristic.} Consider $2^n$ matrices $A_1,\dots,A_{2^n}$ with
$|\det A_i|=b$, singular values $b/M_i\le M_i$, $M_i\in[1,M_0]$, and the
sequences of  pinning $g_{A_1}(v_1)=1$, $v_{i+1}=G_{A_i}(v_i)$,
$g_{A_{i+1}}(v_{i+1})=1$. Assume at each pairing step that the matching
rotation between adjacent singular frames is the identity (perfect
alignment), so that each pairwise product is the diagonal matrix
(\ref{diagonal}), with top singular value at most $b\,M_0$. Regrouping into
blocks of $2^k$ and iterating the pairing, the chain rule transmits the
pinning to the pair-matrices, the determinant exponent doubles at each
level while the singular-value ratios inherit a comparable bound, so a block
of $2^k$ matrices has top singular value at most $b^{2^{k-1}}M_0$ under
perfect alignment. The criterion $b^{2^{k-1}}M_0<1$, that is
$b<M_0^{-1/2^{k-1}}$, relaxes towards $b<1$ as $k\to\infty$.

\emph{What a proof would require.} Making this rigorous means controlling, at
each level of pairing, the deviation of the matching rotations from the
identity (this is the analogue of the perturbation analysis done  for a
single pair in the proof of proposition \ref{prop:pinning}) with the added
difficulty that the singular-value ratios entering level $k+1$ depend on the
alignment errors accumulated through level $k$. The errors must be shown to
decay fast enough across levels that their product remains controlled; we do
not know whether the geometric gain $b^{2^{k-1}}$ in the determinant always
dominates this accumulation. We leave this analysis for future work.



% The threshold $\sqrt b\, M_0 < 1$ in proposition \ref{prop:pinning} can plausibly be relaxed when the chain is extended to more matrices, as the following non-rigorous calculation suggests. Consider $2^n$ matrices $A_1, \ldots, A_{2^n}$ with $|\det A_i| = b$, singular values $b/M_i \leq M_i$, $M_i \in [1, M_0]$, and the chained pinning $g_{A_1}(v_1) = 1$, $v_{i+1} = G_{A_i}(v_i)$, $g_{A_{i+1}}(v_{i+1}) = 1$.

%  At each pairing step described below, one can assume that  the relevant matching rotation between adjacent singular frames is  the identity (a perfect alignment  between most expanding direction and most contracting ones across pairs), so that each pairwise product is captured by the diagonal matrix (\ref{diagonal}). The diagonal matrices have top singular values bounded by $b\,\max(M_{2i-1}/M_{2i}, M_{2i}/M_{2i-1}) \leq b\, M_0$. The criterion for $\|B^{(1)}_i\| < 1$ in this regime is $b\, M_0 < 1$.

% Now, we consider regrouping  matrices in blocks of $2^k$, (level $k$). The chain rule for projective derivatives transmits the pinning condition to the pair-matrices $B^{(1)}_i$, so the construction iterates: pair the $B^{(1)}_i$ into blocks of four, then eight, and so on. At each level, the determinant exponent doubles ($b, b^2, b^4, \cdots$), while the singular-value ratio at the next level inherits a comparable bound (controlled by the orderings carried through the pairing). The resulting top singular value of a block of $2^k$ matrices, under perfect alignment, is bounded by $b^{2^{k-1}}\, M_0$, providing  the criterion
% \[
% b^{2^{k-1}}\, M_0 < 1, \quad \text{equivalently} \quad b < M_0^{\frac{1}{2^{k-1}}}.
% \]



% A rigorous version of this heuristic would require tracking, at each level of pairing, the corrections from the rotations not being exactly identity, along the lines of the proof of proposition \ref{prop:pinning} but with the additional complication that the singular-value ratios at subsequent levels depend on the previous level. We leave this analysis for future work.




% \begin{lemma}\label{gn some growth}Let $f\in Diff^1(M)$ and let $\La$ be a compact invariant such that $f_{|\La}$ is sufficiently dissipative. If  $(x, v)$ is a critical point,  then there exist two infinite sequences of integers $(n_i^+), (n_i^-)$ such that 
% \begin{enumerate}

% \item[--]  $g^{n_i^+}(v)> e^{\ga.n_i^+}$,

% \item[--]  $g^{-n_i^+}(v)> e^{\ga.n_i^-}$.


% \end{enumerate}

 

% \end{lemma}






% \begin{quest}
% Could we give a more quantitative result?
% For instance does there exists $K\geq 1$ such that for any $N\geq1$,
% there exists $n\in[N,KN]$ such that $g^n(v)\geq \lambda^n$?
% \end{quest}


% \bigskip

 % Therefore, if $A=\left(\begin{array}{ccccc} a & b \\  c & d \end{array}\right);$ then, the two directions are given by the values of $\theta$ such that  $\tan(2\theta)=\frac{2(ab+cd)}{a^2+c^2-(b^2+d^2)}$. 

\section{Properties of the critical set}\label{sec:properties}  
\label{s.properties}

% To prove proposition \ref{properties} we need some preliminaries relating  critical direction with eigenspaces and  singular  directions.

% Given $A\in GL(2, \RR)$, if $v$ is an eigendirection with eigenvalue $\la$ and  there is another eigendirection with eigenvalue $\si$ it follows that $g(v)=\frac{\si}{\la}.$
% On the other hand, one can use the  singular frames   defined before $e^1, e^2$, in particular, it follows that   $g(e^2)= \frac{1}{g(e^1)}$ and 
% \begin{equation}\label{min contraction}\frac{1}{g(e^1)}\leq g(w)\leq g(e^1), \,\,\,\mbox{for any}\,\,w.
%  \end{equation}


\begin{proof}[ Proof of proposition \ref{properties}:]   

% To prove the first  item, for each positive integer $n$ we use the singular frame $\{e^1_n, e^2_n\}$ of $df^n$ and  observe that there exists arbitrarily large $n$ such that $g^n(e^2_n)<(1-\de)^n$ and so $g^n(e^1_n)>\frac{1}{(1-\de)^n}.$ Let $a=|Df^n(e^1_n)|$ and $b=|Df^n(e^2_n)|$,  using the frames  $\{e^1_n, e^2_n \}$ and $\{f^1_n,f^2_n\}$ then  $Df^n=\left(\begin{array}{ccccc} a &  0\\  0 & b \end{array}
% \right);$ therefore, given a vector $v=(1, \de)$, then $Df^n(v)= (a, b.\de)$ and $g^n(v)= \frac{a.b}{a^2+b^2\de}$. If $v$ is the critical direction, it follows that $g^n(v)= \frac{a.b}{a^2+b^2\de}\geq 1$ and so $\de< \frac{a}{b}.\sqrt{1-\frac{a}{b}}$ and given that $g^n(e^2_n)=\frac{a}{b}$, which is assumed to be smaller than $(1-\de)^n$, it follows that  $\de<(1-\de)^n$ and therefore, $v$ is $(1-\de)^n-$close to the most contractive direction of $Df^n.$ 

To prove the first item, observe that if $f^n(c)=c$ and $v$ is the critical direction, then $|g^{k.n}(v)|\geq 1$ for any integer $k$. On the other hand, since $G^n:\TT^1\to \TT^1$ is a Mobius transformation, then $G^n$ is either hyperbolic, elliptic or parabolic. If it is hyperbolic, then either $v$ is a fixed point of $G^n$ (one repeller and one attractor, at the repeller the projective derivative grows exponentially, and in the other decrease exponentially to zero) or its forward orbit converges to the attracting  fixed point of $G^n$ and so the projective derivative  decreases exponentially to zero; in both cases, it can not hold that $|g^{k.n}(v)|\geq 1$ for every integer $k$. Therefore, $G^n$  can only be either parabolic or elliptic and so the first item of the proposition is proved.

The second item follows from the fact that if all iterates $(f^n(c), G^n(v))$ critical points and critical directions, then $g(G^n(v))=1$ for any integer $n$ and so $g^n(v)=1$ for all $n$ contradicting that $f$ is far from homotheties.


The last item follows immediately by continuity: if $(c_n)$ is a sequence of critical points of a sequence of diffeomorphisms $(g_n)$ converging to a diffeomorphism $f$, it follows that any accumulation point of the sequence $(c_n)$ is a critical point of $f.$   


\end{proof}

One can get a quantitative version of  the second item of proposition \ref{properties}:


% \begin{lemma}\label{en} Let $f\in Diff^1(M)$ and let $\La$ be a compact invariant and $\ga>0$ such that for any invariant measure, the difference of the lyapunov exponents is larger than $\ga$. 

% If  $(x, v)$ a critical point,  then there exist two infinite sequences of integers $(n_i^+), (n_i^-)$ such that 
% \begin{enumerate}

% \item  $slope(e_{n_i^+}, v)< e^{-\ga.n_i^+}$ where $e_{n_i^+}$ is the maximal contracting direction of $Df^{n_i^+},$

% \item  $slope(e_{n_i^-}, v)< e^{-\ga.n_i^-}$ where $e_{n_i^-}$ is the maximal contracting direction of $Df^{-(n_i^-)}.$


% \end{enumerate}

% \end{lemma}


\begin{proof}[Proof of  proposition \ref{pr:compatibility}] Observe that given $\eps'>0$ there exists a positive integer $N_1$ such that if $n$ is large enough such that $f^{-j}(x)\notin B_\de({\rm Crit}(f,\La))$ for any $0\leq j\leq n$ then there is $n^*\leq N_1$ such that $f^{-n^*}(x)\in B_{\eps'}(\La_\de).$ Moreover, there is $m^*= m^*(n)$ such that $m^*(n)\to +\infty$ as $n\to \infty$ and for any $n^*\leq  j \leq m^*$ it holds that $f^{-j}(x)\in B_{\eps'}(\La_\de).$

Now, provided $\eps$, there is  $\eps'$ small $0<\ga<1$ and  $N_2$, such  that if  $f^{-j}(y)\in B_{\eps'}(\La_\de)$ for all $0\leq j\leq n$ with $n\geq N_2$ then for any $w\in T_yM$ such that $slope(w, F)>\eps$ follows $g^{-n}(w)<\ga^n.$

Now, we take $N_0$ such that if $n\geq N_0$ then it holds that for any critical $c$ if $f^{-j}(c)\notin B_\de({\rm Crit}(f,\La)), 0\leq j\leq n$ then 
\begin{itemize}
    \item[--] there is a positive integer $n^*\leq N_1$ such that  $f^{-n^*}(c)\in B_{\eps'}(\La_\de),$
    \item[--] $f^{-j}(c)\in B_{\eps'}(\La_\de)$ for any $n^*\leq j\leq m^*,$
    \item[--] $m^*-n^*\geq N_2.$
\end{itemize}
Since $n^*$ is uniformly bounded for any critical point, $N_2$ can be taken large such that  $g^{-n^*}(v_c)\ga^{N_2}<1,$ therefore, if $slope(G^{-n^*}(v_c), F)>\eps,$ it follows that  $g^{m^*}(v_c)= g^{m^*+n^*}(G^{-n^*}(v_c)).g^{-n^*}(v_c)<1,$ contradicting that $v_c$ is the critical direction.

The proof for forward iterations is similar.

\end{proof}

\vskip 5pt


Corollary \ref{cor:compatibility} follows from the proof of the previous proposition. Suppose  the conclusion of the corollary. Take $\de_2$ sufficiently small,   one gets a critical point whose pre-iterate   (the first backward visit to a neighborhood of $\La_{\de_1}$) carries a critical direction not close to the center-unstable subbundle. This contradicts Proposition \ref{pr:compatibility}.

\vskip 5pt

 \begin{proof}[Proof of proposition \ref{exponential recurrence}]
 
 
 \noindent 1. We consider $\varepsilon>0$ and a
 critical point $c\in \La$ associated with a direction $u_0$, and with recurrence $(1-\varepsilon)^n$: for some large times $n$,
it satisfies $f^n(c)\in B(c,(1-\varepsilon)^n)$.
\medskip

\noindent
Since $\La$ is dissipative, $|\det Df^j |_\La|<b^j$ for some $b\in (0,1)$ and $j\geq 1$.
\medskip

\noindent
Let $\Delta$ be an upper bound on derivatives of $G$ and $\log\|Df\|$ and $\lambda$ a lower one on $Df$.
\bigskip

\noindent
2. We choose successively $\delta,\delta',\alpha,\widehat \delta,\gamma,\beta,\rho,\eta>0$ small such that
$\widehat \delta<\delta'<\delta<\varepsilon/2$ and
\begin{equation}\label{e.choice-alpha}
(1-\varepsilon)<(1-\delta')\lambda^{\alpha},\qquad b^\alpha(1-\rho)^{-\alpha}<1-\widehat \delta,\qquad \Delta^{\gamma\alpha}(1-\delta')(1+\eta)^{\alpha}<(1-\widehat \delta),
\end{equation}
\begin{equation}\label{e.choice-beta}
\Delta^{\beta} (1-\widehat \delta)<1, \qquad b^{\beta/4}(1+\eta)<(1-\eta)^2,
\end{equation}
\begin{equation}\label{e.choice-eta}
1-\rho>b^{\gamma/3}(1+\eta)^{(1-\gamma)},\qquad (1+\eta)^2(1-\rho)^{\min(\alpha,\beta)}<1-\eta.
\end{equation}

\noindent
3. We will consider a large return time $n$ and introduce the ball $B:=B(c,(1-\delta)^n)$.\newline
%It remains at comparable size during time $2n$:
Since $\underset{k\to +\infty}\limsup \tfrac 1 k \log \|Df^k|_\La\|\leq 0$,
there is $N\geq 1$ such that for $x$ in a neighborhood of $\La$,
\begin{equation}\label{e.bounds}
\|Df^{N}(x)\|\leq (1+\eta)^N.
%\text{ and } \|Df^{-N}(x)\|^{-1}\geq b^N.
\end{equation}
In particular, there exists $k_0$ such that
$$\forall k_0\leq k\leq 2n,\qquad f^k(B(c,(1+\delta)^n))\subset B(f^k(c),(1+\delta')^n).$$

\noindent
4. We define $m:=[\alpha n]$ and the ball $B':=B(f^{n-m}(c),(1-\delta')^n)$ so that $f^{n-m}(B)\subset B'$.\hspace{-1cm}\mbox{}
\smallskip

\noindent
From the choice of $\alpha$ in~\eqref{e.choice-alpha}, the bounds in~\eqref{e.bounds}, the recurrence in \S1, we have $c\in f^{m}(B')$.
\bigskip

\noindent
5. Let us assume that $\|Df^m\|<(1-\rho)^m$ on $B'$. The choice of $m$ and estimates in~\eqref{e.choice-eta}, \eqref{e.bounds} give
$\|Df^n|_B\|<(1+\eta)^n(1-\rho)^{\alpha n}<1.$
Since $0<\delta<\varepsilon$ and $d(f^n(c),c)<(1-\varepsilon)^n$,
the map $f^n$ is a contraction from $B$ into itself, $B$ is included in the basin of a sink and the proposition holds in this case.
\smallskip

\noindent
One can thus assume: \emph{There exist $y_0\in B'$ and a unit vector $v_0$ with} 
\begin{equation}\label{e.g}
\|Df^m.v_0\|\geq (1-\rho)^m.
\end{equation}



\noindent
6. {\bf Claim 1.}
\emph{For any $z\in f^{m}(B')$, the linear map $Df^{-m}(z)$ is not a homothety.}
\medskip

\noindent
Let $e_z$ be the less expanded direction for $Df^{-m}(z)$  (it maximizes $g^{-m}$).
Note that from \S 5, the direction $e_c$ is defined at $c$.
\medskip

\noindent
\quad {\bf Claim 2.}
\emph{The oscillations of $e_z$ on $f^{m}(B')$ are bounded by $4(1-\widehat \delta)^n$.}
\bigskip

\begin{proof}[{\rm 7.} Proof of claim 1.]
We consider $y_0$ with a direction $v_0$ as in \S5. At any $y'\in B'$ let
$v'_0$ be the direction parallel to $v_0$.
The iterates of $v_0,v'_0$ under $G^k$ are denoted $v_k,v'_k$.
\medskip

\noindent
The choice of $y_0,v_0$ and~\eqref{e.g} give $|g^m(v_0)|/|\det Df^m(y_0)|\leq (1-\rho)^{-2m}.$
\smallskip

\noindent
We decompose the orbits of $y_0, v_0$ up to time $m$
into blocks of size $N$. More precisely, to control $||Df^m||$ we split the orbit into segments of length $N$, on each of which the upper bound (\ref{e.bounds}) gives a uniform estimate. The fraction $\chi$ of 'bad' blocks (where the bound is violated) is controlled by a chain-rule estimate combined with the fact that  $|\det Df^m| < b^m$. By~\eqref{e.bounds} for any $x$ in a neighborhood of $\La$ and any direction $v$,
one has $|g^N(v)|/|\det Df^N(x)|\geq (1+\eta)^{-2N}.$
\smallskip

\noindent
Hence the density $1-\chi$ of blocks satisfying $|g^N(v_k)|/|\det Df^N(y_k)|\leq b^{-2N/3}$ satisfies
$$ (1-\rho)^{-2}\geq b^{-\tfrac 2 3\chi}(1+\eta)^{-2(1-\chi)}.$$
From~\eqref{e.choice-beta}, $\chi<\gamma$.
Moreover, for these good blocks, one gets a contraction $|g^N(v_k)|\leq b^{N/3}$.\hspace{-1cm}\mbox{}
\medskip

\noindent
Let us inductively compare $v_k$ and $v'_k$.
Note that $y_k,y'_k$ are at distance smaller than $(1-\delta')^n(1+\eta)^k$, which bounds the difference between
$Df^N(y_k)$ and $Df^N(y'_k)$.
For all the good blocks the angle between $v_k$ and $v'_k$ decreases after iteration by $Df^N(y_k)$.
For the other blocks, the angle may increase by a factor bounded by $\Delta^N$.
\medskip

\noindent
The bound $\gamma$ on the density and~\eqref{e.choice-alpha} imply that $\angle(v_k,v'_k)$ for $0\leq k\leq m$ is bounded by
$$\Delta^{\gamma m}(1-\delta')^n(1+\eta)^m<(1-\widehat \delta)^n.$$
Since the orbits of $v_0,v'_0$ remain $(1-\widehat \delta)^n$-close during $m$ iterates,
\begin{equation}\label{e.expansion}
\|Df^m.v_0'\|\geq \exp(m.\Delta.(1-\widehat \delta)^n)\|Df^m.v_0\|\geq \tfrac 1 2 (1-\rho)^m.
\end{equation}
Hence $Df^m(y')$ is not an homothety. Since $y'$ is arbitrary in $B'$, one gets claim 1.
\end{proof}
\medskip

\begin{proof}[{\rm 8.} Proof of claim 2.]
As in the previous proof, let $z=f^m(y')$ and $\theta:=\angle(v'_m,e_{z})$. Then
$$\|Df^{-m}(v'_m)\|\simeq \theta \|Df^{-m}(e^\perp_{z})\|\geq \theta |\det Df^{-m}(z)|/\|Df^{-m}(e_{z})\|.$$
$$\text{With~\eqref{e.expansion} and~\eqref{e.choice-alpha},}\qquad \theta\leq |\det Df^{m}(y')|\cdot \|Df^{m}(v'_0)\|^{-2}\leq 4b^m (1-\rho)^{-2m}
<(1-\widehat \delta)^n.$$
We have shown in the proof of claim 1 that $\angle(v_m,v'_m)<(1-\widehat \delta)^n$.
One thus concludes that the fluctuations of $e_z$ are bounded by $4 (1-\widehat \delta)^n$.
\end{proof}
\bigskip

\noindent
9. {\bf Claim 3.} \emph{$\angle(u_0,e_c)<(1-\widehat \delta)^n$.}
\begin{proof}
From~\eqref{e.expansion} in the proof of claim 1, $\|Df^{-m}.e_c\|\leq 2(1-\rho)^{-m}$.
Equation~\eqref{e.g} gives:
$$|g^{-m}(u_0)|\leq |\det Df^{-m}(c)|/(\theta\|Df^{-m}.e_c^\perp\|)^2=\|Df^{-m}.e_c\|^2|\det Df^{m}(f^{-m}(c))|/\theta^2 .$$
Since $u_0$ is the critical direction at $c$, one has $|g^{-m}(u_0)|\geq 1$, and with~\eqref{e.choice-alpha},
\begin{equation*}\theta^2\leq 4(1-\rho)^{-2m}|\det Df^{m}(f^{-m}(c))|\leq 4(1-\rho)^{-2m}b^m\leq (1-\widehat \delta)^{2n}.\tag*{\qedhere}
\end{equation*}
\end{proof}
\bigskip

\noindent
10. {\bf Claim 4.} \emph{For any $x\in B$ with a unit vector $w$
satisfying $\|Df^m(G^{n-m}(w))\|\geq (1-\rho)^m$
$$\angle(G^n(w),e_{f^n(x)})<(1-\widehat \delta)^n .$$}
The proof is the same as for claim 2.
\bigskip

\noindent
11. For any unit vector $w$ at points $x\in B$, we will prove $\|Df^{2n}.w\|<(1-\eta)^{2n}$.
\bigskip

\noindent
12. If $\|Df^m.(G^{n-m}(w))\|\leq (1-\rho)^m$, then $\|Df^{2n}.w\|\leq (1-\rho)^m(1+\eta)^{2n}\leq (1-\eta)^{2n}$ from~\eqref{e.choice-eta}.
We are thus reduced to the case where $\|Df^n.w\|> (1-\rho)^n$.
\bigskip

\noindent
13. From \S3, and claims 2, 3, 4, $d(f^n(x),c)$ and $\angle(G^n(w),u_0)$
are smaller than $6(1-\widehat \delta)^n$.
Let $\ell:=[\beta n]$.
For any $1\leq k\leq \ell$, both $d(f^{n+k}(x),f^k(c))$ and
$\angle(G^{n+k}(w),G^k(u_0))$ are bounded by $6\Delta^k.(1-\widehat \delta)^n$.
The definition of $\ell$ and~\eqref{e.choice-beta} give:
$$\|Df^\ell(G^n(w))\|\leq \|Df^\ell(u_0)\|[\exp(6\ell\Delta^\ell.(1-\widehat \delta)^n)]\leq 
2\|Df^\ell(u_0)\|.$$
Note that by~\eqref{e.g} and since $|g^\ell(u_0)|\geq 1$,
we have $\|Df^\ell(u_0)\|\leq |\det Df^\ell(c)|^{1/2}$.
One deduces $\|Df^\ell.G^n(w)\|\leq 2 |\det Df^\ell(c)|^{1/2}\leq b^{\ell/2}$, then
$\|Df^{2n}.w\|\leq b^{\ell/2} (1+\rho)^{2n}$.
Our choice of $\ell$ and \eqref{e.choice-beta} then imply $\|Df^{2n}.w\|\leq (1-\eta)^{2n}$
as required.
\bigskip

\noindent
14. We have proved the inequality $\|Df^{2n}\|<(1-\eta)^{2n}$ on the ball $B=B(c,(1-\delta)^n)$.
Since $d(c,f^n(c))<(1-\varepsilon)^n$, we have $d(f^n(c),f^{2n}(c))<(1-\varepsilon)^n(1+\eta)^n$
and $d(c,f^{2n}(c))<(1-\varepsilon)^n[1+(1+\eta)^n]$. Since $\delta<\varepsilon$,
the ball $B$ is mapped into itself by $f^{2n}$, is contracted, and $c$ belongs to the basin of a sink.
The proposition is proved.

\end{proof}





%\begin{proof} If $Df^n$ is not an homothetie then it follows that either $G^{nk}(v)\to e^+$ as $k\to +\infty$ where $e^+$ is the eigendirection with the largest eigenvalue in absolute value or $v=e^-$, where $e^-$ is the eigendirection with the smallest eigenvalue in absolute value. The first case cannot hold since $g^{n}(e^+)< 1$ and $g^{nk}(v)\geq 1.$ The second case cannot hold either since it would follows that $G^n(v)=v$ and so $g^n(v)=1$ ($g^n(v)\geq 1$ and $g^{-n}(v)\geq 1$)  but $g^{n}(e^-)<1.$  A contradiction.
 
%\end{proof}



%However, in the context of dissipative dynamics, the last property listed in proposition \ref{properties} can be improved. More precisely, we now assume the following:
%\begin{description}
%\item[(H)] There exists $b\in(0,1)$ such that
%for any invariant ergodic measure $\mu$ on $\Lambda$ the difference between the two Lyapunov exponents is larger than $-\log b$.
%\end{description}
%This is the case for instance if the jacobian of $f$ on the set $\Lambda$ is smaller than $b$
%and if $\Lambda$ does not contain any sink.

%\begin{proposition}\label{crit exp}
%Let $\Lambda$ satisfies (H) and let $\lambda\in (1,b^{-1/2})$.
%For any $x\in\Lambda$, $v\in T_x^1M$ satisfying
%$\forall n\in \NN, \; g^n(v)\geq 1$,
%there exists infinitely many $n\geq 1$
%such that $g^n(v)> \lambda^n$.

%For any $x\in\Lambda$, $v\in T_x^1M$ satisfying
%$\forall n\in \NN, \; g^{-n}(v)\geq 1$,
%there exists infinitely many $n\geq 1$
%such that $g^{-n}(v)> \lambda^n$.

%\end{proposition}



\section{Critical points for  mildly dissipative diffeomorphisms of the disk. Proof of theorem \ref{disk2}}\label{sec:disk}



In what follows, we assume that $\DD$ is an open disk that is properly invariant under $f$, that is, $\overline{f(\DD)}\subset \DD$ and $f: \DD\to f(\DD)$ is a diffeomorphism and $f_{|\DD}$ is mildly dissipative.

In the dissipative context, given a fixed (or periodic) point $p$  there is  a well defined local center-unstable manifold denoted as $W^{cu}_\e(p)$, and   $W^{cu}_\e(p)\setminus\{p\}$ has two connected components.  The unstable set $W^u(p)$ of $p$ is the set of points whose backward orbit
converges to $p$. The local unstable set is defined as the set of points whose backward orbit
converges to $p$ and remains in a small neighborhood of $p;$  when this set is not reduced to $p$, it is a $C^1-$curve which
contains $p$ (not necessarily in the interior). In this case, $p$ is called a saddle point. At least one of the center-unstable branches at $p$ is contained in the unstable manifold of $
p$. The other branch, if not contained in the unstable manifold, is either a normally hyperbolic arc or contained in the basin of attraction of $p$; in the latter case, $p$ is called semi-attracting.  Similarly, as  in the case of the center-unstable manifold, each connected component $\Ga^u(p)$ of $W^u(p)\setminus\{p\}$ is called an unstable branch of  $p$. 

\vskip 5pt



The proof of theorem \ref{disk2} is decomposed into a few lemmas and claims but before getting into details, we provide an outline of the proof that may help the readers as a road map.


\noindent \emph{Outline of the proof of theorem \ref{disk2}:} Let $\La^u$ be the accumulation set of the
unstable branch $\Ga^u(p)$. The argument has three stages. \emph{First}
(the far-from-homotheties reduction),  by theorem \ref{dissipative-FFH}, after removing the basins of
all attracting periodic orbits $(q_i)$ of $\La$, every power of $f$ is far
from homotheties on what remains of $\La^u$. Theorem \ref{p.existence}
applied to this sink-free part gives the first dichotomy: either it already
contains a critical point, and the theorem is proved, or it admits a
dominated splitting that can be extended to a neighborhood $U$
of the sink-free part. In the latter case the splitting implies the
far-from-homotheties property of all powers, in the neighborhood $U$; and the full orbit of every sink \emph{except
finitely many} $\{q_1,\dots,q_k\}$ already lies in $U$. It remains only to
recover the property on the local basins of these finitely many exceptional
sinks, which is the role of claim \ref{cl:att-fp}: an exceptional $q_i$ with
eigenvalues of equal modulus is itself a critical point (and we are done),
while one with eigenvalues of distinct moduli is far from homotheties, for
every power, on a ball of definite radius around its orbit. Adjoining these
finitely many balls (and finitely many of their preimages) to $U$ covers
$\La^u$; hence every power of $f$ is far from homotheties on all of $\La^u$,
and a final application of theorem \ref{p.existence} yields the dichotomy of
the theorem: \emph{either} ${\rm Crit}(f,\La^u)\neq\emptyset$, \emph{or}
$\La^u$ admits a dominated splitting $E^s\oplus F$, which we assume
henceforth. \emph{Second} (the geometry of $\La^u$),
assuming there is no critical point we analyse how $\La^u$ sits with respect
to the center-unstable foliation: lemma \ref{pre-cu-omega} shows that at
least one center-unstable branch of each point lies in $\La^u$; lemma
\ref{cu-omega} upgrades this along $\omega$-limits; and lemma
\ref{whole-centerunstable} concludes that \emph{both} branches lie in
$\La^u$ except at finitely many boundary periodic points. \emph{Third} (the
trapping region), after a perturbation that turns those finitely many
boundary points into (semi-)attracting fixed points, the local stable arcs
over $\La^u$ assemble into an open, connected, simply connected set $B$ with
$\overline{f(B)}\subset B$ and $\bigcap_{n>0}f^n(B)=\La^u$. The
decomposition of remark \ref{decomposition} applies to $f_{|B}$, and a
Poincar\'e--Bendixson argument rules out homoclinic classes and normally
hyperbolic closed arcs unless a tangency and hence a critical point in
$\Ga^u(p)$ is produced; what survives is exactly the dichotomy of the
theorem. The critical points and normally hyperbolic arcs found for the
perturbed map are valid for $f$, since the two coincide on $\La^u$.

\medskip

Now, let us proceed to formalize that outline. Let $\La^u:= \overline{\cup_{n>0} f^n(\gamma^u)}$  be the accumulation set of the branch $\Ga^u,$ where $\gamma^u$ is a fundamental domain of $\Gamma^u$.   From the fact that $f$ is dissipative it follows that $\Lambda^u$ is a cellular set (a compact connected set of the plane  whose  complement is connected). The sufficiently dissipation condition implies that after removing all the attracting periodic points $(q_i)$ in $\La$ it holds that  for any positive integer $N$, $f^N_{|\La^u\cap [\cup_{i} W^s(q_i)]^c}$ 
is far from homotheties and so one can apply theorem \ref{p.existence}: if there is a critical point in $\La^u\cap [\cup_{i} W^s(q_i)]^c$ then the proof of theorem \ref{disk2} is over, if not, then $\La^u\cap [\cup_{i} W^s(q_i)]^c$ has a dominated splitting.    Observe that this splitting can be extended to a neighborhood $U$ of $\La^u\cap [\cup_{i} W^s(q_i)]^c$, in particular, for any positive integer $N$, $f^N$ is far from homotheties in $U.$  Observe also that except for   a finite number of attracting periodic points $\{q_1, \dots, q_k\}$ in $\La^u$, the orbits of the others are fully contained in $U$. Since the result is about $\La^u$ one has to analyze now what happens in the complement  of $U$ that intersects the basins of attraction. For that, we use the following claim (the proof is straightforward from the definitions and left to the reader, see remark \ref{rmk:why-powers}).

\begin{claim}\label{cl:att-fp} Let $q$ be an attracting fixed point and let $\la$ and $\si$ be the eigenvalues of $D_qf.$
\begin{itemize}
    \item[--] if $|\la|<|\si|$ then all the powers of $f$ are far from homotheties in the local basin of attraction of $q$;

     \item[--] if $|\la|=|\si|$ then  $q$ is a critical point.
\end{itemize}
    
\end{claim}

If one of the attracting periodic points $\{q_1, \dots, q_k\}$ (taking an iterate, we may  assume that they  are   fixed) has  eigenvalues of equal moduli, then $\La^u$  has a critical point. So, we have to consider the case that all the attracting  periodic points $\{q_1, \dots, q_k\}$  have eigenvalues of distinct moduli. In that case, for $r>0$ small enough it follows that for any positive integer $N$ we have $f^N$ restricted to $B=\cup_{i=1}^k B_{r}(O(q_i))$   is far from homotheties.  Fixing a positive integer $m$, we also see that $f^N$ restricted to $B_m=\cup_{i=1}^mf^{-i}(B)$ is also far from homotheties for any positive integer $N$.  Taking $m$ large, $B_m\cup U$ covers $\La^u$, therefore $f^N_{|\La^u}$ is far from homotheties for any positive integer $N$ and so one can apply theorem \ref{p.existence}: either there are critical points in $\La^u$ or $\La^u$ has a dominated splitting. So, either the proof of theorem \ref{disk2} is complete or we continue assuming that $\La^u$ has a dominated splitting.

\medskip

The following claim was proved in \cite{CPT} (see proposition 2.7).

\begin{claim}\label{cpt-lemma} Let $f$ be mildly dissipative and let $p,q$ be  two saddle points and $\Ga^u(p)$ and $\Ga^u(q)$ two of  their branches. If the accumulation set of $\Ga^u(p)$ intersects $\Ga^u(q)$ then it follows that the accumulation set of $\Ga^u(q)$ is contained in the accumulation set of  $\Ga^u(p).$ 
    
\end{claim}

From the previous claim it follows that if $q$ is a periodic point such that one branch  of the unstable manifold of $q$ intersects $\La^u$, then the whole branch is in $\La^u.$


If $\Lambda^u$ is a single point $q$, then $\La^u$  is a fixed point (if $\La^u$ is  finite, then it is a periodic point). The local center-unstable branch of $q$ that is accumulated by $\Ga^u$ can neither be contained in a normally hyperbolic arc nor be contained in the unstable manifold of $q$  otherwise, by claim \ref{cpt-lemma}, the accumulation set of $\Gamma^u$ should contain one of the unstable branches of the saddle $q$ or a normally hyperbolic invariant arc. Therefore, the fixed point is a (semi)attracting fixed point and its basin of attraction contains $\Gamma^u$. 

To conclude let us assume that $\La^u$ is not a single point and does not contain critical points. Therefore,  $f_{|\La^u}$ has dominated splitting and from the fact that $f$ is dissipative then it holds that  the splitting exhibits a contractive direction ($E^s\oplus F$). In particular, for any point in $\La^u$ there is a local stable manifold and two  center-unstable branches. The next lemma is the key lemma to prove theorem \ref{disk2}.



\begin{lemma}\label{whole-centerunstable} For any $x\in \La^u$ it holds that  there exists $\e(x)>0$ such that

\begin{itemize}
    \item[--]either  $W^{cu}_{\e(x)}(x)\subset \La^u,$ meaning that both (sufficiently small) branches of the center-unstable of $x$ are contained in $\La^u$, 

    \item[--] or only  one of the center-unstable branches  is in $\La^u$ and in that case $x$ is a periodic point, that we call boundary periodic point. 
    
    % only one branch $\Ga^{cu}_{\e(x)}(x)\subset \La^u$ and there is a periodic point $q$ such that $ x\in W^s(q).$
\end{itemize}

    Moreover, there are only a finite number of boundary periodic points.
    
\end{lemma}
The proof of the previous lemma goes through two lemmas; the first one  states that at least one center-unstable branch is in $\La^u$ (lemma \ref{pre-cu-omega} whose proof needs claim \ref{attracting}), the second  states a similar statement to lemma \ref{whole-centerunstable} but for points in the omega limit of another point in $\La^u$ (lemma \ref{cu-omega},  its proof also needs claim \ref{attracting}).



\begin{claim}\label{attracting} Let $z$ be a (semi)attracting periodic point and let $x$ be a point in $\La^u$ and in the local basin of attraction of $z.$ Then there is a normally hyperbolic arc contained in $\La^u$ such that one of the extreme point is $z.$ 
    
\end{claim}



\begin{proof} Let us assume that $z$ is fixed. Let $l$ be an arc contained in the unstable branch of $p$ (the fixed point from the hypothesis of theorem \ref{disk2}) with one endpoint is close to $x$ and the other is in the boundary of the local basin of attraction of $z$; that point exists otherwise $\La^u$ would be reduced to $z$. The forward iterates of $l$ is  a sequence of arcs $(l_n)$ that are in the unstable manifold of $p$ such that one of their extreme points is in the boundary of the local basin of $z$ and the other converges to $z$. Let $\hat l$ be a continuous arc that is limit of the arcs $(l_n)$ and observe that the interior of $\hat l$ is contained in the local basin of $z$, one  of the boundary points of $\hat l$ is $z$ and the other, named $y$, is in the boundary of the local basin of $z$.  Since $y\in \La^u$ let us consider the local stable manifold of $y$; let  $y'$ be a point sufficiently close to $y$ in $l$,   and observe that the local stable manifold of the  forward iterates of $y'$ converge to  $z$ and since its stable manifold separated the disk, it follows that the stable manifold of an iterate of $y'$ also separates the disk and so also the stable manifold of $y'$; given that $y'$ can be taken arbitrary close to $y$ it follows that also the stable of $y$ separates the disk. Let  $\ga_n$ the connected  part of the stable manifold of $f^{-n}(y)$ intersected with the disk  that contains $y$ and observe that $\ga_n$ separates the disk; observe that $\ga_1$ is in the connected component of $\DD\setminus\{\ga_0\}$ that does not contain $z$ and similarly $\ga_{n+1}$ is in the connected component of $\DD\setminus\{\ga_n\}$ that does not contain  $\ga_{n-1}$. In that way, we get a sequence $(\ga_n)$ of stable leaves ordered by the  inclusion  of the connected components and so there is a limit  $\hat \ga$ of the sequence of stable arcs $(\ga_n)$; since $f(\ga_n)\subset \ga_{n-1}$ it follows that $\hat \ga$ is a continuous invariant arc and so it has a fixed point $z'$ that is also the limit of $f^{-n}(y)$. Observe that the connected component $\hat \DD$ of $\DD$ bounded by $W^s_\DD(z)$ and $\hat \ga$ is a forward invariant disk foliated by local stable manifolds of points in   $\tilde l= \cup_{n>0}f^{-n}(\hat l).$  Since $\hat l$ is forward invariant, then $\tilde  l$  is also forward invariant so it is a normally hyperbolic arc. 


\end{proof}

\begin{lemma}\label{pre-cu-omega} For any $x\in \La^u$ there is $\e(x)>0$ such that one of the branches of $W^{cu}_{\e(x)}(x)$ is contained in $\La^u.$
    
\end{lemma}

\begin{proof}
If $x$ is a (semi)attracting periodic point, the lemma follows from claim \ref{attracting}. So, let us assume that  $x$ is not a (semi)attracting periodic point.
For any point $x\in \La^u$ and $\e$ small, let $\La^{u}_\e(x)$ be the connected component of $\La^u\cap B_\e(x)$ that contains $x$ and observe that  is a continuous set that contains $x$. If $\La^u_\e(x)$ is not in one of the branches of the  local center-unstable of $x$, by lemma \ref{centerunstable}, it holds that  backward iterates of $\La^u_\e(x)$ converge to the whole stable set of a point in $\La^u$, but that is not possible since by the assumption of mild dissipation, the stable manifold cannot be contained in the disk. 
\end{proof}



\begin{lemma}\label{cu-omega} Let $x\in \La^u$; for any $z\in \omega(x)$ it follows

\begin{itemize}
    \item[--] either  there exists $\e(z)>0$ such that $W^{cu}_{\e(z)}(z)\subset \La^u$, meaning that both (sufficiently small) branches of the center-unstable of $z$ are contained in $\La^u$;

    \item[--] or $z$ is a periodic point in $\La^u$ such that $x\in W^s(z)$ and at least one center-unstable branch of $z$ is contained in $\La^u.$
    
   %  Moreover

   %  \begin{itemize}
   %      \item either $z$ is an attracting periodic point and the forward iterates of $W^{cu}_{\e(x)}(x)$ are contained in the local basin of attraction of $z$,

   % \item or  $z$ a semi-attracting periodic point and the interior of the forward iterates of $W^{cu}_{\e(x)}(x)$ are contained in the interior of the local basin of attraction of $z$.

        

   %  \end{itemize}
    
    
    
    
\end{itemize}


    Moreover, there are only a finite number of periodic points that satisfy that only one center-unstable branch is in $\La^u$.
\end{lemma}

% Now, we consider $W^s_\de(\La^u)$ and we get the following claim:

% \begin{claim}\label{center-continuous} For any $x\in \La^u$ 
% we claim that there is $\e(x)$ such that 
% \begin{itemize}
%     \item[--] either $W^{cu}_{\e(x)}(x)$ is contained in $W^{cu}_\e(x)\cap W^s_\de(\La^u),$

%     \item[--] or there are finite number of periodic points such that $x$ is in the stable manifold of one of these periodic point and $W^{cu}_\e(x)\cap W^s_\de(\La^u)$ coincides with  only one center-unstable branch of $x.$
% \end{itemize}
    
% \end{claim}



\begin{proof}
    

By  lemma \ref{pre-cu-omega}, given $x$ there is one   branch $J$ of $W^{cu}_\e(x)$  contained in $\La^u$. Either $length(f^n(J))\to 0$ or there exists a sequences of integer $(n_k)$ such that $length(f^{n_k}(J))$ is uniformly bounded by below. 

In the first case, by lemma \ref{centerunstable2} it follows that $(f^n(J))$ converges to a (semi)attracting periodic point $z$. If the point is attracting, then $f^n(J)$ is in the interior of the basin of attraction of $z$; if $z$ is semi-attracting, then  the interior of  $f^n(J)$ is in the interior of the local  basin of attraction of $z$ and $x \in W^s(z)$ (meaning that the iterates of $x$ in in the boundary of the basin of attraction of $z$). In both cases, one can apply claim \ref{attracting} and so one center-unstable branch of $z$ is contained in $\La^u.$


% By lemma \ref{centerunstable2} it follows that either their length of the forward iterates go to infinity or the lemma is proved.


% . In the first case, $x$ either the length of the forward iterates of the center-unstable branch goes to zero, or there is a subsequences of iterates that their length is bounded by below. In the case that the length go to zero, one get a  (semi)attracting periodic point and the whole arc is in the basin of attraction of that point.

 

In the second case, one can assume also that there is $z$  such that $f^{n_k}(x)\to z.$ Since $length(f^{n_k}(J))$ are uniformly bounded by below it follows that here is an arc $I$ containing $z$ and contained  in one center-unstable branch of $z$ that  is also  contained in the limit of $f^{n_k}(J)$ and so $I\subset \La^u.$ 

If there are two integers $n_k< n_{k'}$ such that  $f^{n_{k'}}(x)\in W^s_\de(f^{n_k}(x))$ it follows the there is a periodic point $z'$ in the local stable manifold of $f^{n_k}(x)$ and so $z=z'$  and in that case    $x$ is in the stable manifold of it and by lemma \ref{pre-cu-omega} one of its center-unstable branches is in $\La^u$. In the other case, let $I'$ be an arc containing $z$ in the other center-unstable branch of $z$ disjoints with $I.$ Either 

\begin{itemize}
    \item[--] i) there is an  iterate  $f^{n_k}(x)$ close to $z$  in the interior of $W^s_\de(I)$,

    \item[--] ii) or there are infinitely many iterates  $f^{n_k}(x)$ accumulating on $z$ in the interior of $W^s_\de(I').$
\end{itemize}

  In the first case, let $\hat J$ be an arc in the center-unstable of $f^{n_{k_0}}(x)$ that one extreme point is given by that point and the other by the intersection of the local stable of $z$ and the center-unstable of $f^{n_{k_0}}(x).$ It follows immediately that $\hat J$ is disjoint with $f^{n_{k_0}}(J)$ and observe that since $\hat J$ is in the local basin of $I$ then its forward iterates remain in a neighborhood of $\La^u$ and so  we can apply again lemma \ref{centerunstable2}; if the length of the iterates of $\hat J$ go to infinity, then the accumulation of the  forward iterates of $f^{n_k}(\hat J)$ are in $\La^u$ and  contain the other center-unstable branch $I'$ of $z$ disjoint with $I$ and concluding that the whole center-unstable of $z$ is in $\La^u.$ If the lengths of the iterates of of $\hat J$ does not go to infinite,  $x$ is in the basin of a normally hyperbolic arc and $z$ is a periodic point such that one of its center-unstable branch is contained in $\La^u.$  
  
  In the second case, i.e. that $f^{n_k}(x)$ is in the local basin of $\hat I$, then we can take $n_{k'}$ such that $f^{n_{k'}}(x)$ is closer to $z$ and again we can consider an arc $\hat J$ is the center-unstable branch of $f^{n_{k'}}(x)$ that one extreme point is given by that point and the other by the intersection of the local stable of $f^{n_k}(x)$ and the center-unstable of $f^{n_{k'}}(x).$ It follows immediately that $\hat J$ is disjoint with $f^{n_{k'}}(J)$ and observe that since $\hat J$ is in the local stable manifold of  $f^{n_k}(J)$, that is contained in $\La^u,$ then its forward iterates remain in a neighborhood of $\La^u$ and so  we can apply again lemma \ref{centerunstable2} and argue as before.





% and in that case both center-unstable branches of $x$ intersect $W^s_\de(I)$ or for iterates close to $y$ it holds that  $W^s_\de(f^{n_k}(J))$ intersects both center-unstable branches of a point $f^{n_{k'}}(x)$ much closer to $y$.





% the forward iterates of $J.$ Either the length of all the forward iterates remain bounded by above or there are sub-arcs $J_k\subset J$ and positive integers $n_k$  such that $x\in J_k,$  $f^{i}(J_k)$ are in a small neighborhood of $\La^u$ for any $i\leq n_k$ and  $length(f^{n_k}(J))$ are uniformly bounded by below. In the first case, by lemma \ref{bounded-arcs}, it follows that the forwards iterates of $J$ converge to a normally hyperbolic arc and $x$ is in the stable manifold of one of the boundary points of the periodic arc. In the second case, let us consider an arc $I$ that is limit of the iterates $f^{n_k}(J_k)$;  it follows that $I$ is contained in $\La^u$, it has positive length and  one of the extreme points is the limit of $f^{n_k}(x)$. If $f^{n_k}(x)$ are not in the stable manifold of $y$, it follows that there are iterates  $f^{n_k}(x)$ close to $y$ in the interior of $W^s_\de(I)$ and in that case both center-unstable branches of $x$ intersect $W^s_\de(I)$ or for iterates close to $y$ it holds that  $W^s_\de(f^{n_k}(J))$ intersects both center-unstable branches of a point $f^{n_{k'}}(x)$ much closer to $y$.  Since we are assuming that $W^s_\de(\La^u)$ intersects only one center-unstable branch of $x$  (and therefore also of any $f^n(x)$), it follows that $f^{n_k}(x)\in W^s_\de(y)$ and so $f^{n_{k'}}\in W^s_\de(f^{n_k}(x))$ for $n_{k'}> n_k$ which implies that there is a periodic point in that local stable manifold and so $x$ is in the stable manifold of that periodic point.


To prove the last claim of the lemma (about the finiteness of periodic points satisfying that  only one center-unstable branch is contained in $\La^u$),  first we prove that their periods are uniformly bounded. That is immediate for  periodic points that are either   (semi)attracting or contained in a normally hyperbolic arc. For the remaining periodic points, by proposition \ref{PS2-1}   any  pair of periodic points with large period and close enough are homoclinically related, but this implies that given three periodic points sufficiently close, it follows that the local stable manifold of  one  of them, named $q$, intersects the unstable branch of another that is contained in $\La^u$ and so by the classical inclination lemma it follows that the unstable of $q$ is contained in $\La^u$, a contradiction.  
So, if there are infinitely many periodic points, given that their periods are uniformly bounded, one can regroup them into a finite number of normally periodic arcs that are contained in $\La^u.$
\end{proof}


    
% \end{lemma}

The next remark is a simple observation  that follows immediately from the dynamics of fixed points in the dissipative context. 

\begin{remark}\label{att-NH} Let $z$ be a periodic point  and let us choose one of the center-unstable branches. Either that branch is

\begin{itemize}
    \item[--]in the basin of attraction of $z$ or,

    \item[--] a normally hyperbolic arc or, 

    \item[--] contained in the unstable manifold of $z.$
\end{itemize} 
Moreover, given a sequence of points accumulating on $z$ and inside the connected component of $B_\e(z)\setminus W^s_\e(z)$ that intersects the chosen center-unstable branch, then either 


\begin{itemize}
    \item[--]the forward iterates of that sequence converge to $z$  or to points in a  normally hyperbolic arc,

    \item[--] or there is a sequence of forward iterates that accumulates in a fundamental domain of the unstable manifold of $z.$
\end{itemize} 
 
    
\end{remark}

% Let us assume that only one center-unstable branch of $x$ is in $\La^u.$ 

\begin{proof}[Proof of lemma \ref{whole-centerunstable}]


Let $y\in \alpha(x)$ and let $z\in \omega(y)$; on one hand it follows immediately that $z\in \alpha(x)$, on the other, we can apply lemma \ref{cu-omega} to the point $z$. 

Let us consider first the case that  both center  unstable of $z$  are in $\La^u$.   Let $(n_k)$ be a sequence of positive integers  such that $f^{-n_k}(x)\to z$. It could hold that  
\begin{itemize}
\item[--] either there are iterates $(f^{-n_{k}}(x))$ accumulating on $z$  that are contained in the  local center-unstable of $z,$


\item[--] or there are two iterates $f^{-n_k}(x), f^{-n_{k'}}(x)$ that belong to the local stable of $z,$


\item[--] or there are infinitely many iterates $f^{-n_k}(x)$ that are in the interior of a  quadrant $Q$ determined  by the local stable and center-unstable of  $z.$

\end{itemize}

In the first case since both branches of $z$ are in $\La^u$ it holds by continuity of the center-unstable manifold along the orbit the same holds for $f^{-n_k}(x)$ and so for $x$. In the second case, $f^{n_{k'}-n_k}(x)$ is in the local stable manifold of $x$ and so $x$ is a periodic point and therefore $x=z$ and again both center-unstable  branches of $x$ are in $\La^u.$ In the third case, let  $J_k^{+}$ be one arc in the center-unstable of $f^{-n_k}(x)$ and $J_k^{-}$  in the other, and let us treat them separately. For each $J^+_k$ (the arcs contained in the quadrant $Q$), let $\hat J^+_k$ be the arc in the center-unstable given by the intersection of the local stable manifold of $J^+_k$ with the center-unstable of $z$ and observe that this arc in $\La^u$. If it holds that there is a subsequence of $(f^{n_{k'}}(J^+_{k'}))$ with length uniformly bounded by below, since $f^{n_{k'}}(\hat J^+)\in W^s_{\la^{n_{k'}}.\de}(f^{n_{k'}}(J_{k'}^+))$ it follows also that the length of the subsequence $(f^{n_{k'}}(\hat J^+_{k'}))$ is also bounded by below, therefore  there is an arc $\hat J$ containing $x$ and contained  in one of the branches of the center-unstable of $x$ (observe that $f^{n_{k'}}(J^+_{k'})$ contains $x$) that is in $\La^u.$ If the lengths of $f^{n_k}(J^+_k)$ goes to zero, for each $k$ one takes $k'$ larger than $k$ such that $f^{-n_{k'}}(x)$ is closer to $z$ than $f^{-n_k}(x)$ and since are accumulating in the same quadrant let $x_{k'}$ in the intersection of the local stable of $f^{n_k}(x)$ and $J^+_{n_{k'}};$ from the fact that $f^{-n_k}(x)$  is in the local stable manifold of $x_{k'}$ and the length of $f^{n_{k'}}(J_{k'})$ goes to zero follows that that $f^{n_{k'}}(x_{k'})$ is arbitrary close to $x$ and $f^{n_{k'}-n_k}(x)$ provided that $k$ and $k'$ are large enough; therefore $x\in \omega(x)$ and so by lemma \ref{cu-omega} it follows that both center-unstable branches of $x$ is in $\La^u$ or it is a periodic point with at least  one center-unstable branch in $\La^u$. A similar argument can be made for the other branch $J^-_k$ with a caveat; if the length of the iterates of $f^{n_k}(J^-_k)$ go to zero, one considers  the arcs  $\hat J^-_k$  in the center-unstable of $z$ given by the intersection of the local stable manifold of $J^-_k$ with the center-unstable of $z$ and observe that these arcs are in $\La^u$ and each of them contain one local center-unstable branch of $z$, named $J_z$; therefore, if the length of the iterates of $f^{n_k}(J^-_k)$ go to zero, also the length of $f^{n_k}(J_z)$ and by lemma \ref{centerunstable2} follows that $z$ is in the basin of attraction of a (semi)attracting periodic point or it is in the basin of a normally hyperbolic arc and so converging to a periodic point; since $z$ is  in $\alpha(x)$ it follows that $x$ is  a periodic point.





% one takes $x_{k'}$ in the intersection of the stable manifold of $f^{-n_{k'}}(x)$ with $J^-_k$ and again taken $k'$ and $k$ large enough since $f^{n_{k'}-n_k}(x_{k'})$ is close to $x$ (because $x_{k'}$ is in the local stable manifold of $f^{-n_{k'}}(x)$) and also close to $f^{n_{k'}-n_k}(x)$ (because  



% Let  $J$ be an arc in  the center-unstable of $z$ contained in $\La^u$ and  containing $z$  in its interior and let $J^+$ and $J^-$ be the closure of the two connected components of $J\setminus\{z\}.$   and let  $J_k^{\pm}$ be  the arc in the  center-unstable of $f^{-n_k}(x)$ obtained as the intersection of the center-unstable of $f^{-n_k}(x)$ with $W^s(J^\pm).$ 

% % $length(f^{n_k}(J_k^\pm))\to \infty$

% Now, we can apply lemma \ref{centerunstable2} to $J^+$ and $J^-$. If both branches satisfy that length of their forward iterates goes to infinity,   since $f^{n_k}(J^\pm)\in W^s_{\la^{n_k}\de}(f^{n_k}(J_k^\pm))$ then   the lengths of the $n_k-$iterate of both branches of $J_k$ also go to infinity. Since each $f^{n_k}(J_k^\pm)$ contains a center-unstable branch of $x$ each one at each side of $x$ that are accumulated by $f^{n_k}(J^\pm)$, it follows that both center-unstable branches of $x$ are in $\La^u.$ 

% If $length(f^n(J^\pm))$ does not converge to infinity for one of the branches  then forward iterates of one of the branches $J^\pm$  converge either to a (semi)attracting periodic point or to a normally hyperbolic arc. By claim \ref{attracting}, in both cases $x$ is a periodic point and  belongs to a normally hyperbolic arc.  

Let us consider the second option of lemma  \ref{cu-omega}, i.e., $z$ is a periodic point. If $z$ is an attracting periodic point, and $z\in \alpha(x)$ then $x=z$ and so $x$ is an attracting periodic point in $\La^u$ and by claim \ref{attracting} it follows also  that $x$ belongs to a normally hyperbolic arc in $\La^u$. If $z$ is a semi-attracting periodic point then  the pre-iterates of $x$ that accumulate on $z$ are either in the basin of attraction of $z$ or they belong to the connected component of $B_\e(z)\setminus W^s_\e(z)$ that is not in the basin of attraction of $z.$ In the first case, it follows that $x=z$ and again by claim \ref{attracting} it follows also  that $x$ belongs to a normally hyperbolic arc in $\La^u$. In the second case, by remark \ref{att-NH} it follows that either $x=z$, or  it belongs to an invariant normally hyperbolic arc or there exists $z'\in W^u_{\e}(z)$ accumulated by pre-iterates of $x$ and so $z'$ is in $\La^u$ and by claim \ref{cpt-lemma} it follows that the unstable branch of $W^u(z)$ that contains $z'$ is in $\La^u$ and so there are pre-iterates of $x$ that accumulates at a point such that both center-unstable branches are in $\La^u;$ therefore we can argue as done before.

% To finish, we have to consider the last option of claim \ref{cu-omega},  that is,  $z$ is a periodic point such that only one center-unstable branch of $z$ is in $\La^u$. Again, by remark \ref{cu-omega}, we can argue as in the previous  paragraph.


\end{proof}





\begin{proof}[Proof of theorem \ref{disk2}]






Given a point $x\in \La^u$, let $\hat W^{cu}_{\e(x)}(x)$ be the piece of the center-unstable manifold of $x$ that is contained in $\La^u.$ From lemma \ref{whole-centerunstable} it follows that the whole center-unstable $W^{cu}_{\e(x)}(x)$ is contained in $\La^u$ or $x$ is a boundary periodic. Now, we consider the box $B(x)$ defined as follows  $$B(x):=\cup_{y\in W^{cu}_{\eps(x)}(x)}W^s_\de(y).$$ 


It follows that either $x$ is in the interior of $B(x)$, or $x$ is a boundary
periodic point lying in the boundary of $B(x)$. In the latter case we modify
the diffeomorphism as follows. Recall (lemma \ref{whole-centerunstable})
that there are only finitely many boundary periodic points; taking an
iterate of $f$, we assume that they are fixed. For each boundary fixed point
$x$, let $C_x$ be the connected component of $B_\e(x)\setminus W^s_\e(x)$
that does not intersect the center-unstable branch of $x$ contained in
$\La^u$; in particular $C_x\cap \La^u=\emptyset$, and the components $C_x$
can be chosen pairwise disjoint. We choose a $C^2$ dissipative
diffeomorphism $g$ of the disk such that:
\begin{itemize}
\item[--] $g=f$ on the complement of $\bigcup_x C_x$;
\item[--] for each boundary fixed point $x$, the center-unstable branch of
$x$ contained in $C_x$ lies, for $g$, in the basin of attraction of $x$;
that is, $x$ is a (semi-)attracting fixed point of $g$.
\end{itemize}

The following three observations justify replacing $f$ by $g$ in the rest of
the argument.

First, the perturbation does not affect $\La^u$, neither topologically nor
at the level of the derivative: $g=f$ on $\La^u$, so $\La^u$ is
$g$-invariant; and since $g$ and $f$ coincide on the closure of the
complement of $\bigcup_x C_x$, which accumulates on every point of $\La^u$,
the continuity of the derivatives gives $D_zg=D_zf$ for every $z\in\La^u$.
Hence the derivative cocycles of $f$ and $g$ along $\La^u$ coincide; in
particular $\La^u$ carries the same dominated splitting for both maps. Moreover, since
$\Ga^u(p)\cup\{p\}\subset \La^u$ is disjoint from $\bigcup_x C_x$, the
unstable branch of $p$ is the same for $f$ and for $g$, and $\La^u$ is the
accumulation set of that branch for both maps.

Second, $g$ is again mildly dissipative: an ergodic $g$-invariant measure
either gives no weight to $\bigcup_x C_x$, in which case it is
$f$-invariant with the same stable branches at almost every point, or
charges some $C_x$; in the latter case, since every point of $C_x$ is
attracted under $g$ to the fixed point $x$, the measure is the Dirac measure
at $x$, and the stable branches of $x$ for $g$ contain its stable branches
for $f$, which reach the boundary of the disk by the mild dissipation of
$f$.

Third, by construction each boundary fixed point $x$ is a (semi-)attracting
fixed point of $g$, so that the extended box $B(x)$ obtained by adding
to $B(x)$ the component $C_x$  which is attracted to $x$, contains $x$ in
its interior.


% It follows that either $x$ is in the interior of $B(x)$ or $x$ is a boundary periodic point in the boundary of $B(x)$; in that case, consider the connected component of $B_\e(x)\setminus W^s_\e(x)$ that does not intersect the center-unstable branch contained in $\La^u;$ since that component is disjoint with $\La^u$, one can deform the dynamics in that component without affecting the dynamic in $\La^u$ and  assuming that $x$ becomes a (semi)attracting periodic point. In both situations, the box $B(x)$ can be extended to a new box that we keep noting as $ B(x)$ by adding that connected component that is attracted to $x$ and therefore again $x$ is in the interior of the extended box. Since there is only finite number  of boundary periodic points (lemma \ref{whole-centerunstable}), we can assume that they are also fixed.

After this modification (we keep calling $g$ as $f$), we can cover $\La^u$ by $$B=\cup_{z\in \La^u} \overset{\circ}{B}(z).$$

Observe that $B$ satisfies: 



\begin{itemize}
     \item[--] it is open and connected, 
\item[--] $\overline{f(B)}\subset B,$
 \item[--] it is simply connected.
\end{itemize}



 
 
The  first item follows from the fact that $B$  is an open union of neighborhood  of all the points of a connected set. To prove the second point, observe that 
 given $y\in B$ there is $x\in \La^u$ such that $y\in B(x)$, so  there is $x'$ in the interior of $ W^{cu}_{\e(x)}(x)\cap \La^u$ such that $y\in W^s_{\e}(x')$ and so $f(y)\in W^s_{\la.\e}(f(x'))\subset W^s_\e(f(x'))$ and so $f(y)\in B(f(x'))\subset B.$ If $B$ is not simply connected, there are connected components of $\DD\setminus B$ that do not intersect the boundary of $\DD$; let $U$ be the union of these components and since $\overline{f(B)}\subset B$ it follows that $f(U)$ contains $U$, contradicting  the fact that $f_{|\DD}$ is dissipative.
 
 % or $x$ is a  boundary periodic point and so $f(y)\in B(x).$ 
 
 % It also holds that the intersection of $\La^u$ with the boundary of $B$ is given by a finite number of arcs contained in the stable manifold of a finite number of fixed points. 

Observe that $\cap_{n>0} f^n(B)=\La^u$ and so the limit set of $f_{|B}$ has dominated splitting. Now we can apply the remark \ref{decomposition}. 
If there is a non-trivial homoclinic class in the limit set of $f$ inside $B$,  then there is a loop $\ga$ formed by an unstable arc $\ga^u$ and a stable arc $\ga^s$ contained in $B$. On the other hand, observe that the stable foliation in $B$ is transversal to $\ga^u$ and $\ga^s$ is a stable leaf of that foliation; however, there is a point in $\ga^u$ tangent to the stable foliation; a contradiction: a closed curve transversal to a foliation on a disk would imply the loop bounds a disk filled by leaves (by Poincar\'e–Bendixson), which would force a leaf to terminate inside the disk; contradiction with the leaves being stable manifolds of points with no terminating dynamics. So, the limit set is just formed by a  finite number of periodic points and finite normally hyperbolic  periodic arcs. Taking an iterate of $f$, one can assume that the periodic points and arcs are fixed. Let us first conclude that there are not normally hyperbolic arcs. If the $\Ga^u(p)$ is not a normally hyperbolic arc (otherwise the theorem is proved) but  accumulates on a non-trivial normally hyperbolic invariant arc $J$, $\Ga^u(p)$ has to intersect the stable local manifold of of one of the  endpoints of $J$, that is a fixed point $q$; this implies that there is a closed arc $\ga=\ga^u\cup \ga^s$ formed by an arc  $\ga^u\subset \Ga^u(p) $ and an arc $\ga^s\subset W^s_\DD(q)$ that is contained in $B,$ and arguing as before one gets a tangency between $\ga^u$ and the stable foliation in $B$ and therefore a critical point in $\Ga^u(p)$. Now we  conclude that there is only one fixed point $q$ in $\La^u$ and it is a (semi-)attracting fixed point: if in $\La^u$ there is a saddle fixed point $q$ such that one of its unstable branches is in $\La^u$, then $\Ga^u(p)$ intersects the local stable of $q$ (if it accumulates without intersecting one gets a point in the local stable manifold that it is in the limit set of $\La^u$ but is not a fixed point); if $\Ga^u(p)$ intersects $W^s_\DD(q)$ without crossing, one gets a tangency between $\Ga^u(p)$ and $W^s_\DD(q)$ and so a critical point in $\Ga^u(p);$  if $\Ga^u(p)$ crosses  $W^s_\DD(q)$, arguing as before we get a closed arc in $B$ and again a tangency between $\ga^u$ and the stable foliation in $B$. So, there are only (semi)attracting fixed points. Let $z$ be one of these (semi)attracting fixed points and let $x$ be in $\La^u$ and in the basin of attraction of $z;$ by claim \ref{attracting}, there is a normally hyperbolic arc $\Ga$ in $\La^u$ that one of its extreme points is  $z$  and arguing as before, $\Ga\subset \Ga^u(p)$ and if it does not coincide with $\Ga^u(p)$ then one gets a  critical point in $\Ga^u(p).$



Finally, the conclusions obtained for the modification of $f$ transfer to $f$. The critical
points produced by the tangency arguments above lie in $\La^u$ (they are
obtained either on $\Ga^u(p)$ or on its accumulation set), and along
$\La^u$ the derivative cocycles of $f$ and its modification coincide, so they are
critical points of $f$ as well. Likewise, a normally hyperbolic invariant
arc contained in $\La^u$ for the modification of $f$  is normally hyperbolic for $f$, since
normal hyperbolicity only involves the restriction of the derivative cocycle
to the arc.

\end{proof}



% Two parallel boundaries of these rectangles are given by local stable manifolds (call it the stable boundary),  and by the mild dissipation condition one can assume that the other two boundaries (call it center boundaries) are given by two arcs that are disjoint with $\La^u.$ These rectangles either contain $x$ in its or one of its boundary is given by the local stable manifold  of $x$ and observe that in that case $x$ is also in the stable manifold of some periodic points that satisfies the second item of claim \ref{center-continuous}. Given that $\La^u$ is compact  one can cover it by a finite union $B:=\cup_{in \mathcal{I}} B(x_i)$ of these rectangles in such a way that any point in $\La^u$ is either in the interior or it is in the stable boundary of one of them. 


% it holds that any periodic point of sufficiently  large period is a hyperbolic saddle periodic point and $W^{cu}_\e(p)=W^u_\e(p)$; therefore, for any periodic with large period and satisfies the second item of the claim it holds that $\La^u(p)\cap W^u_\e(p)$ is an arc that contains $p$ and it is contained in one of the branches; by the fact that $\La^u$ is forward invariant and $p$ is periodic it follows that $\La^u(p)\cap W^u_\e(p)$ coincides with one of the branches of $\W^u_\e(p).$ Again, by theorem \ref{PS2-1} 
% In the first case, from the mild dissipative property it follows that the center-unstable manifold is contained inside $\La^u$ which is a celular set. In other words, the celullar set is saturated by the center-unstable manifold. 
% Now, for each point in $\La^u$, that is a celullar set, one consider the local stable manifold at that point, and consider the union of all the local stable manifold and that provides a neighborhood of $\La^u$ that is mapped properly inside.

% Moreover, for each point of $\La^u$ it follows that its local stable manifold intersects  in a unique point, otherwise, one can form a loop using a an arc contained in $\La^u$ and a piece of local stable manifold: in that case, either one gets a normally hyperbolic circle (which is not possible in the dissipative case) or one gets a critical point.   



% \begin{claim} Let $f\in Diff^{2}(M)$ and let $\La$ be a compact invariant set of $f$ with dominated splitting. If for any point, the center-unstable manifold is contained in $\La$ then it follows that $\La$ is a maximal invariant attractor.

% \end{claim}

% \begin{proof}[Proof of theorem \ref{disk}] Let $\La=\cap_{n>0}f^{n}(\DD)$ and let us   assume that $\La$ does not have critical points; so, it has a dominated splitting. Now, one can apply the spectral decomposition theorem in \cite{PS2} and since the dynamic is dissipative it cannot contain normally hyperbolic invariant closed curves and therefore   the Limit set is decomposed into a finite number of homoclinic classes and normally hyperbolic arcs. From the fact that there are not semiattracting periodic points, it follows that the limit set is given by only one homoclinic class that also is locally maximal and therefore, it coincides with the attractor $\La.$  Moreover, one can  prove that $f$ restricted to the class is conjugate to a hyperbolic set and this is a contradiction since there are not non-trivial hyperbolic attractors in the disk. 

% To conclude the second part, let $\Ga$ be an unstable branch of a periodic point $q$ and taking an iterate of the map one can assume that $\Ga$ is fixed. Let $\De$ be the accumulation set of $\Ga$ and let us assume that it does not have critical points and so it has dominated splitting. Let $U$ be a neighborhood on  $\De$ can be extended to $U$ and observe that for large $N$ it follows that $\Ga_N:=\cup_{n\geq N} f^n( W^u_{loc}(q))\cap \Ga$ is contained in $U.$ 

% \begin{claim} There exists $\eps>0$ such that for any $\de>0$ there a positive integer $N$  such that if $z\in \Ga_N$ the arc $\Ga_\eps(z)$ defined as the connected component of $\Ga_N\cap B_\eps(z)$ that contains $z$ is $\de$ $C^1-$close to $W^{cu}(z)$.

% \end{claim}

% As a consequences of previous claim it follows that the center-unstable manifold for any point in $\De$ is contained in $\De$ and this implies that $\De$ is an attractor: Let $z$ be a point nearby $\De$ and let $z'$ be a point in $\De$ close to $z;$ it follows that $z\in W^s_{\e}(z'')$ for some $z''$ contained in $W^{cu}(z')$ and since the forward iterates of $z$ gets closer to $z''$ and the orbit of $z''$ is in $\De$ it is concluded that the forward iterates of $z$ are in $U.$ 

% Now, using that $f$ is dissipative, it follows that  $U$ is a disk and as in the first part of the proof of the theorem, one gets a contradiction.


% \end{proof}


% \section{ Misiurewicz}


% Let $x\in \Omega$ and $\La$ be a basic piece.  It is said that $x$ is transitive  related to $\La$ if for any neighborhood  $U$ of $x$ it holds that $Acc^+(U)$ contains the local unstable manifold of $\La$   and $Acc^-(U)$ contains the local stable manifold of $\La$ (observe that this definition is invariant by iteration).  Let us called the closure of the  set of points that are transitive related to $\La$ as the extended basic piece and  denote that set as $\hat \La$. It follows immediately that $\hat \La$ is closed and transitive. 

% It is said that $x$ is forward related to $\La$ if for any neighborhood  $U$ of $x$ it holds that $Acc^+(U)$ contains the local unstable manifold of $\La$   and it is said that $x$ is backward related to $\La$ if for any neighborhood  $U$ of $x$ it holds that $Acc^-(U)$ contains the local stable manifold of $\La$ (observe that this definition is invariant by iteration).



% \begin{claim}If there is a finite number of basic pieces such that any point in the nonwandering is belong to one of the extended basic pieces then it follows that the nonwandering is decomposed ina finite number of transitive pieces.
 
% \end{claim}





% Observe that each $\La_c$ is a set with dominated splitting that is conjugate to a basic hyperbolic set. Moreover, the number of disjoint sets is finite.


% \begin{claim} Let $U$ be an open set such that $f(U)\subset U$ then it follows that $U\cap Crit(f)\neq \emptyset.$
 
% \end{claim}


% \begin{proof} Let us consider $\Ga=\cap_{n\geq 0} f^n(U)$ and observe that it is an invariant set. Let us assume that it does not intersect the critical set. Therefore, it can be decomposed in finite basic pieces. Moreover, since $f(U)\subset U$ at least one unstable branch of these pieces is contained in $U$ (it could hold that a point in $\Ga$ is in the boundary of $U$) and therefore in $\Gamma$. That implies one of the pieces is a (semi-)attracting basic piece in the disk, but this implies that there is a semi-attracting periodico point. A contradiction.
 
 
% \end{proof}


% \begin{claim} One can assume that for any $c$ it holds that $\omega(c)\cap interior(\La_c).$ 
 
 
% \end{claim}

% Maybe previous claim needs to assume mild dissipation or Pixton genericity.

% \begin{proof} By assumption, each critical point accumulate on a basic piece. The unstable branches of these pieces, aslo accumulates on critical points. Since there are only finite basic pieces  involved, then there are cycles involving these pieces and the critical points.  

% Let us assume that a critical point does not accumulate on the iterior of a basic piece. Then, it accumulates on a periodic point that it is either isolated or is a boundary point of a basic piece. There are two options, or the critical point is contained in the stable manifold of the periodic point or not.  If that critical point is contained in that cycle that at least one critical point accumulae on the interior of the class, one can carry the intersection of stable and unstable arc to extend the basic pieces to larger ones in such a way that the critical accumulates in the interior. If not, there are just cycles involving critical points accumulating in boundary points. Arguing as in \cite{CPT} or using Pixton generecity we get the result.
 
 
 
 
% \end{proof}


% Given a critical point $c$ and the associated set $\La_c$, there exists and open neighborhood $U$ of $Crit(f)$   such that $ \La_c\subset U^c.$ 
% As a consequence of previous claim it follows that for any critical point $c$   there is a rectangle $R_c$ (arbitrary small) with boundaries given by stable and unstable local arcs of $\La_c$  such that an iterate of $c$ is in the interior of $R_c.$ In particular, for each $c$ there is a open neighborhood $U_c$, a rectangle $R_c$ related to $\La_c$  and a positive integer  $n_c$ such that $f^{n_c}(U_c)\subset R_c$. By compactness, there is a finite number of open sets $(U_i)$, finite number of basic sets $\La_i$, positive integers $(n_i)$ and rectangles $(R_i)$ related to $\La_i$ such that $Crit(f)\subset \cup_i U_i$ and $f^{n_i}(U_i)\subset R_i$. Moreover, one can assume that $ \cup_i U_i\subset U.$


% \begin{claim}\label{insideUi}
 
% If $x\in \Omega$ and $ x\in \cup U_i$ then $x$ is transitive related to  $\La_i$.
% \end{claim}

% \begin{proof} Let $I$ be an open set that contains $x$ and let $n$ be a positive integer such that $f^n(I)\cap I\neq \emp.$ Let $J=\cup_{n\geq 0} f^{n.m}(I)$ and  observe that $J$ is connected and $f^n(J)\subset J$; so, $L=\cup_{i=0}^{n-1}f^i(J)$ verifies $f(L)\subset L.$ Therefore, it follows that $\bar L\cap Crit(f)\neq \emp$ and so there is $c\in Crit(f)$ and $j>0$ such that $f^j(c)\in \bar J.$  In particular, for any $l>0$ it follows that $f^{n.l+j}(c))\in \bar J$. From the fact that the there is a subsequences $(n_k)$ such that $f^{n_k}(c)\to interior \La_c$ then there exists a sequences of positive integers $(l_i)$ such that $f^{n.l_i+j}(c)\to interior(\La_c)$ and so $\bar J \cap interior(\La_c).$ In particular, for  an arbitrary large $l$ such that $f^{n.l}(f^j(c))$ is inside a small rectangle $R_c$. Since $x\in U\cap J$, $f^{n.l}(f^j(c))\in \bar J$ and $\bar J$ is connected, then it follows that $\bar J$ intersects the boundary of $R$.  In particular, there is $m$ such that $f^{n.m}(I)$ intersect the boundary of $R$ and observe that if $I$ was choosen sufficiently small, then $m$ is large. One claim that  $f^{n.m}(I)$ intersect the stable arc that bounds $R$: if it interects  the unstable one, it follows that $f^{-mn}$ of that arc is inside $U$ but backward iterates of that arc gets closer to $\La_c$ that is contained in the complement of $U.$
% Therefore,  $J$ interects the the local stable arc of $R$ and is disjoint with the unstable one; therefore,   it disconnects $J$  and so, by a type of $\la-$lemma, follows that  $J$ accumulates on the local  unstable manifold 
% of $\La_c.$ If $\La_c$ coincides with $\La_i$ then the proof is finished. If $\La_c$ is a different basic set of $\La_i$, recall that $\bar J \cap interior(\La_c)$ and since each basic piece is invariant (it is considered the whole orbit of a class), then $f^{n_i}(\bar J)\cap interior(\La_c)$ where $n_i$ is such that $f^{n_i}(U_i)\subset R_i$; in particular, $\bar J$ intersect the $\La_c$ and is contained in $R_i$ and so it has to intersect the boundary of $R_i$; arguing as before, it has to intersect the stable boundary and so $\bar J$ contains the unstable manifold of $\La_i.$ 

% To prove that is backward connected, let now an open  (sufficiently small) neighborhood $I$ of $f^{n_i}(x)$ (that is contained in $R_i$) disjoint with $\La_i$ (recall that $x\in \La_i$) and let $n$ a positive integer  such that $f^{-n}(I)$ intersects $I$. As before, let $J=\cup_{m>0}f^{-n.m}(I)$ and since $f$ is dissipative in  $\DD$ (so, the jacobian of $f^{-1}$ is larger than $1$) and $f^{-n}(J)\subset J$ , it follows that $J$ cannot be contained in $\DD$ and in particular it has to cross the boundary of $R_i$; it is claimed that $J$ does  not interects the stable boundary, otherwise there is a large $l$ such that $f^{-l.n}(I)$ intersect the stable boundary of $R_i$ that is given by a local stable arc of $\La_i$;  therefore $I$ intersects the iterate of that arc that now is closer to $\La_i$; since $n.l$ is arbitrary large provided that $I$ is sufficiently small, it would imply that $f^{n_i}(x)$ is in $\La_i$, a contradiction. So, $J$ cross the unstable boundary of the rectangle and therefore by a $\la-$lemma argument it follows that $\bar J$ contains the unstable manifold of $\La_i.$


 
 
 
% \end{proof}


% \noi{\em Proof of theorem \ref{almostMisiu}:} observe that any point in $ \Omega$ such that its orbits is in the union of the open sets $U_i$, by claim \ref{insideUi} is related to one of the sets $\La_i$. On the other hand, the the part of the nonwandering that its orbit do not go through the union of the open sets $U_i$, its formed by a finite number of basic pieces. 

% \qed


\section{Misiurewicz diffeomorphisms.  Proof of theorem \ref{thm:finite classes}}\label{sec:misiu}

%\subsection{Preliminaries }
% In the next subsection, we recall a series of results and immediate consequences about the dynamics with dominated splitting.
% In the last subsection, we provide the proof of theorem \ref{thm:finite classes}.

% \subsection{Proof of theorem \ref{thm:finite classes}}

Observe that by theorem \ref{disk2}, any periodic point in a non-trivial homoclinic class has an unstable branch that accumulates on the critical set.  Therefore, 
for the proof of theorem \ref{thm:finite classes}, we first address saddle periodic points for which at least one unstable branch accumulates on the critical set, and show that they can be grouped into a finite number of homoclinic classes; later it is shown that the remaining periodic points   are contained in the interior of   normally hyperbolic arcs. Throughout this section, for a saddle periodic point $q$ and one of its
unstable branches $\Ga^u(q)$, we write $\La^u(q)$ for the accumulation set
of that branch, in the sense of section \ref{sec:disk}.

Before going  into the proof, let us clarify a few definitions. Recall that for mildly dissipative diffeomorphisms, any branch of a stable manifold of a point is not contained in the disk $\DD$; in particular, one can define $W^s_\DD(x)$ as the connected component of $W^s(x)\cap \DD$ that contains $x$; similarly, given a stable branch $\Ga^s(x)$ one can define $\Ga^s_\DD(x).$ Observe that by the mildly dissipative hypothesis, $\DD\setminus W^s_\DD(x)$ has two connected components $\DD^+_x$ and $\DD^-_x$. 

\begin{definition}\label{defi:crossing} Let $p$ and $q$ be two saddle periodic points and let $x\in W^s(p)\cap W^u(q)$; we say that $W^u(q)$ {\it crosses} $W^s(p)$ at $x$,  if there are two connected arcs $\ga^+$ and $\ga^-$ in the unstable manifold of $q$ and $x$ is an extreme point of both arcs satisfying that 

\begin{itemize}
    \item[--] $\ga^+$ intersects $\DD^+_x$ and is contained in the closure of  $\DD^+_x$ and 
    \item[--] $\ga^-$ intersects $\DD^-_x$ and is contained in the closure of  $\DD^-_x.$
    
\end{itemize}

The point $x$  is also called a crossing intersection of $W^s(p)$ and $W^u(q).$

It is said that the {\it unstable manifold of $q$ crosses the stable of $p$} if there exists $x\in W^s(p)$ such that $W^u(q)$ intersects both components $\DD^+_x$ and $\DD^-_x.$
    
\end{definition}


Observe that the unstable of a periodic point crosses the stable of other periodic point if and only if  there is a crossing intersection.


Classically, a homoclinic class of a periodic point is defined as the closure of all the transversal intersections between the stable and unstable manifold of that periodic point. We extend that definition as follows:

\begin{definition}\label{defi:generalizedHC}
Given a periodic point $p$, the {\it generalized homoclinic class} is the closure of all the crossing points of the stable and unstable of $p.$

Given two periodic points $p$ and $q$, it is said that they are generalized homoclinically related if the unstable manifold of $q$ crosses the stable of $p$ and vice versa. 
    
\end{definition}

If the  diffeomorphism is $C^\infty$, it follows that if two periodic points are generalized homoclinically related then they are also homoclinically related in the sense that there is an intersection between the stable and unstable manifold that is transversal (see \cite{BCS}).



For brevity, we keep saying  homoclinic class when we are actually referring to generalized homoclinic class.




Coming back to the proof of theorem \ref{thm:finite classes}, we prove the following next proposition and lemma.


\begin{proposition}\label{critical-Misiu} 
Under the hypothesis of theorem \ref{thm:finite classes} it holds that all saddle periodic points that are close enough to a same critical point and for which at least one unstable branch accumulates on the critical set are homoclinically related.     
\end{proposition}


For the other  saddle periodic points, one can prove the following next lemma.


\begin{lemma}\label{critical-Misiu2} Under the hypothesis of theorem \ref{thm:finite classes}, any saddle periodic point such that none of its unstable branches accumulates on a critical point belongs to the interior of a normally hyperbolic periodic arc.
    
\end{lemma}


\begin{proof}


    
% \paragraph{\it Proof of lemma \ref{critical-Misiu2}:}

The proof follows immediately from theorem \ref{disk2}. In fact,  if unstable set of the periodic point is trivial, then it  belongs to a normally hyperbolic periodic arc; its unstable branches accumulate either on a semi-attracting periodic point or on a normally hyperbolic periodic arc. Adding these points or arcs to the unstable branches yields a closed periodic curve without critical points, which is therefore a normally hyperbolic arc.


\end{proof}




Before proving proposition \ref{critical-Misiu}, we show how the previous proposition and lemma combined with proposition \ref{PS2-1} provide the proof of theorem \ref{thm:finite classes}.


\paragraph{\it Proof of theorem \ref{thm:finite classes}:} After lemma \ref{critical-Misiu2} it is enough to consider saddle periodic points with branches accumulating on the critical set. By proposition \ref{critical-Misiu} for each critical point $c$ there is $\e(c)>0$ such that  the periodic points  whose orbits intersect $B_{\e(c)}(c)$ are homoclinically related and so they belong to a homoclinic class. There is a finite number of critical points $\{c_1,\dots, c_n\}$ such that  $B=\cup_{i=1}^n B_{\e(c_i)}(c_i)$ contains the critical set and all periodic points with their orbits intersecting $B$ are therefore  contained in a finite number of homoclinic classes. The closure of the orbits of the saddle  periodic points that do not intersect $B$ has dominated splitting and therefore   from proposition \ref{PS2-1} it is also concluded that they also belong to a finite number of homoclinic classes. 

\qed

\subsection{Proof of proposition \ref{critical-Misiu}.}

% The remainder of the section focuses on the proof of the proposition  \ref{critical-Misiu}.

We split the proof of Proposition \ref{critical-Misiu} into two parts according to whether  a sequence of saddle periodic points accumulates on critical points whose $\omega$ or 
$\alpha$ limits are non-trivial (interior points or heterocycles), or only on boundary critical points (see definition \ref{defi:boundary}). In the first case, treated below, Proposition \ref{critical-Misiu} follows directly from Lemmas \ref{interior} and \ref{HC}. The second case reduces, via the notion of a  $*-$sequence (definition \ref{def:*-seq}), to proposition \ref{prop:*-seq}, whose proof occupies the rest of the section.




% \subsection{Proof of proposition \ref{critical-Misiu}}


For a  Misiurewicz diffeomorphism, it follows that the  omega-limit of any critical point is a set with dominated splitting and so by proposition \ref{dicho-omega} it is either a periodic orbit, or a heterocycle or a set that contains interior points.



\begin{definition}\label{defi:boundary} 
    
Let ${\rm Crit_I}(f)$ be the set of critical points $c$ such that  $\omega(c)$ or $\alpha(c)$  either contains interior points, is a heterocycle, is contained in a normally hyperbolic arc, or lies in the interior of the basin of attraction of a (semi-)attracting periodic point. Let us call ${\rm Crit_{II}}(f)$ its complement. 

The critical points in ${\rm Crit_{II}}(f)$ are called {\it boundary critical points}  and the periodic points whose stable or unstable manifolds contain a boundary critical point are called {\it boundary saddle periodic points}. 
\end{definition}



Now, let us consider a sequences $(p_n)$ of saddles that accumulates on a critical point. Observe that if there is a sequence $(p_n)$ of periodic points with unbounded period that accumulates at a periodic point, then that periodic point has to be a saddle periodic point that can not be contained in the interior of a normally hyperbolic arc. Moreover, if a sequence of periodic points $(p_n)$ accumulates at a critical point, then their periods must be unbounded.

Given a sequence $(p_n)$ of saddle periodic points, let us define ${\rm Crit}(f, (p_n))$ the set of the critical points accumulated by $(p_n)$ and they can be divided into two sets: ${\rm Crit_I}(f, (p_n))= {\rm Crit_I}(f)\cap {\rm Crit}(f, (p_n))$ and ${\rm Crit_{II}}(f, (p_n))= {\rm Crit_{II}}(f)\cap {\rm Crit}(f, (p_n))$

\paragraph{Case ${\rm Crit_I}(f, (p_n))\neq \emptyset$: }


% let $\La$ be a locally maximal invariant set with dominated splitting contained in a homoclinic class. It is said that a point $x\in \La$ is an interior point if $\La$ intersects both  connected component of $W^{s}_{loc}(x)\setminus\{x\}$. Otherwise, it is said  that $x$ is a boundary point. Recall that if $x$ is a boundary point then there is a  boundary periodic point $p\in \La$ such that $x\in W^s(p).$

% Form proposition \ref{dicho-omega} it follows that 
% Recall that given  $y\in L(f)$ if  the set $\omega(y)$ has dominated splitting then this set  is contained in a locally maximal invariant set $\La$ with dominated splitting. Moreover, if $\omega(y)$ does not contain interior points in $\La$ then $\omega(y)$ is a boundary point and therefore $y$ belongs to the stable manifold of a periodic boundary point. Since the omega limit of any critical point has dominated splitting, it follows that either that critical point belongs to stable manifold of an interior point or belongs to the stable manifold of a boundary periodic points.

% Observe that, the number of boundary periodic points that its stable manifold contains a critical point is finite. Also, the number of stable boundary points fr away the critical set is finite.



% There is a finite number of locally maximal hyperbolic invariant set $\La^+_1,\dots \La^+_k$  such that $Crit(f)\cap W^s_{loc}(\La^+_i)\neq\emp$ and $Crit(f)\subset \cup W^s_{loc}(\La^+_i)$, these sets are called stable catching sets. Simultaneously,  there is a finite number of locally maximal hyperbolic invariant set $\La^-_1,\dots \La^-_m$  such that $Crit(f)\cap W^u_{loc}(\La^-_i)\neq\emp$ and $Crit(f)\subset \cup W^s_{loc}(\La^-_i)$, these sets are called unstable catching sets. All these classes are called catching sets.



%We now provide a proposition and before proving let us show how using that proposition it is concluded   theorem \ref{thm:finite classes}.




%\paragraph{Proof of theorem \ref{thm:finite classes}:}

%\qed




% \begin{remark} Any Markov partition with sufficiently small rectangles do not intersect the critical set.
    
% \end{remark}



% % \begin{remark} Given a boundary periodic point, then there are two quadrants that each one is disjoint with the set.  Let's call this quadrant, free quadrants.

    
% % \end{remark}

% % \begin{definition} It is said that $x\in \cup \La_i$ is a stable catching point if there is a critical point belonging to the local stable set of $x$. Similarly, it is said that it is said that is an unstable catching point if there is a critical point belonging to its local unstable set.. and its forward orbit is in a free quadrant and its backward  orbit is in a free quadrant.
    
% % \end{definition}





% % Type of catching points:

% % \begin{itemize}

% % \item[--] Interior points

% % \item[--] stable boundary points


% % \item[--] unstable boundary points


% % \item[--] corner points


    
% % \end{itemize}

% %quadrants/branches


% \begin{lemma}\label{go to the interior} Let $\cal R$ be a Markov partition of $\La$. Let $x$ be a point in $\La$ that does not belong to the stable manifold of a boundary point.  Then it holds that 

% \begin{itemize}
% \item[--]for any $n$ sufficiently large follows that $f^n(x)$ is in the interior of the partition 

% \item[--] there are  infinitely many forward iterates of $x$ that are at a bounded distance of the unstable boundary of $\cal R.$
% \end{itemize}
    
% \end{lemma}


% \begin{proof} Let us prove first that for large $n$ it holds that $f^n(x)$ is in the interior of the partition; by hypothesis, for any $n$ it holds that  $f^n(x)$ is not in the stable boundary; so if it is not the in the interior it holds that  belongs to the unstable boundary and therefore belong to a finite number of compact arc of some unstable branch of finite number of  boundary periodic points; by hypothesis, $x$ is not one of these  periodic point and therefore, any large iterate does not belong to the compact arcs and therefore is not in the unstable boundary and so the first part of the lemma follows. 

% Now if there is an accumulation points $z$ of the forward orbit of $x$ that does not belong to the unstable boundary, the second part of the lemma follows. But as it was claimed before,  there is an iterate of $z$  that does not belong to the unstable boundary. 

% % branch of un but in that case , it follows that $\omega(x)$ is contained in the boundary of the rectangle. If $\omega(x)$ intersects the stable boundary of the Markov partition, then there is a point in $\omega(x)$ that belongs to the local stable manifold of of a boundary periodic point and since $x$ is not in that stable manifolds, there are iterates of $x$ that accumulates on a point $z$ in the unstable manifold of the periodic point. If $z$ is in the interior of the 

    
% \end{proof}

% \begin{remark} A corner point is a saddle periodic point. Moreover, there is a quadrant that does not intersect  the hyperbolic set that contains the corner point.

    
% \end{remark}


\begin{lemma}\label{interior} Under the hypothesis of theorem \ref{thm:finite classes}, let $(p_n)$ be  a sequence of saddle periodic points such for  each of them at least one unstable branch accumulates on the critical set and they accumulate on an interior point $x$ of a   set $\La$ with dominated splitting; it follows  that the periodic points of the sequence $(p_n)$ close enough to $x$  are homoclinically related. 

\end{lemma}

 \begin{proof}

From the fact that $\La$ has dominated splitting and $x$ is an interior point,  
one can get a small rectangle $R$   satisfying: 

\begin{itemize}
    \item[--]  it is bounded by local stable and unstable arc of recurrent points in $\La,$
    \item[--] $x$ is in the interior of $R$,
    \item[--] it is disjoint from the  critical points.

    
\end{itemize}

Given a periodic point $p_n$ let us consider $g=f^{per(p_n)}$. Observe that by mild dissipation, the stable branch of any $p_n$ close enough to $x$ is not contained in $R$ and therefore it intersects the unstable boundary of $R.$  Similarly, if $R$ is sufficiently small then there are no critical points in $R$ and so if at least one of the  unstable branches of $p_n$  accumulates on the critical points it follows  that  the unstable branch is  not contained in $R$ and so  it  intersect the stable boundary of $R.$ From that fact, the $\la-$lemma, and using   that the points that provides the stable and unstable arcs that bound the rectangle are recurrent, it follows that there are   stable and unstable arcs of any periodic $p_n$ close to $x$ that accumulates on the stable and unstable arcs that bound the rectangle; that implies that the stable and unstable manifold on any pair of periodic points close to $x$ intersects each other.

    
\end{proof}

\begin{lemma}\label{HC} Under the hypothesis of theorem \ref{thm:finite classes}, let $(p_n)$ be  a sequence of saddle periodic points such that  for each of them at least one unstable branch accumulates on the critical set and the sequence $(p_n)$ accumulates at a point $x$ such that $\omega(x)$ is a heterocycle; it follows  that the  periodic points in the sequence $(p_n)$ close enough to $x$  are homoclinically related. 

\end{lemma}

\begin{proof} From the remark \ref{rmk:heterocycle} one gets  a rectangle $R$  containing a periodic point $q$ of the heterocycle in one corner and bounded by two parallel stable arcs and two parallel unstable arcs of that point. Arguing  as in lemma \ref{interior} it follows that  any periodic $p_n$ close to the heterocycle is homoclinically related to $q$ and so homoclinically related to each other.


%Let $q$ be a periodic point in the heterocycle. Using the $\la-$lemma one can get a small rectangle that the boundary is given by two arcs in  the local stable and unstable branches of $q$ that contains $q,$ and two arcs in the stable and unstable branches of $q$. Moreover, that rectangle contains forward iterates of $x$ and does not contain critical points. Arguing as in lemma \ref{interior}, the lemma is concluded. 

    
\end{proof}

%one can get a sequence of  rectangles $(R_n)$ such that each one satisfies: 

% \begin{itemize}
%     \item[--]  it is bounded by stable and stable arc of the boundary saddle periodic points;
%     \item[--] $p_n\in R_n$;
%     \item[--] it is disjoint of the accumulation points of the orbits of any critical point.

    
% \end{itemize}

% Given a periodic point $p_n$ let us consider $g=f^{per(p_n)}$. Observe that by mild dissipativeness, the stable branch of $p_n$ is not contained in $R_n$ and therefore it intersects the unstable boundary of $R_n.$  Similarly, by theorem ...  the unstable branch of $p_n$ accumulates on the forward orbit of a critical point  and therefore the unstable branch is not contained in $R_n$ and so  it intersects the stable boundary of $R_n.$ 


% \begin{definition}\label{defi:boundary} 
    
% Let $Crit_I(f)$ be the set of critical points $c$ such that  $\omega(c)$ or $\alpha(c)$  either contains interior points, is a heterocycle, is contained in a normally hyperbolic arc, or lies in the interior of the basin of attraction of a (semi-)attracting periodic point. Let us call $Crit_{II}(f)$ its complement. 

% The critical points in $Crit_{II}(f)$ are called {\it boundary critical points}  and the periodic points whose stable or unstable manifolds contain a boundary critical point are called {\it boundary saddle periodic points}. 
% \end{definition}

As a consequence of the previous two lemmas, it follows an immediate corollary:

\begin{corollary}\label{cor:crit1}
     Let $(p_n)$ be  a sequence of saddle periodic points such that for  each of them  at least one unstable branch accumulates on the critical set and  the orbits of the sequence $(p_n)$ accumulate on $Crit_I(f)$; it follows that  they belong to a finite number of homoclinic classes.
\end{corollary}


% \begin{remark}\label{rmk:finite-boundary}
% It follows from proposition \ref{PS2-1} that the number of  boundary saddle periodic points is finite.     
% \end{remark}


\paragraph{Case ${\rm Crit_I}(f, (p_n))=\emptyset$.} In that case, ${\rm Crit}(f, (p_n))= {\rm Crit_{II}}(f, (p_n)).$ By proposition \ref{PS2-1} there are only a finite number of boundary saddle periodic points that contain the critical points in $Crit_{II}(f).$  So, to conclude proposition \ref{critical-Misiu} it remains to consider the sequence of periodic points that accumulate only on boundary critical points.


\begin{definition}\label{def:*-seq}  
Let $(p_n)$ be a sequence of saddle periodic points such that for each of them at least one of its unstable branches accumulates at a critical point. Let $\La$ be the accumulation set of the orbits of the periodic points $(p_n)$; we say that the sequence $(p_n)$ is a {\it $*-$sequence} if 

\begin{itemize}
    \item[--] the set of critical points in $\La$ is finite and they are all boundary  critical points,

    \item[--] any  compact invariant set with dominated splitting  contained in $\La$ is formed by a finite number of saddle periodic orbits.
\end{itemize}


\end{definition}


Therefore, as a consequence of  corollary \ref{cor:crit1},  to finish the proof of proposition \ref{critical-Misiu} it remains to prove the following proposition.


\begin{proposition}\label{prop:*-seq} Under the hypothesis of theorem \ref{thm:finite classes}, let $(p_n)$ be a $*-$sequence, then it holds that they belong to a finite number of homoclinic classes.
\end{proposition} 



To prove proposition \ref{prop:*-seq}, the goal is to show that there is a ``touching cycle" (defined below) between  periodic points that ``encloses" some iterates of  the $*-$sequence and that is used to show  that the periodic points in  $*-$sequence are all homoclinically related.

\paragraph{Quadrants, sequences and cycles.} Given a saddle (or semi-attracting) periodic point  $q$, its local  stable and center-unstable branches (denoted as $\Ga_l^s(q)$ and $\Ga^{cu}_l(q)$ respectively) define at most $4$ quadrants ($2$ in case that $q$ is semi-attracting); more precisely, the complement of the union of the local  stable and center-unstable set $q$ in a small neighborhood of it has $4$ connected components; let us denote by $(q,Q)$ one of these quadrants that can also be identified with the triplet $(q, \Gamma^s_l(q), \Gamma^{cu}_l(q))$ that indicates  the branches that bounds the quadrant. When the center-unstable branch is an unstable branch, we denote $\Gamma^{cu}_l(q)$ with $\Ga^u_l(q).$ 
% (meaning that the sequence is in  one connected component of a small neighborhood of $x$ minus the local branch of $x$),

Given a point $x$ that belongs to the local stable branch of a periodic saddle $q$ and a sequence $(p_n)$ that accumulates at $x$ on one side of the local stable manifold  of $q$, there exists a quadrant $Q$ bounded by branches $\Gamma^s_l(q), \Gamma^u_l(q)$ such that 
\begin{itemize}
    \item[--] the sequence $(p_n)$ is contained in $Q$,
    \item[--] there exist $y\in \Gamma^u_l(q)$ and positive integers $(k_n)$ such that $f^{k_n}(p_n)\to y$,
    \item[--] for any $n$ and any  $0\leq j \leq k_n$ it holds that $f^j(p_n)\in Q.$
\end{itemize}

% such that   and that the sequence $(p_n)$ belong to only one component of   , we consider the quadrant $Q$ such that $x$ belong to the local stable branch of $q$ that contains $x$ and the iterates of $(p_n)$ accumulate on $q$ and intersect $Q.$ Let us denote the triplet $(q, x, Q)$.   




\begin{definition}\label{def-cycle-points}
Given a finite set of saddle periodic points $(q_i)$  and respective quadrants $(q_i, \Ga^s_l(q_i), \Ga^u_l(q_i))$, we say that there is a {\it cycle} if for each $i$, 
\begin{itemize}
    \item[--] there is a point $(y_i)$ in $\Ga^{u}_l(q_i),$
    \item[--] there is a point $(x_i)$ in $\Ga^{s}_l(q_i),$
     \item[--] a sequence of forward iterates $(f^{n_j}(y_i))$ contained in the quadrant bounded by  $\Ga^s_l(q_{i+1})$ and $\Ga^u_l(q_{i+1})$ and  accumulating on  $x_{i+1}\in \Ga^{s}_l(q_{i+1})$.
\end{itemize}


  


\noindent We say that the cycle is a {\it critical  cycle} if the points $(x_i)$ are critical points. 


\noindent If for any $i$ there exists $n_i>0$ such that $f^{n_i}(y_i)=x_{i+1} \in\Ga^s_l(q_{i+1})$ and there is an arc in $\Ga^u(q_i)$ that contains $x_{i+1}$ and intersects $Q_{i+1}$, we say that the cycle is a {\it touching cycle}. 




\end{definition}

To incorporate all the data defined above, a cycle between periodic points and quadrants is noted as a finite sequence of quadruplets $(q_i, Q_i, x_i, y_i).$

% where  $Q_i$ is the quadrant bounded by $\Ga^s_l(q_i)$ and $\Ga^u_l(q_i)$, it is said that there is a cycle between the sequence    of saddle periodic points, quadrants and marked points $y_i\in \Gamma^u_l(q_i), x_i\in \Gamma^s_l(q_i)$.  


% \begin{definition}\label{def-cycle-points-bis}
% Given a finite set of saddle periodic points and local stable and unstable branches $(q_i, \Ga^s(q_i), \Ga^u(q_i))$, it is said that there is a cycle if for each $i$ holds that the iterates of $\Ga^u(q_i)$ accumulates on $q_{i+1}$ and the accumulation holds inside the quadrant bounded by  $\Ga^s(q_{i+1})$ and $\Ga^u(q_{i+1})$. It is also said that the cycle is a touching cycle if there is an iterate of $\Ga^u(q_i)$ that intersects $\Ga^s(q_{i+1})$. 
% \end{definition}



\begin{figure}[h]
\centering
\includegraphics[scale=0.45]{fig_cycle.pdf}
\caption{Quadrants and cycles}
\label{cycles}
\end{figure}

\begin{proposition}\label{touching-cycle}  Given a touching cycle  of a mildly dissipative diffeomorphisms of the disk, for any saddle $q$ in the cycle there is a sequence $(q_n)$ of hyperbolic saddle periodic points  and  a point $x\in W^s_{\DD}(q)$ such that 

\begin{itemize}
\item[--] the sequence  $(q_n)$ accumulates at $x$ and  is contained in the quadrant $Q$ of $q$ involved in the cycle,
    \item[--] all saddles $(q_n)$ are homoclinically related, 
    \item[--]$W^s_{\DD}(q_n)$ intersects transversely $ \Ga^u_{l}(q)$ (the unstable branch of $q$ that bounds $Q$)  and  $W^s_{\DD}(q_n)\to W^s_{\DD}(q)$, 
    \item[--]  there is an arc $\ga^u_n\subset W^u(q_n)$ contained in $Q$ that joins $q_n$ with a point in  $W^s_\DD(q_{n+1}).$
\end{itemize}


    
\end{proposition}

% \begin{proof} The proof is an extension of the proof of  lemma 5.8 in \cite{CPT}. In that lemma, given $q$  (that one can assume that is fixed) and $x\in W^s(q)$ it  is shown that there is a sequences of saddle periodic points $q_n$, each one contained in a non-trivial homoclinic class, that  accumulate on $x.$ The proof proceeds by showing that there are small rectangles $R_n$ parallel to the local stable manifold of $q$ that accumulates at $x$ verifying that there exist iterates $r_n$ such that $f^{r_n}(R_n)$ intersects transversely $R_n$ in at least two connected components. Using the dissipation, one can also show that the rectangles can be chosen not only that its iterates intersect $R_n$ but also $R_{n+1}.$


% \end{proof}

% \begin{proof} The proof extends that of lemma 5.8 in \cite{CPT}. We recall
% its mechanism and indicate the additional input. Given the saddle $q$ in the
% cycle (assumed fixed) and $x\in W^s_\DD(q)$ accumulated, inside the quadrant
% $Q$, by the unstable branch $\Ga^u_l(q)$ involved in the cycle, lemma 5.8 of
% \cite{CPT} produces small rectangles $R_n$ that accumulate on $x$,
% parallel to $W^s_\DD(q)$, together with return times $r_n$ such that
% $f^{r_n}(R_n)$ crosses $R_n$ transversally in at least two components. Each
% such crossing yields, by the classical horseshoe construction, a hyperbolic
% saddle $q_n$ contained in a non-trivial homoclinic class, with $q_n\to x$;
% moreover the local stable manifolds satisfy $W^s_\DD(q_n)\to W^s_\DD(q)$ and
% cross $\Ga^u_l(q)$ transversally, and consecutive $q_n,q_{n+1}$ are
% homoclinically related (their invariant manifolds cross inside $Q$), so all
% the $q_n$ lie in one homoclinic class.

% The additional point, not needed in \cite{CPT}, is the last item of the
% statement: that an unstable arc $\ga^u_n\subset W^u(q_n)$ inside $Q$ joins
% $q_n$ to a point of $W^s_\DD(q_{n+1})$. This is obtained from the dissipation
% exactly as the transversal crossing above: the rectangles $R_n$ can be
% chosen so that $f^{r_n}(R_n)$ meets not only $R_n$ but also $R_{n+1}$ (the
% returns accumulate on $x$ from the same side, and $R_{n+1}\subset R_n$ up to
% the stable direction), and the component of $W^u(q_n)\cap Q$ realizing the
% crossing of $R_{n+1}$ furnishes the arc $\ga^u_n$ reaching
% $W^s_\DD(q_{n+1})$.
% \end{proof}

\begin{proof} The proof extends that of lemma 5.8 in \cite{CPT}. We recall
its mechanism and indicate the additional input. Given the saddle $q$ in the
cycle (assumed fixed) and $x\in W^s_\DD(q)$ accumulated, inside the quadrant
$Q$, by the unstable branch $\Ga^u_l(q)$ involved in the cycle, lemma 5.8 of
\cite{CPT} produces small rectangles $R_n$ that accumulate on $x$,
parallel to $W^s_\DD(q)$, together with return times $r_n$ such that
$f^{r_n}(R_n)$ crosses $R_n$ transversally in at least two components. Each
such crossing yields, by the classical horseshoe construction, a hyperbolic
saddle $q_n$ contained in a non-trivial homoclinic class, with $q_n\to x$;
moreover the local stable manifolds satisfy $W^s_\DD(q_n)\to W^s_\DD(q)$ and
cross $\Ga^u_l(q)$ transversally.

The additional point, not needed in \cite{CPT}, is the last item of the
statement: that an unstable arc $\ga^u_n\subset W^u(q_n)$ inside $Q$ joins
$q_n$ to a point of $W^s_\DD(q_{n+1})$. This is obtained from the dissipation
exactly as the transversal crossing above: the rectangles $R_n$ can be
chosen so that $f^{r_n}(R_n)$ meets not only $R_n$ but also $R_{n+1}$, and
the component of $W^u(q_n)\cap Q$ realizing the crossing of $R_{n+1}$
furnishes the arc $\ga^u_n$ reaching $W^s_\DD(q_{n+1})$. In particular
$W^u(q_n)$ crosses $W^s_\DD(q_{n+1})$ and, symmetrically, $W^u(q_{n+1})$
crosses $W^s_\DD(q_n)$; hence consecutive saddles $q_n,q_{n+1}$ are
homoclinically related, and therefore all the $q_n$ belong to a single
homoclinic class.

\end{proof}


\begin{corollary}\label{touching} Under the hypothesis of theorem \ref{thm:finite classes}, given a touching cycle, it follows that  all saddle periodic points satisfying 

\begin{itemize}
    \item[--]  for  each of them  at least one unstable branch accumulates on the critical set,
    \item[--] they are close enough to at least one saddle involved in the  cycle and  are  contained in the correspondent quadrant involved in the cycle,
\end{itemize}
are homoclinically related.
    
\end{corollary}

\begin{proof} Let $(q, Q)$ be a saddle and a quadrant involved in the touching cycle and observe that nearby $q$ there are no critical points. Let $(q_n)$ be the sequence of saddle periodic points provided by proposition \ref{touching-cycle}. Given a saddle periodic point $p$ sufficiently close to $q$, one can get two consecutive saddles $q'_m$ and $q'_{m+1}$ that are two forward iterates of two points in the sequence $(q_n)$ such that the arcs $W^s_\DD(q'_m), W^s_\DD(q'_{m+1}), {\ga^u_n}'$ and $\Ga^u_q$ provided by proposition \ref{touching-cycle} form a rectangle $R_n$ close to $q$ that contains $p$ in its interior. Since there are no critical points nearby $q$ and the unstable manifold of $p$ accumulates at a critical point then it has to cross either $W^s_\DD(q_m)$ or $W^s_\DD(q_{m+1});$ on the other hand, $W^s_\DD(p)$ has to cross either ${\ga^u_m}'$ of $\Ga^u_l(q)$ and observe that  if $W^s_\DD(p)$ crosses $\Ga^u_l(q)$ then the stable manifold of $p$ accumulates on $W^s_\DD(q)$ from the side of $Q$ and so it crosses an arc $\ga^u_k$ of some saddle point $q_k$     in the sequence $(q_n).$ Since all the saddles $(q_n)$ are homoclinically related it follows that $p$ is also homoclinically related to them.

\end{proof}

To conclude proposition \ref{critical-Misiu}  it is enough to prove the following proposition.

\begin{proposition}\label{prop:*-touching} Under the hypothesis of theorem \ref{thm:finite classes} given a $*-$sequence it follows that  there exists a touching cycle such that the $*-$sequence satisfies the hypothesis of corollary \ref{touching}.

    
\end{proposition}



% show that given the sequences of the saddle periodic point $(p_n)$ there is a touching cycle such that some iterates of the saddles $(p_n)$ satisfy the hypothesis of the previous corollary.

% GOAL: get a cycle such that the  orbits of the sequences of periodic points intersects the quadrants involved in the cycle.


\subsubsection{Proof of proposition \ref{prop:*-touching}}

The proof proceeds in three stages. First, starting from the 
$*-$sequence we iteratively construct a chain of boundary saddle periodic points and quadrants; this chain closes into a {\it maximal critical cycle} with a marked sequence of boundary critical points. Second, we develop the geometric tools (a rectangle lemma, Pixton disks, and a transversality dichotomy for unstable branches (Proposition \ref{prop trans})) that govern how unstable manifolds interact with the boundary structure produced by the cycle. Third, we show that any maximal critical cycle is either itself a touching cycle or can be enlarged to one (Proposition \ref{existence of touching}); combined with Proposition \ref{touching-cycle} and Corollary \ref{touching}, this concludes the proof.

\paragraph{Building the maximal critical cycle.} The next corollary follows from  claim \ref{cpt-lemma}.

% \begin{lemma}\label{inclination} Let $f$ be mildly dissipative and let $p,q$ be  two saddle points and $\Ga_p$ and $\Ga_q$ two of  their branches. If the accumulation set of $\Ga_p$ intersects $\Ga_q$ then it follows that the accumulation set of $\Ga_q$ is contained in the accumulation set of  $\Ga_p.$
    
% \end{lemma}

\begin{corollary}\label{jumping} Let $f$ be mildly dissipative and let $q_0, q_1, q_2$ be  three  saddle points and  $Q_1, Q_2$ two quadrants of $q_1, q_2$. Let  $\Ga^u_l(q_1)$ and $\Ga^u_l(q_2)$ be two  local unstable branches of $q_1$ and $q_2$ respectively, such that $\Ga^u(q_1)$  bounds $Q_1$ and $\Ga^u(q_2)$  bounds $Q_2$. If the iterates of $\Ga^u(q_0)$ accumulate on $q_1$ along $Q_1$ and $\Ga^u(q_1)$ accumulates on $q_2$ along $Q_2$, it follows that $\Ga^u(q_0)$ accumulates at $q_2$ along $Q_2$ and the accumulation set of $\Ga^u(q_0)$ contains $\Ga^u_l(q_2).$ 
    
\end{corollary}


Let $(p_n)$ be a $*-$sequence of periodic points, and let $c_0$ be a boundary critical point accumulated by the sequence. Let $q_0$ be the boundary periodic point that contains $c_0$ in its stable manifold;  let $Q_0$ be a  quadrant of $q_0$ that contains infinitely many points of the sequence $(p_n)$ and taking a subsequence one can assume that they all belong to the same quadrant.  Observe that there exists  a point $y_0$ in the local unstable branch $\Ga^u(q_0)$ that bounds $Q_0$ such that $f^{k_n}(p_n)\to y_0$ and $f^j(p_n)\in Q_0$ for any $0\leq j\leq k_n$ (observe that  there is also a sequence of iterates $(f^{j_n}(p_n))$ such that  $f^{j_n}(p_n)\to q$ and $0< j_n< k_n$).  Let us consider $\omega(y_0)$ and observe that any point in that set is also accumulated by iterates of $(p_n)$; either $\omega(y_0)$ contains a critical point $c_1$ that belongs to the stable manifold of a boundary periodic point or $\omega(y_0)$ has dominated splitting.  In the first case, it has to be a boundary critical point that belongs to one of the finite boundary periodic points. In the second case by the definition of $*-$sequence it follows that $\omega(y_0)$ is  a periodic orbit and so, there is a saddle periodic point $q_1$ and a forward  iterate of $y_0$ named $c_1$ that is in the local stable manifold of $q_1$.     In conclusion,  we obtain a new  saddle periodic point $q_1$,  a quadrant $Q_1$ and a point $c_1$ satisfying:

\begin{itemize}
    \item[--] $q_1$ and $c_1$ are  accumulated by  iterates of the sequence $(p_n)$ that are  contained in $Q_1$ and  

    \item[--] $q_1$ is accumulated  by $\Ga^u(q_0)$ along the quadrant $Q_1$.
\end{itemize}
Moreover, it follows that either 
\begin{itemize}
    \item[--] $c_1$ is a critical boundary point that belongs to the local stable manifold of $q_1$ (a boundary saddle period point)  and infinite forward iterates of $y_0$ accumulate on $c_1$, let us call this case type 1,  or 
    %$c_1$ is a critical boundary saddle periodic point, 

    \item[--] $q_1$ is a   saddle periodic point (that could be a boundary saddle periodic point) and  there is a forward iterate of $y_0$ equal to  $c_1$ (that could be a boundary critical point), let us call this case type 2.
\end{itemize}

Informally, type 1 is the favorable case: the construction has reached a
 boundary critical point $c_1$ that belongs  in the stable manifold of the
boundary saddle $q_1$, and the chain can close there. Type 2 is the
inconclusive case: the forward orbit of $y_0$ has only passed through a
saddle $q_1$ without yet landing on a boundary critical point, so the
construction must be continued from $q_1$.

We repeat the argument from $q_1$, producing a chain of periodic points and
quadrants $(q_n,Q_n)$ each of type 1 or type 2; the goal is to show that the
chain either closes into a touching cycle directly or, after removing  the
inconclusive consecutive  type 2  via corollary \ref{jumping}, one gets a critical cycle
of boundary critical points. More precisely, if there are finitely many consecutive ones in that chain that are of type 2 followed by a saddle of type 1, using corollary \ref{jumping}, one can remove those consecutive points of type 2 and the next one in the chain is a critical boundary saddle point (type 1).  If there are infinitely many consecutive ones  in that chain  that are of type 2 and such that their orbits are uniformly bounded away from the critical set, by the definition of $*-$sequence it follows that they have to repeat and they form a touching cycle proving proposition \ref{prop:*-touching}.  If there are infinitely many consecutive ones in that chain that are of type 2 and such that their orbits accumulate on a  critical point, given that the periodic points in the chain are accumulated by the orbits of the sequence $(p_n)$, it follows that  their orbits also accumulate on that critical point and therefore it is a boundary critical point and, again, using corollary \ref{jumping} one can replace the consecutive chain of periodic points of type 2 by a critical boundary point (type 1). 

Therefore either  a touching cycle that encloses the $*-$sequence is obtained, or using the fact that the number of critical boundary points is finite there is a cycle of quadrants and marked points $(q_i, Q_i, c_i, y_i)$ such that  
\begin{itemize}
\item[--] $q_i$ is a boundary saddle periodic point,

\item[--] $Q_i$ is a  quadrant bounded by $\Gamma^s_l(q_i)$ and $\Gamma^u_l(q_i)$ that contains $c_i$ and $y_i$ respectively,

\item[--]  $c_i$ is a boundary critical point,

\item[--] the  forward  iterates of $y_i$ accumulate  on $c_{i+1}$ inside $Q_{i+1}$.

    
\end{itemize}
moreover,  the orbits of the sequence $(p_n)$ satisfy the following properties (named $(P)$)
 \begin{itemize}
\item[--] accumulates on the boundary critical points $(c_i)$ inside $Q_i$,

\item[--] accumulates on the periodic points $(q_i)$ inside $Q_i$.

% \item[--]  intersects the quadrants $(Q_i)$. 

    
\end{itemize}
 
The next remark shows that any critical cycle satisfying property (P) can be enlarged to other critical cycle also satisfying property (P) and thereby obtaining a maximal critical cycle.

\begin{remark}\label{enlarging}
    
Observe now that given a critical cycle, $(q_i, Q_i,c_i, y_i)$ if there are two consecutive periodic points  $q_i, q_{i+1}$ in the cycle and a third periodic point $\bar q$ such that 

\begin{itemize}
    \item[--] there is a critical point $\bar c$ that belongs to some stable branch of $\bar q$ that  is  accumulated by iterates of $y_i,$

    \item[--] there is a point $\bar y$ in the unstable branch of $\bar q$ that is also accumulated by iterates of $y_i$ such that forward iterates of $\bar y$ accumulate on $c_{i+1},$


\end{itemize}

 then   one can enlarge the cycle including  $(\bar q, \bar Q, \bar c, \bar y)$ and that the orbits of the sequence $(p_n)$ also satisfy the properties  $(P)$ listed above, and so  $\bar c $ is a boundary critical point. In that way, using that the boundary critical points are finite,  a critical cycle can be enlarged to a {\it 
 maximal critical cycle}.
 \end{remark}

Proposition \ref{prop:*-touching} follows from next  proposition. 

% and by proposition \ref{touching-cycle} the proof of proposition \ref{critical-Misiu} and so  theorem \ref{thm:finite classes} are concluded.


\begin{proposition}\label{existence of touching} Given a maximal critical cycle then either the cycle is touching or it is contained in a larger touching cycle (that it is not necessarily a critical cycle).


\end{proposition}

The rest of the section focuses on proving the previous proposition.  For that, first we need to introduce a few tools. 

\paragraph{Rectangles, Pixton disks and transversality  dichotomy.} 
The next three results give the geometric infrastructure used in the proof of Proposition \ref{existence of touching}: a rectangle lemma controlling the local stable geometry near an accumulation point, the notion of a Pixton disk associated with a non-touching cycle, and a transversality dichotomy for unstable branches.

We start with a simple lemma for mildly dissipative diffeomorphisms of the disk.





% \begin{remark} If $\La^u(p_i)\cap W^s_{loc}(\La_j)\cap Crit(f)\neq \emp$ either $\Gamma^u(p_i)$ intersects transversaly the stable of of the boundary points of $\La_j$ or there is a critical point catch by a boundary points of $\La_j$ that its accumulated by the unstable branch of $p_i$
    
% \end{remark}


% CRITICAL CYCLE AND MAXIMAL CRITICAL CYCLE


\begin{lemma}\label{rectangle} For any  mildly dissipative diffeomorphism of the disk it follows that given a saddle periodic point $q$ and $x\in W^s_{\DD}(q)$ that is accumulated by the unstable branch of a saddle periodic point $p$, there exists a rectangle $R$ arbitrarily close to $x$ (see figure \ref{rectangle-fig}) bounded by two pairs of parallel  arcs $(\ga_1, \ga_2)$ and $(\be_1, \be_2)$ satisfying that
\begin{itemize}
    \item[--] $\ga_1$ is contained in the local stable manifold of $q$ and $x$ belongs to the interior of $\ga_1,$
    \item[--] there is a sequences of connected arcs $(\si_n)$ contained in the unstable branch of $p$ such that 
    \begin{itemize}
     \item[--] each $\si_n$ is contained in $R$ and is disjoint with $\be_1$ and $\be_2,$
     \item[--] one extreme point of the arc $\si_n$ belongs to $\ga_2$ and the other extreme points $(x_n)$ satisfy that $x_n\to x.$

    
    
    \end{itemize}

\end{itemize}

\end{lemma}

\begin{proof} One can choose an open arc $\delta$ in $W^s_\DD(q)$ whose backward iterates are contained in the complement of the disk. Then, one can choose two consecutive iterates $\delta_i, \delta_{i+1}$ of $\delta$ such that the point $x$ is in the interior of an  arc $\gamma$ contained in the local stable branch bounded by $\delta_i$ and $\delta_{i+1}$ (observe that since  $x$ is accumulated by the unstable branch of a periodic point $p$, its full orbit is contained in the disk). Extending the arc $\gamma$ to an arc $\gamma_1$  (also inside the local stable branch of $q$), one can assume that the extreme points of $\gamma_1$ are contained in $\delta_i$ and $\delta_{i+1}$ respectively. Now, one can choose two parallel transversal  arcs $\beta_1, \beta_2$ in $Q$ (the quadrant of $q$ where the unstable branch of $p$ accumulates at $x$) and such that one of the  extreme points of each arc coincides with the opposite extreme points of $\gamma_1$; observe that if both arcs are small enough then their backward iterates are fully contained in the complement of the disk. Then, one chooses an arc $\gamma_2$ parallel to $\gamma_1$ and close to it, to form a rectangle $R$. The unstable branch $\Ga^u(q)$ that  accumulates at $x$ is contained in the disk so it cannot intersect the boundaries $\beta_1, \beta_2$ and given that $p$ is not in $R$, its unstable branch intersects  $\gamma_2$; now, one can take all the connected components $\Ga^u(q)$ intersected with $R$ and   some of those connected components  accumulate on $x.$
    
\end{proof}

\begin{figure}[h]
\centering
\includegraphics[scale=0.75]{fig_rectangle.pdf}
\caption{Rectangle in lemma \ref{rectangle}}
\label{rectangle-fig}
\end{figure}

Now we introduce a definition that extends one provided in \cite{CPT} (see definition 5.9 in that paper).

\begin{definition} Given a cycle  of periodic points and quadrants  $(q_i, \Ga^s_l(q_i), \Ga^u_l(q_i))$ such that $\Ga^u(q_i)$ does not intersect  $\Ga^s(q_{i+1})$, a {\it Pixton disk} (see figure \ref{fig_pixton}) associated with that sequence  is a topological disk $D$ bounded by a collection of arcs $(\ga^s_i, \ga^u_i, \de^i)$ such that

\begin{itemize}
   \item[--] $\ga^u_i$ is an arc in $\Ga^u(q_i)$, 


   \item[--] $\ga^s_i$ is an arc in $W^s_\DD(q_{i})$ with one endpoint is $q_i$ and another endpoint is a point $x_i$ that is accumulated by $\Ga^u(q_{i-1}),$

   \item[--] $\de_i$ is an arc that joins $\ga^u_i$ and a point $x_{i+1}$ in $\ga^s_{i+1}$ satisfying that $f(\de_i)\cap \de_i=\emptyset$ (it could be the case that $\de_i$ is empty and in that case $x_{i+1}\in \ga^u_i\cap \ga^s_{i+1}$).
\end{itemize}

    
\end{definition}

\begin{figure}[h]
\centering
\includegraphics[scale=0.75]{fig_pixton.pdf}
\caption{Pixton disk involving three saddles.}
\label{fig_pixton}
\end{figure}
A version of the next  lemma  is stated in \cite{CPT} and its proof follows from lemma 5.10. For completeness (and because the assumption in that lemma are not quite similar) we include the proof.





\begin{lemma}\label{rmk:pixton} Given a  Pixton disk, there  is a point $x$ in one of the unstable branches such that  its forward iterates are not in the interior of the Pixton disk and therefore neither  $\omega(x)$.
    
\end{lemma}


 \begin{proof}
     
The proof follows from showing that $K=\alpha(x_i)$ is disjoint from the interior of the Pixton disk. In fact, if for some $x_i$ it holds that its $\alpha$-limit is inside $D$, let $f^{-n_k}(x_i)\in D$. From the mild dissipativeness and the fact that the unstable branches of the points in the cycle  do not intersect the stable branches, it follows that $W^s_\DD(f^{-n_k}(x_i))$ intersects some $\de_j$. Since $f(\de_j)\subset D$ (provided that $\de_j$ is small) it follows that $f(W^s_\DD(f^{-n_k}(x_i)))\cap D\neq\emptyset$ and since $f(W^s_\DD(f^{-n_k}(x_i)))\subset W^s_\DD(f^{-n_k+1}(x_i))$ it follows again that $W^s_\DD(f^{-n_k+1}(x_i))\cap \de\neq \emptyset.$ Inductively, it follows that $W^s_\DD(x_i)\cap \de_j\neq \emptyset,$ a contradiction.
\end{proof}

The next proposition is a preparatory one to prove proposition \ref{existence of touching} and states that under the hypothesis of theorem \ref{thm:finite classes} any unstable branch of a saddle periodic point that is not contained  in the basin of attraction of a (semi)attracting periodic point has to cross the stable of another saddle. The proof uses the dichotomy of  theorem \ref{disk2} and that if the critical points accumulated by the unstable branch of the saddle are not interior critical points, then there is a cycle involving boundary saddle periodic points and therefore a Pixton disk as defined before. 

\begin{proposition} \label{prop trans} Under the hypothesis of theorem \ref{thm:finite classes},  given a saddle periodic point $p$  and one of its unstable branches $\Ga^u(p)$, it holds  

\begin{itemize}
  \item[--] either it  crosses the stable manifold   of a saddle periodic point or, 

 \item[--] there is a (semi)attracting periodic point such $\Ga^u(p)$   intersects the interior of its  basin of attraction.
    
\end{itemize}



    
\end{proposition}


\begin{proof} 

 If  the closure of the unstable branch of $p$ does not have critical points then by theorem \ref{disk2} either it is contained in the basin of attraction of a (semi)attracting periodic point or it is a non-trivial normally hyperbolic attracting interval; in the last case the unstable branch crosses the stable manifold of a saddle in the normally hyperbolic arc. 

If it does have  critical points $c$  then $\omega(c)$ has dominated splitting so either it contains interior points, or is a heterocycle, or it is a periodic point.  In the first two  cases the proposition is proved.  In the last case, it has to be a boundary saddle periodic point. 

So, it remains to consider the case that all the critical points in the closure of the unstable branch are  boundary critical points. In that case, recall that the boundary critical points  belong to a finite number of boundary saddle points, named $\{q_1
\dots q_n\}$. 

\emph{Iterating the dichotomy.} For each boundary saddle point $q_i$, one considers its unstable branch  $\Ga^u(q_i)$ that is accumulated by the unstable branch of $p$ , and apply the same dichotomies as before: if  some  branch $\Gamma^u(q_i)$  does not have critical points, again by theorem \ref{disk2}  the proposition is concluded. So, it remains to consider the case where for all $q_i$ the set $\La^u(q_i)$ (the closure of any $\Ga^u(q_i)$ that is accumulated by $\Ga^u(p)$) has critical points that all are boundary critical points.



\emph{Creating cycles.} Reasoning as  before, for each critical  boundary point $c_i$  there is  a critical cycle (named ${\mathcal C}_i$) of saddle boundary points such that there is no transversal homoclinic intersection between the unstable and stable branches of those saddle points (otherwise the proposition is also proved); moreover, by remark \ref{enlarging} one can assume that the cycle ${\cal C}_i$  is maximal. If the cycle is a touching cycle, by proposition \ref{touching} the proposition  is concluded. 

\emph{Closing into Pixton disks and escaping from them.} It remains to  consider the case that all the  maximal cycles ${\cal C}_i$ are not  touching cycles. For each ${\cal C}_i$ one produces a  Pixton disk $D_{j_i}$. By lemma \ref{rmk:pixton} it follows that there is $x_{j_i}$  in the unstable branch of one of the saddles that bounds the Pixton disk $D_{j_i}$ such that $\omega(x_{j_i})$ is disjoint with  the interior of the Pixton disk $D_{j_i}$. If for  some $x_{j_i}$ it follows that $\omega(x_{j_i})$ does not have a critical point then, arguing as before, $\omega(x_{j_i})$ is a periodic orbit $O(\hat q)$ contained in $\La^u(p)$ and  $x_{j_i}\in W^s_\DD(\hat q)$. 

Now we have a simple dichotomy: either i) $\hat q\notin \{q_1\dots q_n\}$ or ii) $\hat q\in \{q_1\dots q_n\}$. In the first case, the unstable branch accumulated by $\Ga^u(p)$ does not have critical points and therefore, $\La^u(\hat q)$ is either a (semi)attracting periodic point or a normally hyperbolic arc and in both cases the proposition is proved.  In the second case,  the unstable branches of the saddles in the cycle ${\cal C}_{j_i}$  accumulate on  another Pixton disk $D_{j_i'}$. One can repeat the argument with the Pixton disk $D_{j_i'}$ and observe that the proposition is either proved  or the unstable branches of the periodic points in the cycle ${\cal C}_{j_i'}$ involved in the  Pixton disk has to accumulate on another Pixton disk;   since the cycle in each Pixton disk is maximal there cannot be a cycle between Pixton disk, then it follows that for some $x_{j_i}$ its $\omega(x_{j_i})$ does not have critical points and the proposition is concluded.






\end{proof}



\paragraph{Proof of proposition \ref{existence of touching}.} 

Let us consider two saddle $q_0$ and $q_1$ involved in the maximal critical cycle and let $c_1$ be the critical point in  $\Gamma^s(q_1)$ (one of the branches of $q_1$) that is accumulated by $\Gamma^u(q_0)$ and the orbits of the $*-$sequence $(p_n)$. The goal is to show that either $\Gamma^u(q_0)$ intersects $\Gamma^s(q_1)$ at $c_1$ or there is another saddle  periodic point $\bar q$ such that $\Gamma^u(q_0)$ intersects $\Gamma^s(\bar q)$ and  $\Gamma^u(\bar q)$ intersects $\Gamma^s(q_1).$ Applying this argument to each consecutive pair in the cycle, one concludes that the cycle is either a critical touching cycle or can be enlarged to a touching cycle. 

Since $\al(c_1)$ has dominated splitting and since the orbit of the $*-$sequence $(p_n)$ also accumulates at $\alpha(c_1)$ it follows that $\al(c_1)$ is a saddle periodic point named $\bar q$.  If $\bar q$ is equal to $q_0$ then $\Ga^u(q_0)$ intersects $\Ga^s(q_1)$. If $\bar q\neq  q_0$, observing that $\Ga^u(q_0)$ and the orbit of the $*-$sequence $(p_n)$ also accumulate on $\bar q$,  there is a point $\bar x$ in $\Ga^s_\DD(\bar q)$ (one of the local stable branches of $\bar q$) that is accumulated by $\Ga^u(p_0)$ and the orbits of $(p_n)$ along a same quadrant of $\bar q.$ If $\bar x$ is a critical point, it follows that adding $\bar q$ to the cycle one gets a larger critical cycle contradicting the fact that the initial critical cycle was  maximal.

To continue, one has to consider the case that $\bar x$ is not a critical point, and now we consider different possibilities related to $\al(\bar x)$: either it does not contain a critical point or it has it. Similarly, as before, the goal is to prove that $\Ga^u(q_0)$ intersects $\Ga^s(\bar q)$ and in that way we enlarge the cycle adding a periodic point such that now that part of the cycle is touching.

If $\al(\bar x)$ does not have  a critical point then it has a dominated splitting and by the fact that the  pre-iterates of the sequence $(p_n)$ accumulate on $\al(\bar x)$ it follows by the definition of $*-$sequence that $\alpha(\bar x)$ is a saddle periodic orbit (that one can assume it is fixed), named $\hat q$. In that case, $\bar x$ belongs to $\Ga^u(\hat q)$. The intersection of the unstable manifold of $\hat q$ and $\bar q$ at $\bar x$ is transversal, otherwise by lemma \ref{lem:tang-crit} there is a point in the orbit of $\bar x$ that is a critical point; since the $*-$sequence and $\Ga^u(q_0)$ also accumulate on the orbit of $\bar x$ we get a contradiction with the assumption that the cycle is maximal.   Since $\Ga^u(q_0)$ accumulates at $\bar x$, it also accumulates at $\hat q$; let $\hat Q$ be the quadrant of $\hat q$ that contains points of $\Ga^u(q_0)$ accumulating on $\hat q.$ Let $\Ga^s_\DD(\hat q)$ be the local stable branch of $\hat q$ that bounds $\hat Q$, and observe that  is accumulated by  connected arcs $\bar \ga_n$ contained in $\Ga^s(\bar q)$ that one extreme points are pre-iterates of $\bar x$ and the others are in the boundary of the disk. Then one gets a sequence of rectangles bounded by $\Ga^s_\DD(\hat q)$, $\bar \ga_n$, arcs in $\Ga^u(\hat q)$ and arcs in the boundary of the disk. These rectangles  contain in the interior iterates of $\Ga^u(q_0)$, and therefore it has to cross  the arcs $\bar \ga_n$ a so $\Ga^u(q_0)$ intersects $\Ga^s(\bar q)$ as we wanted.



It remains to consider the case that  $\al(\bar x)$  has critical points and let $\hat c$ be one of that critical points; since $\hat c$ is also accumulated by $\Ga^u(q_0)$ and pre-iterates of the sequence $(p_n)$ it follows that it is a boundary critical point and so there is a boundary periodic point $\hat q$,  that one can assume that it is fixed,  and $\hat c \in \Ga^s_\DD(\hat q)$.   
By proposition \ref{prop trans} (and the fact that $\Ga^u(\hat q)$ cannot be fully contained in the basin of attraction of a (semi)attracting periodic points since that branch is accumulated by periodic orbits), either there is a saddle $\tilde q$ such that $\Ga^u(\hat q)$ (the unstable branch of $\hat q$ that it is accumulated by the pre-iterates of $\bar x$) that crosses $\Ga^s(\tilde q)$ or there is a (semi)attracting periodic point (also denoted as $\tilde q$) such that the interior of its basin intersects  $\Ga^u(\hat q)$; in both cases, one gets an arc $\tilde \ga$ in  $W^s_\DD(\tilde q)$  that cross $\Ga^u(\hat q).$  Now, one  gets  that $\Ga^s_\DD(\hat q)$   is accumulated by  connected arcs $\tilde \ga_n$ obtained as pre-iterates of arcs in  $\tilde \ga$ whose endpoints are in $\Ga^u(\hat q)$ and in the boundary of the disk. Moreover, these arcs intersect the quadrant $\hat Q$ that contains the pre-iterates of $\bar x$ that accumulate on $\hat c$. Then one gets a sequence of rectangles $R_n$ bounded by $\hat \ga$, $\tilde \ga_n$, arcs in $\Ga^u(\hat q)$ and arcs in the boundary of the disk, that contains pre-iterates of $\bar x$. Let us pick $f^{-k}(\bar x)$ close to $\hat c$ and let $\bar \ga_k\subset W^s_\DD(f^{-k}(\bar x))$ be an arc in $\Ga^s(\bar q)$ that one end point is $\bar q$ and the other is $f^{-k}(x).$ Provided that $n$ and $k$ are  large, the   arc $\bar \ga_k$ has to exit $R_n$ and given that it cannot intersect neither $\hat \ga$ nor $\tilde \ga_n$ then it has to intersect either $\Ga^u(\hat q)$ or the boundary of the disk. In both cases, either pre-iterating or post-iterating $\bar \ga_k$,  there is a connected arc $\bar \ga_k'$ in $\Ga^s(\bar q)$ inside $R_n$ such that one extreme point close to $f(\hat c)$ and the other close to $f^{-1}(\hat c).$ On the other hand let us consider a rectangle $R$ as in lemma \ref{rectangle}, where the arcs $\si_m$ are contained in $\Ga^u(q_0).$ Taking $R\cap R_n$ observe that both arcs $\be_1$ and $\be_2$ are crossed by  $\bar \ga_k'$ and therefore, the arcs $\si_n$ has to cross $\bar \ga_k'$, showing again that $\Ga^u(q_0)$ intersects transversely $\Ga^s(\bar q)$ as we wanted  (see figure \ref{final} for the last configuration).

\qed 


\begin{figure}[h]
\centering
\includegraphics[scale=0.75]{fig_touching_final.pdf}
\caption{Final of the proof of proposition \ref{existence of touching} }
\label{final}
\end{figure}

That concludes the proof of theorem \ref{thm:finite classes}.






% \section{Twist maps, dominated splitting and critical points.}


% Let us define for any $x$ and any positive integers $n$ the Under the assumption that the directions $G^\pm(x)$ are well defined, 
% since  $\La$ is compact, it follows that $G^+_n(x)\to G^+(x)$ and $G^-_n(x)\to G^-(x)$ uniformly. Therefore, the directions $x\to G^\pm(x)$ moves continuously. 

% \vskip 5pt



% \noi{\it Proof of theorem \ref{green domination}:} If there is a dominated splitting $E\oplus F$, it follows immediately that $f_{|\La}$ is twist: it is enough to consider the diagonal bundle $v=E+F$. 

% For the converse, if one assume that there is $\alpha>0$ such that $slope(G^+(x), G^-(x))>\alpha$ for any $x$, one can do a uniform change of metric in such a way that $\{G^-(x), G^+(x)\}_{x\in \La}$ are orthonormal basis, denoted as $\{e_1, e_2\}$ independent of the base point $x$; under this change of coordinates, $Df$ is a diagonal matrix. Without loss of generality, one can assume that the vertical direction is given by the vector $v=-e_1+ e_2$, and since $slope(G(v), v)> \beta$ it follows that $slope(G(v), e_2)<\lambda.slope(v,e_2)$ for some $0<\lambda<1$; this immediately implies that $\frac{|Df_{|G^-}|}{|Df_{|G^+}|}< 1,$ and so the splitting $G^-\oplus G^+$ is a dominated splitting.

% \qed




% %  associated with the vertical direction, we define $H^+_x$ as the connected component in $T_xM\setminus<v^+_x>$ that contains $Df_{f^{-1}x}(v^+).$


% % \begin{remark}
% % For any $x$ the subbundles $G^-_x$  and $G^+_x$ belong to the  interior of $H^+_{f(x)}$ with slope larger than $\al_0$ with $v^+.$ 
% % \end{remark}
% % \begin{claim}\label{change met} Let us assume that $f_{|\La}$ is twist and for any $x$ holds that $\angle(G^-_x, G^+_x)\geq \be_0$. Then there exists a Riemannian metric $m_2$, $K_0=\{\al_0, \be_0\}$ and $\bar \al_0\geq \min\{\al_0, \pi/4\}$ such that 
% % \begin{enumerate}
% %  \item $m_1(v,w) K_0^{-1}< m_2(v,w) <m_1(v,w) K_0$ where $m_1$ is the initial Riemannian metric;
% % \item $G^-_x\perp_{m_2}G^+_x$;
% % \item $||Df(u)||_1=||Df(u)||_2$ where $u=G^-, G^+$; 
% % \item $\angle_2(v^+, G^+)\geq \bar \al_0, \angle_2(v^+, G^-)\geq \bar \al_0;$ 
% % % \item $sl_1(u, G^\pm) K_0^{-1}< sl_1(u, G^\pm) < sl_1(u, G^\pm) K_0$ for any $u$ such that $\angle_2(v^+, F)\geq \bar \al_0.$
% % \item $K_0^{-1} \leq  g^k_2(G^\pm_x)\leq K_0 g^k_1(G^\pm_x).$
% % \end{enumerate}

 
% % \end{claim}



% % \noi{Proof of theorem \ref{green domination}:} Let us assume that for any $x$ follows that $G^-_x$ and $G^+_x$ are transversal directions; so there exists $\be_0$ such that $\angle(G^-_x, G^+_x)\geq \be_0$ for any $x.$
% % We consider now the metric given by  claim \ref{change met} and we can write the $Df^k_x$ as a diagonal basis using the  orthogonal basis $\{G^-_x, G^+_x\}, \{G^-_{f^k(x)}, G^+_{f^k(x)}\}.$
% % In that basis, 
% % $$slope(G^k(v^+), G^+_{f^k(x)})= slope(v^+, G^+_x)\frac{|Df^k(G^-)|}{|Df^k(G^+)|}=slope(v^+, G^+_x)g^k_2(G^+_x);$$
% % since $slope(G^k(v^+), G^+_{f^k(x)})\to 0$ and $slope(v^+, G^+_x)\geq \bar\al_0$ it follows that given $\ga<1$ there is  $k$ large,  such for any $x$ holds $g^k_2(G^+_x)<\ga$.
% % Now, let us consider the direction $v^-= {v^+}^{\perp}$ and since $G^+_x$ is a contractive direction in the projective tangent bundle it follows that $slope(G^+_{f^k(x)}, G^k(v^-))\to 0$ and therefore, the cone  $C$ bounded by $v^+$ and $v^-$ is uiformly contracted and therfore, thee is a dominated splitting. 



% % \qed


% The following lemma is an immediate application, in the context of one-dimensional projective dynamics, of a standard lemma in one-dimensional dynamics.

% \begin{lemma}\label{uniform basin angle}
% Let $\lambda <1$ and $v$ such that $|g^k(v)|<\la^{k}$ for $k=1\dots n$ then there exists $r_0$ (that only depends on the upper bound of the norm of the cocycle) such that  if $slope(v, u)< r_0$ then $slope(G^k(v),G^k(u))< \la^{k} slope(v,u).$
% \end{lemma}




% \begin{lemma}\label{g on G}
% For any twist diffeomorphisms, it follows that $|g^k(G^+(x))|\leq 1.$
% \end{lemma}

% \begin{proof} It follows immediately from the fact that for $n$  and $k$ sufficiently large,  it holds that $slope(G^k(G^+_n(x)), G^k(G^+(x)))\to 0,$ meaning that vectors on one side of $G^+(x)$ are attracted to its forward  iterates.
 
% \end{proof}



% \noi{\em Proof of theorem \ref{green critical}:} Let us assume that $slope(v,G^+_x)> \ga_0$. By lemma \ref{gn some growth}, there is $\la<1$ and $n$ arbitrarily large such that  $g^{-n}(v)>\la^{-n}$ and so $g^n(G^{-n}(v))< \la^n$; by  Pliss's lemma, there is $n_0>0$ such that $g^k(G^{-n+n_0}(v))<\la^{k}$ for $k=1\dots n-n_0$ and $n-n_0$ is arbitrarily large provided that $n$ is chosen sufficiently large.  To avoid notation, let us denote $m_0:=-n+n_0$ and $v_{m_0}:= G^{n_0-n}(v).$ By  lemma \ref{uniform basin angle} follows that  $slope(v_{m_0}, G^+_{f^{n_0-n}(x)}))>r_0$. So, taking the basis $\{v_{m_0}, G^+_{f^{-n+n_0}(x)}\} $ and $\{v, G^+_{x}\} $ by a uniform change in metric such that both bases are orthogonal, we can assume that $Df^{n-n_0}$ is a diagonal matrix  and since $g^{n-n_0}(v_{m_0})< \la^{n-n_0}$ it follows that $g^{n_-n_0}(G^+_{f^{n_0-n}(x)})>\la^{-(n-n_0)}$ contradicts the fact that $g^k(G^+(z))\leq 1$ for any point $z$ (see lemma \ref{g on G}). Similarly, we prove that $G^-_x$ coincide with $v.$
 
% \qed

% \section{Minimal dynamics in the annulus.}
% \label{sec:annulus}

% It has been proved elsewhere
%  the following

%  \begin{theorem*}\label{twist} Let  $f\in \operatorname{Diff}^{1+}(\AA)$, let us assume that  $f_{\AA}$ is a twist diffeomorphism and  $\La$ is  a  minimal  aperiodic  attractor. Then $\La$ is a Lipschitz graph respects to the vertical, i.e. there is a Lipschitz function $\ga: S^1\to \RR$ such that $\La=graph(\ga)$.


% \end{theorem*}

% The next lemma is a simple observation that relates how the Lipschitz cone of a curve is related to the Green's bundle. Recall that given a Lipschitz curve $\gamma$, one can define the Lipschitz cone at a point $x$ (denoted as $C_x$) as the minimal  cone in the tangent bundle that contains all the vectors obtained as limits of $\frac{1}{h}( \gamma(t+h)-\gamma(t)),$ where $\gamma(t)=x.$

% \begin{lemma}\label{curve-crit}Let  $f\in \operatorname{Diff}^{1+}(\AA)$, let us assume that  $f_{\AA}$ is twist let $\La$ is  a  minimal  aperiodic  attractor. Then there are not conjugate points in $\La$ and for any  $x$ holds that $$C_x\subset [G^-_x,G^+_x].$$
% \end{lemma}


% \begin{proof} The fact that the upper bound of $C_x$ is bounded by $G^+_x$ follows immediately form the following: given  two  nearby points $x, x'$  in the curve $\gamma$,  and  letting  $x_{-n}=f^{-n}(x), x_{-n}'=f^{-n}(x')$, if the points are close enough it follows that on the one hand, the vector $v_{-n}=x_{-n}'-x_{-n}$ is to the right to the vertical direction and so $G^n(v_{-n})$ is in the cone bounded by $G^n(v^+)= G_n^+(x)$ and $G^+_x$; and on the other hand $G^n(v_{-n})$ is close to the vector $v=x'-x.$   In a similar way iterating forwardly, one prove that the lower bound of $C_x$ is bounded by $G^-_x.$

    
% \end{proof}

%  Observe that it follows immediately that if $G^+_x=G^-_x$ then the curve $\gamma$ is differentiable at $x$. However, the curve can be differentiable even in the case that $G^-_x\neq G^+_x$; in fact, if $\gamma$ is normally hyperbolic, it follows on the one hand that $\gamma$ is smooth and, on the other hand, by theorem \ref{green domination}, $G^\pm$ coincides with the invariant splitting at the curve.
 
%  %Moreover, there is $\beta:= \beta(\al_0)$ such that 
% % $sl(v^+, C_x)>\beta$ and if $C_x$ is the tangent cone to $\ga$ and   $C_x$ does not have interior, then $\ga$ is differentiable at $x.$
 


% \paragraph{\em Proof of theorem \ref{thm:annulus}:} It follows immediately from the theorems \ref{green critical} and lemma \ref{curve-crit}. 
 
% \qed




% % \begin{lemma}\label{vert pliss} Let $c$ be  a critical point and let $m^-_i, m^+_i$ be the sequences given by cor \ref{crit pliss}.  Then $$slope(G^{m^-_i}(v^+), v)< \la^{m^-_i},$$  $$slope(G^{-m^+_i}(v^+), v)< \la^{m^+_i}.$$  
 

% % \end{lemma}


% % \begin{remark}
% % For any $y$ the bundle $E^s_y$ belong to the positive half $H^+_y$ bounded by $v^+$ and that $G(v^+)$ is in the interior of $H^+_{f(y)}$ with slope larger than $\al_0$ with $v^+.$ Also, $sl(v^+, E^s_y)>\al_0.$
% % \end{remark}
% % % 

% % %Observe that in the new metric, it follows that $g^k_2(E^s_x)=g^k_1(E^s_x).$


% % % 
% % % \begin{definition}
% % % Given $A\in GL(2, \RR)$ let $e, f$ be the  maximally contracting and expanding  directions.  \end{definition}
% % % The maximally contracting and expanding  directions  are the solutions of $$\partial_\theta||A(\cos(\theta), \sin(\theta))||=0$$ which are given by
% % % $$\tan(2\theta)=\frac{2(ab+cd)}{a^2+c^2-(b^2+d^2)}$$
% % % where $$A=\left(\begin{array}{ccccc} a & b \\  c & d \end{array}
% % % \right).$$
% % % It follows that $e\perp f$, $A(e)\perp A(f)$,  $g(f)= \frac{1}{g(e)}$ and $\frac{1}{g(e)}\leq g(w)\leq g(e),$  for any $w.$ In particular, in the basis $\{e, f\}$ 
% % % $$A=\left(\begin{array}{ccccc} a
% % %  & 0 \\  0 & d \end{array}
% % % \right)$$ and $g(e)=\frac{d}{a}, g(f)=\frac{a}{d}.$
% % % 
% % % Let $\{A_i\}$ be a sequences of matrixes in $GL(2,\RR)$  such that $||A_i||<\s$ for a fixed $\s$. Then it follows that if $f_n$ is the maximally expanding direction then 
% % % $g^n(f_n)\geq \frac{b}{\s}^n;$ in fact since $g^n(f_n)=  \frac{a_n}{d_n}$ with $a_nd_n < b_n$ and $d_n < \s^n.$


% % \begin{lemma}\label{uniform angle} Given $b_0$ and $\al_0$, there exists $\la<1$ such that for any $0< b< b_0$ and $\be>0$ if $sl(E^s, F)>\be_0$ then  there exists $n_0= n_0(b, \be_0)$ such that 
% % $$|Df^{n_0}_{E^s}|< \la$$ and $$g^{n_0}(E^s)> 1+\de.$$
% % \end{lemma}

% % \begin{proof} Using the previous change of metric, it follows that
% % $$sl_1(G^k(E^s), G^k(v^+))< sl_2(G^k(E^s), G^k(v^+))K_0< K_0 \frac{b}{\la}^ksl_2(E^s, v^+),
% % <K_0^2 \frac{b}{\la}^ksl_1(E^s, v^+),$$
% % so taking $k$ such that $$K_0^2 \frac{b}{\la}^k<1,$$
% % we get a contradiction.


% % \end{proof}

% % \noindent{\em Proof of theorem \ref{goes to tangency}.} If the thesis of the theorem does not hold, by lemma \ref{uniform angle} we get that for all $n$ follows that $|{Df_n^{n_0}}_{/E_n^s}|< \la$ and $g^{n_0}(E^s_n)>1+\de.$ Now, for any $x\in \La_f$ we take a $E_x$ and accumulation point of $E_n^s$ and it follows that $g^{n_0}(E_x)>1+\de.$ So, we get a uniform cone field and therefore $\La$ has a dominated splitting.

% % \qed

% % \marginpar{Can the  Hall's construction be adapted to get an example such that $g^n$ does not grow?}

\section{Further questions}\label{sec:questions}

 
 The critical set is a new
object, and several natural questions about it are left open by this paper.

\paragraph{Is ${\rm Crit}(f,\La)$ totally disconnected?} In all the examples of
section \ref{examples} the critical set is at most a Cantor set (for the
Wang-Young attractor (paragraph iv-) it is explicitly the intersection of a
nested sequence of rectangles) and none exhibits a critical arc. When
$f_{|\La}$ is far from homotheties the critical direction is unique and
varies continuously, which constrains the set; we do not know whether this,
or a dissipation hypothesis, forces total disconnectedness, nor whether a
continuum of critical points can occur in general.

\paragraph{When is a critical point strong?} Following \cite{CLPY} (paragraph
v-), one may strengthen definition \ref{def:crit} by requiring exponential
growth of the projective cocycle along the critical direction in both time
directions, $|g^n(v)|>(1+\de)^{|n|}$ for all $n$ and some $\de>0$. The
strong notion carries genuine geometry: at such a point there are a local
stable and a local center-stable manifold, tangent at the critical point, so
that the tangency picture of paragraph i- is recovered as actual invariant
manifolds rather than merely as a condition on the cocycle. For critical
points coming from tangencies this strengthening holds automatically
(paragraph i-), and proposition \ref{exponential recurrence} controls the
recurrent case; it would be of interest to identify general conditions,
on the recurrence of the critical set or on the dynamics of $\La$, under
which every critical point is strong, since, as \cite{CLPY} shows, the strong
notion carries substantial structural information.

\paragraph{Critical-orbit conditions and SRB measures.} More speculatively, the
critical set may support, for surface diffeomorphisms, the program that in
one dimension extracts ergodic information from conditions on the critical
orbit: a Collet--Eckmann condition, or summability and slow-recurrence
conditions in the spirit of Misiurewicz, yield absolutely continuous
invariant measures and, more generally, SRB measures. The Misiurewicz
condition of section \ref{sec:misiu} is the non-recurrent extreme of such a
hierarchy; one may ask whether intermediate conditions, call it a two-dimensional
Collet--Eckmann condition along the critical directions, or controlled
recurrence of the critical set, yield SRB measures, as they do in
dimension one and as they do, by different methods, in the work of
Benedicks--Carleson and Wang--Young (paragraph iv-).

\begin{thebibliography}{HPSJ}

{\footnotesize \setlength{\parskip}{.001in}

% % \bibitem[Ar]{Ar}M.-C. Arnaud. 
%  Three results on the regularity of the curves that are
% invariant by an exact symplectic twist map. {\em Publ. Math. Inst. Hautes Études Sci.} No. 109 (2009), 1-17.


\bibitem[BC]{BC} M. Benedicks, L. Carleson. The dynamics of the H\'enon map. {\em Ann. Math.} 133, 73–169 (1991).



\bibitem[BS]{BS} E. Bedford, J. Smillie;
Real polynomial diffeomorphisms with maximal entropy: Tangencies.
{\em Ann. of Math. (2) 160(2004)}, no.1, 1–26.

\bibitem[BS2]{BS2}  E. Bedford, J. Smillie;
Real polynomial diffeomorphisms with maximal entropy. II. Small Jacobian.   {\it Ergodic Theory Dynam. Systems 26(2006)}, no.5, 1259–1283.

\bibitem[BoS]{BoS}
J. Boronski and S. Stimac.
{\it The pruning front conjecture, folding patterns and classification of Hénon maps in the presence of strange attractors.}
ArXiv:2302.12568



\bibitem[BCS]{BCS} J. Buzzi, S. Crovisier, O. Sarig; 
Measures of maximal entropy for surface diffeomorphisms. {\em Annals of Math.} 195 (2022), 421-508.


% \bibitem[C]{C}S. Crovisier. Ensembles de torsion nulle des
%  applications d\'eviant la verticale. {\em Bull. Soc. math. France}
% {\bf 131} (1), 2003, p. 23-39.

\bibitem[CP]{CP} S. Crovisier, E. Pujals; Strongly dissipative surface diffeomorphisms. Comment. Math. Helv., 93 (2018), 377–400.

% \bibitem[CKP]{CKP} S. Crovisier, A. Kocsard, E. Pujals; Strongly dissipative surface diffeomorphisms in the annulus. 


\bibitem[CLPY]{CLPY}S. Crovisier, M. Lyubich, E. Pujals, J. Yang; Renormalization of Unicritical Diffeomorphisms of the Disk, {\emph arxiv:2401.13559.}



\bibitem[CPT]{CPT} S. Crovisier, E. Pujals and C. Tresser.
Strongly dissipative diffeomorphisms of the disc with zero entropy. \emph{Acta Mathematica}, Volume 232 (2024) Number 2, 221-323.


\bibitem[E]{E} H. Enrich, 
A heteroclinic bifurcation of Anosov diffeomorphisms.
\emph{Ergodic Theory Dynam. Systems}, {\bf 18} (1998), no.3, 567-608.

\bibitem[HMS]{HMS} V. Horita, N. Muniz and P.R.  Sabini;
Non-periodic bifurcations for surface diffeomorphisms.
\emph{Trans. Amer. Math. Soc.} {\bf 367} (2015), no.12, 8279-8300.

\bibitem[M]{M} R. Ma{\~n}{\'e}, \emph{Hyperbolicity, sinks and measure in one-dimensional
dynamics}, Commun. Math. Phys., {\bf 100} (1985), 495--524.

\bibitem[N]{N} S. Newhouse, Non-density of Axiom A(a) on S2. {\em Proc. Am. Math. Soc. Symp. Pure Math.} 14,
191–202 (1970).


% \bibitem[Pi]{Pi} D. Pixton, 
% Planar homoclinic points.
% {\em J. Differential Equations 44 (1982)}, no. 3, 365--382.

\bibitem[PRH]{PRH} E. Pujals, F. Rodriguez Hertz. Critical points for surfaces maps, {\em
Journal of Modern Dynamics} Volume 1, No. 4, 2007.

\bibitem[PS1]{PS1} E. Pujals, M. Sambarino, Homoclinic
tangencies and hyperbolicity for surface diffeomorphisms, {\em
Annals of Math,  151, (2000)}, 961-1023.
\bibitem[PS2]{PS2} E. Pujals, M. Sambarino, On the dynamics of dominated splitting, {\em
Annals of Math,  169, (2009)}, 675-704.


\bibitem[R]{R} I. Lug\~ao Rios, 
Unfolding homoclinic tangencies inside horseshoes: hyperbolicity, fractal dimensions and persistent tangencies,
{\it Nonlinearity 14(2001)}, no.3, 431-462.


\bibitem[T]{T}  H. Takahasi,
Prevalent dynamics at the first bifurcation of Hénon-like families.
{\it Comm. Math. Phys.312 (2012)}, no.1, 37–85.

\bibitem[V]{V} F. Valenzuela-Henr\'quez   
 On critical point for two-dimensional holomorphic systems. {\em Ergodic theory and dynamical systems} (37),  2017 , 2276-2312


\bibitem[WY]{WY} Q. Wang, L-S. Young, Toward a theory of rank one attractors. {\em Ann. of Math.} (2) 167
(2008), no.2, 349–480.

}
\end{thebibliography}


\begin{tabular}{l l l}
\emph{Sylvain Crovisier} &&
\emph{Enrique Pujals}\\

Laboratoire de Math\'ematiques d'Orsay &&
Graduate Center \\

CNRS - Univ. Paris-Saclay &&
CUNY\\

Orsay, France &&
New York, USA\\

&&\\


\end{tabular}

\end{document}









\paragraph{\it Twist and conjugate to horseshoes}

Let $f_0$ be a surface diffeomorphisms such that $\La_0=\cap_{n\in \ZZ}f_0^n(U)$ is a hyperbolic invariant set.  Let $f$ be a dissipative diffeomorphisms in $U$ such that  $\La=\cap_{n\in \ZZ}f^n(U)$ is conjugate to $\La_0$; moreover, we assume that the conjugacy holds in $U$, i.e. there exists a homeomoprhisms $h: U\to U$ such that $h\circ f_0=f\circ h.$

Let $f\in Diff^1(M)$, an open neghborhoof $U$ not necessary invariant such that $\La=\cap_{n\in \ZZ}f^n(U)$ is a compact invariant set inside $U$. Let us assume that $f_{/U}$ is twist. We say that there is not conjugate points in $U$, if for any $x\in U$ if 
$n_x^-=\sup\{n >0 f^{-j}(x)\in U, j<n\}$   and  $n_x^+=\inf\{n >0 f^{j}(x)\in U, j<n\}$ and
$$\tilde G^+(x)=\lim_{n\to n^-_x} G^n(v^+),\,\,\,\,\tilde G^-(x)=\lim_{n\to n^+_x} G^{-n}(v^+)$$ it  follows that $$\tilde G^-(x)\leq \tilde G^+(x).$$





\begin{theorem} \label{tang-horseshoes} Let $f_{\La}$ be a twist diffeomorphisms such that 
\begin{enumerate}
 \item $\La$ is conjugate to to a hyperbolic set in a neighborhood;
\item $\La$ is not a hyperbolic set and the periodic points are hyperbolic;
\item $f$ has  not conjugate points in $U;$ 
\end{enumerate}
 then there is a tangency; i.e. there is a point $x\in \La$ such that its  local stable and unstable set are differentiable at $x$ and its tangent spaces coincide.
 
\end{theorem}

\begin{proof} It follows from the simple fact that at the critical point $G^+(x)=G^-(x)$ and 
the unstable and stable sets are Lipschitz graph such that  each point $z$ the Lipschitz cone is contained  in $[\tilde G^-(z), \tilde G^+(z)].$

\end{proof}


Let us define $D^{++}_x= \cap_{n\geq 0} f^n(H^+_{f^{-n}x}\cap U)$ and $D^{+-}_x= \cap_{n\geq 0} f^{-n}(H^+_{f^{n}x}\cap U)$. In the same way we define $D^{-+}_x= \cap_{n\geq 0} f^n(H^-_{f^{-n}x}\cap U)$ and $D^{--}_x= \cap_{n\geq 0} f^{-n}(H^-_{f^{n}x}\cap U)$.

\begin{lemma} The sets $D^+_x=D^{++}_x\cap D^{+-}_x$ and $D^-_x=D^{-+}_x\cap D^{--}_x$ are  not empty and contain interior.
 
\end{lemma}

\begin{lemma} There are  one branch of the local unstable and local stable set contained in $D^+_x$ and the others  branches of the local unstable and local stable set are contained in $D^-_x.$
 
\end{lemma}


% 
% \begin{theorem}\label{green-crit-horseshoes} The critical points coincides with the points where $G^+=G^-.$
% \end{theorem}

\paragraph{\it Twist and homotopic to hyperbolic}





Let $\{f_t\}_{t\in [0,1]}$ be a one parameter family of surface diffeomorphisms such that

\begin{enumerate}
\item  for any $t\in [0,1]$ there is a maximal invariant set $\La_t=\cap_{n\in \ZZ}f_t^n(U),$ 
 \item for any $t\in [0,1]$, ${f_t}{/U}$ is twist;
\item for any $t<1$ the set $\La_t$ is hyperbolic;
\item for $t=1$ all the periodic points in $\La_1$ are hyperbolic but $\La_1$ is not a hyperbolic set;
\item  $f_1$ has not conjugate points in $U.$
\end{enumerate}

\begin{remark}\label{nc open}
 Observe that the non-conjugate property is an open property.
\end{remark}


\begin{lemma}  Let $\{f_t\}_{t\in [0,1]}$ be a one parameter family of surface diffeomorphisms satisfying the  properties (1,2,3,5) listed above.  Then,   there exists $W^{cu(cs)}: \La_1\to Emb^{Lip}((-1,1), M)$ such for any $\e>0$ there exists $\de$ that if $x\in \La_1$ follows 
$$f^{-n}(W^{cu}_\e(x))\subset W^{cu}_\de(f^{-n}(x)).$$
 
\end{lemma}

\begin{proof} Let $x\in \La$ and let $x_t\in \La_t\to x$. For any $t$ close to $1$, there exist $W^u_t: \La_1\to Emb^{r}((-1,1), M)$ such that $W^u_{t,\e}(x_t)$ is the local unstable manifold of $x_t\in \La_t.$ Let us consider $W^u_{t, U}(x_t)$ as the connected component of $W^u_{t}(x_t)$ contained in $U$ that contains $x.$ From remark \ref{nc open}, it follows that $W^u_{t, U}(x_t)$ is a a $C^r$ curve such that $T_zW^u_{t, U}(x)\subset [\tilde G^-(z), \tilde G^+(z)].$ Moreover, since for any $t$ follows that  $x_t$ is far from the boundary of $U$, the curve has $W^u_{t, U}(x_t)$ has a length bounded by below. Taking a limit of that curves, we get $W^{cu}_U(x).$
 

\end{proof}

\begin{theorem}  Let $\{f_t\}_{t\in [0,1]}$ be a one parameter family of surface diffeomorphisms satisfying the  properties (1,2,3,4,5) and  also that $f_1$ is conjugated to $f_0$.  Then,   there exists a tangency between a stable and unstable local set.  
 
\end{theorem}




\begin{lemma} for any $t<1$ follows that $G^-_{x_t, t}< G^+_{x_t, t}$ and $G^-_{x_t, t}< G^+_{x_t, t}$ where $x_t=h_t(x).$
 
\end{lemma}

\begin{proof} The lemma holds for $t<1$ since ${f_t}{/\La_t}$ is twist and $\La_t$ is hyperbolic and therefore $G^-$ and $G^+$ coincide with the hyperbolic subbundles. For $t=1$ observe that fixed $n$ if $t\to 1$ folows that $x_t\to x_1$ and  $G^n_{f^{-n}_t(x_t)}(v^+)\to G^n_{f^{-n}_1(x_1)}(v^+)$ and so, if there exists $x_1$ such that $G^-_{x_1, 1}> G^+_{x_1, 1}$ then there exists $n$ such that $G^{-n}_{f^{n}_1(x_1)}(v^+)> G^n_{f^{-n}_1(x_1)}(v^+)$ and the same would hold for $t$ close to $1.$ A contradicion.
 
\end{proof}



\begin{lemma}
 For any $t$ follows that the sets  $D^+_{x,t}=D^{++}_{x,t}\cap D^{+-}_{x,t}$ and $D^-_{x,t}=D^{-+}_{x,t}\cap D^{--}_{x,t}$ are  not empty and contain interior.
\end{lemma}

\begin{proof} For $t<1$ it follows from the fact that $\La_t$ is hyperbolic. 
For $t=1$, it follows that $D^+_{x,1}=D^{++}_{x,1}\cap D^{+-}_{x,1}$ and $D^-_{x,1}=D^{-+}_{x,1}\cap D^{--}_{x,1}$ are the Hausdorff limit of the respetive sets for $t<1$.
Therefore, for $t=1$ the sets are non-empty. By... the local stable and unstable sets are contained in ... and is the boundary touches, the ocal stable and unstable would intersect in another point different than $x.$ A contradiction.

 
\end{proof}




\paragraph{Dissipative case.}







The following consequence allows to recover one of the main result of [PR].

Comparison with definition in \cite{PRH}.

\begin{corollary}\label{some growth}
If the set $\Lambda$ satisfies (H), then for any $x\in \Lambda$
and $v\in T^1M$ satisfying the definition of critical point as in Proposition 1,
there exists a sequence $n^+_i\to +\infty$ such that
$(f^{n^+_i}(x),g^{n_i}(v))$
converges  to a pair $(z^+,w^+)$ such that for each $n\geq 1$,
$$g^n_z(w^+)\geq b^{n},\quad g^{-n}_z(w^+)\geq b^{-n}.$$
There exists another sequence $n^-_i\to +\infty$ such that
$(f^{n^-_i}(x),g^{n_i}(v))$
converges  to a pair $(z^-,w^-)$ such that for each $n\geq 1$,
$$g^{-n}_z(w^-)\geq b^{n},\quad g^{n}_z(w^-)\geq b^{-n}.$$
Moreover, there exists a sequence $k_i$ such that $f^{k_i}(z^-)\to z_+$ and $G^{k_i}(w^-)\to w_+.$
\end{corollary}


From  previous corollary relates the present definition of critical points with the one provided in \cite{PRH}. On the other hand, the  forward orbit of $z^-$ accumulates on the critical set and the backward orbit  of $z^+$ accumulates on the critical set, where $z^-$ and $z^+$ are the points in previous corollary.


\begin{corollary}\label{crit pliss} Let $c$ be a critical point and let $n_i^-, n_i^+$ be the positive integers given by proposition \ref{crit exp}; then there exist sequences $m_i^-, m_i^+$ such that
$$|g^j(G^-{m_i^-}(v))|< \la^j\,\,\,\, 0<j < m_i^-,$$
$$|g^{-j}(G^{m_i^+}(v))|< \la^j\,\,\,\, 0<j < m_i^+.$$

 
\end{corollary}



\paragraph{l-- Perturbations of minimal dynamics.}


\begin{definition}Let  $f\in \operatorname{Diff}^{1}(\AA)$  where $\AA$ is an annulus, $f$ is twist, $f(\AA)\subset \AA$ and let assume that $\La$ is a non-hyperbolic minimal  aperiodic  attractor. It is said that $f$ satisfies an $r-$good distribution condition if there exists infinitely many $q_n$ such that 
\begin{enumerate}
 \item $dist(f^{q_n}(c), c) > \la^{\frac{q_n}{r+1}}$,

 \item there exists $x$ such that $dist(f^{q_n}(x), c) > \la^{\frac{q_n}{r+1}}$ and $dist(f^{q_n}(x), x)< \ga^n,$ for some $\ga<1.$

\end{enumerate}

 
\end{definition}



\begin{theorem}\label{perturbation minimal } Let  $f\in \operatorname{Diff}^{1}(\AA)$  where $\AA$ is an annulus, $f$ is twist, $f(\AA)\subset \AA$ and let assume that $\La$ is a non-hyperbolic minimal  aperiodic  attractor satisfying the $r-$good distribution condition. There exists $g$ $C^r-$arbitrary close to $f$ such that $\La_g$ is not a graph.
 
\end{theorem}





\begin{theorem}\label{perturbation minimal entropy } Let  $f\in \operatorname{Diff}^{1}(\AA)$  where $\AA$ is an annulus, $f$ is twist, $f(\AA)\subset \AA$ and let assume that $\La$ is a non-hyperbolic minimal  aperiodic  attractor satisfying the $r-$good distribution condition. There exists $g$ $C^r-$arbitrary close to $f$ such that $g/\La_g$ has positive topological entropy.
\end{theorem}






\noi{\em Proof of theorem \ref{perturbation minimal}:}

1- Take $m^-$ and $m^+$ large such that $f^{-m^-}(c)$ and $f^{-m^+}(c)$ satisfy lemma \ref{vert pliss}, and let $q_n$ such  $q_{n-1}< \max\{m^-, m^+\}.$
  
2- We consider a perturbation $\hat f$ of $f$ in the following way: i- take  ball $B_1$ of radius $\la^{\frac{q_n}{r+1}}$ around $c$ and a ball $B_2$ of radius $\la^{{q_n}}$ around $c$; in the complement of $B_1$ $g=f$ and in $B_2$ $\hat f=R_a\circ f$ where $R_a$ is a rotation of radius $\la^{q_n}$ centered at $c.$

3- Observe that from the first condition of the $r-$good distribution property $f^i(c)\notin B_1$ for $0<|i|\leq \max\{m^-, m^+\}.$ In particular, $G^i_{f}(v^+)=G^i_{\hat f}(v^+)$ and therefore it holds that  $slope(G^{m^-}(v^+), v)< \la^{m^-}$ and $slope(G^{-m^+}(v^+), v)< \la^{m^+}.$

4- At the point $f^{-m^-}(c)$ there is a stable arc $W^s_m$ such that points in that arc are contracted by rate $\la'<\la$ and such that it intersects the boundaries of the annulus. Moreover, if $x\in W^s_m$ it follows that $dist(f^{i}(x), c)< {\la'}^{i}$ and  $slope(G^{m^-}_x(v^+), v)<\la^{m^-}.$ In particular $f^i(x)\notin B_2$ for any $0\leq i< m^-.$

5- Therefore, for $G^{m^-}_{\hat f}(v^+)< G^{-m^+}_{\hat f}(v^+)$ so $\hat f^{m^-}(\gamma)$ cannot be a graph (where $\gamma $ is one boundary of the annulus) contradicting remark \ref{graph bound}.

\qed





% 
% \subsection{Conjugate points and torsion}
% 
% 
% We also assume that $f$ is twist in $U$ respect to the vertical and it has no torsion.
% Let $z\in U$ and $v\in T_zM\setminus\{0\}$. We denote with $\theta_0(v)$  the lift to $\RR$ of the angle contained in $(-1, 0]$ of the orientated angle between the vertical $v^+$ and the vector $v.$ We denote $\theta_z(v^+) = \theta_0(D_zf(v^+))$ and more generally $\theta_1(v)$ the lift to $\RR$  contained in  $(\theta_z(v^+)-1, \theta_z(v^+)]$ of the orientated angle between  $v^+$ and $D_zfv$. We define $$\theta(v) = \theta_1(v)-\theta_0 (v).$$
% The map $\theta: T_z M\setminus \{0\} \to \RR$is the unique continuous function that lift the angle between $v$ and $D_zfv$ such that  $\theta(v^+) \in (-1 , 0].$ Moreover, $\theta(v)$ depends continuosly on $v$ and $z$. It is also defined
% 
% $$\theta_n(v)=\left\{\begin{array}{clcr}\sum_{0\leq k\leq n-1} \theta(D_zf^kv)\,\,\,\mbox{if}\,\, n\geq 0\\ -\theta_{-n}(Df^nv)\,\,\,\mbox{if} \,\,n\leq 0\end{array}\right.$$
% 
% \begin{definition}{\bf Null torsion} It is said that the torsion is null if for any $z$ and $v$ follows that $$\frac{\theta_n(v)}{n}\to 0.$$
%  
% \end{definition}
% 
% 
% \begin{lemma} Let $f_{\La}$ be a twist diffeomorphsims such that the torsion is null for any point in a neighborhood $U$ of $\La.$ Then there is not conjugate points in $U.$
%  
% \end{lemma}
% 
% 
% \vskip 10pt
% \noindent
% \\
% Maybe work with a two dimensioal version of the  the family $2x+\omega+a\sin(2\pi x)$.
% \smallskip
% \noindent
% \\
% Or even the  henon map.


% Let us assume that it is false, i.e. for any $n$ there exists $x$ such that $|Df_{/E^s_x}|>\la> \la^n$. Then we can chose $x$ such that $|Df_{/E^s_x}|>\la^k$ for any $k=1\dots n$ which implies that $g^k(E^s_x)< [\frac{b}{\la}]^k$ for any $k=1\dots n.$  Observe that there exists another direction $\hat E_x$ such that for any vector $w$ in the positive quadrant follows that $sl(G^k(w), G^k(E^s_x))\to 0$ and for any vector $w$ in the negative quadrant follows that $sl(G^k(w), G^k(-E^s_x))\to 0.$ Moreover, if $sl(w,\hat E_x)>d_0$ for some positive $d_0$ it follows that $sl(G^k(w), G^k(E^s_x))< c_0[\frac{b}{\la}]^k sl(E^s_x, w)$ for some constant $c_0$ that depends on $d_0$.  Since $-E^s_x\in H^-_x$ and its iterate also belongs to the negative half, follows that that $G^k(\hat E_x) $ belongs to $H^-_{f^k(x)}$ and in particular, $v^+$ belongs to the positive quadrant bounded by $\hat E_x$ and $E^s_x$; moreover, observe that, $sl(G(v^+), G(\hat E_x))>\al_0$ 
%  so it follows that $sl(G^k(v^+), G^k(E^s_x))< c_0[\frac{b}{\la}]^k sl(v^+, E^s_x)< c_0[\frac{b}{\la}]^k\al_0$ but since $G^k(v^+)$ converge to the subbundle $F$ and the slope between $F$ and $E^s$ is bounded away by $\be_0$ therefore the slope between $G^k(v^+)$ and the stable subbundle is larger than $\be_0/2$ so it follows that $k<\frac{\log(\al_0/c_0)}{\log(b/\la)}.$
